% L'aide au développement, les échanges commerciaux et la bourse des valeurs — Allemand Tle A — Cours signé OnBuch+
% Programme officiel MINESEC. Compiler avec : tectonic AL09-l-aide-au-developpement-les-echanges-commerciaux-et-la-bours.tex
\newcommand{\DOCMATIERE}{Allemand}
\newcommand{\DOCNIVEAU}{Tle A}
\newcommand{\DOCMODULE}{Chapitre 2 : Vie économique}
\newcommand{\DOCLECON}{AL09}
\newcommand{\DOCTITRE}{L'aide au développement, les échanges commerciaux et la bourse des valeurs}
% OnBuch+ — préambule « cours signés » ENRICHI (compilé avec tectonic / XeLaTeX)
% Système visuel « Pop » d'OnBuch+ : fond crème (jamais blanc), bordures encre
% épaisses, ombres « dures » décalées, titres Archivo Black en capitales.
% Version MATHÉMATIQUES : graphes (pgfplots/tikz), boîtes pédagogiques
% (définition, propriété, méthode, attention, le savais-tu, exercice résolu,
% exercice, TP, bilan), page de garde et signature OnBuch+ ancrée en bas.
%
% Macros attendues dans main.tex AVANT \input{preamble} :
%   \DOCMATIERE  \DOCNIVEAU  \DOCMODULE  \DOCLECON (numéro)  \DOCTITRE
\documentclass[11pt]{article}

\usepackage{fontspec}
\usepackage[ngerman,french]{babel} % français = langue principale ; l'allemand via \all{..} / textebox
\usepackage[a4paper,margin=2cm,top=3.4cm,bottom=2.8cm,headheight=1.4cm,footskip=1.3cm]{geometry}
\usepackage{amsmath,amssymb,amsfonts}
\usepackage{mathtools} % \xmapsto, \coloneqq... souvent utilisés par les modèles
\usepackage{cancel} % simplifications barrées (bilans de forces)
% \abs{z} (module, valeur absolue) et \norm{u} (norme), utilisés par les modèles
\providecommand{\abs}{}\let\abs\relax\DeclarePairedDelimiter{\abs}{\lvert}{\rvert}
\providecommand{\norm}{}\let\norm\relax\DeclarePairedDelimiter{\norm}{\lVert}{\rVert}
\usepackage[table]{xcolor} % [table] : fournit aussi \rowcolors (alternance de lignes)
\usepackage{pagecolor}
\usepackage{tikz}
\usetikzlibrary{arrows.meta,positioning,calc,shapes.geometric,shapes.misc,shapes.arrows,decorations.pathmorphing,decorations.pathreplacing,decorations.markings,patterns,shadows,mindmap,trees,fit,backgrounds,matrix,intersections,angles,quotes}
\usepackage{pgfplots}
\usepackage[version=4]{mhchem} % équations nucléaires \ce{^{235}_{92}U}
\usepackage[european,siunitx]{circuitikz} % schémas électriques
\pgfplotsset{compat=1.18}
\usepgfplotslibrary{fillbetween,groupplots} % aires entre courbes (calcul intégral)
\usepackage{enumitem}
\usepackage{fancyhdr}
% [most] charge toutes les bibliothèques tcolorbox (listings, minted, raster,
% documentation...) et gonfle inutilement la mémoire TeX sur un document long
% (~300 boîtes) : on ne charge que ce qui est réellement utilisé.
\usepackage[skins,breakable]{tcolorbox}
\usepackage{graphicx}
\usepackage{colortbl}
\usepackage{booktabs}
\usepackage{tabularx}
\usepackage{diagbox} % case d'en-tête diagonale des tableaux à double entrée
% Colonnes extensibles centrée / gauche / droite (C, L, R), souvent utilisées par les modèles
\newcolumntype{C}{>{\centering\arraybackslash}X}
\newcolumntype{L}{>{\raggedright\arraybackslash}X}
\newcolumntype{R}{>{\raggedleft\arraybackslash}X}
\usepackage{array}
\usepackage{multicol}
\usepackage{multirow}
\usepackage{siunitx}
% parse-numbers=false : accepte aussi \SI{2,5 \times 10^{-3}}{...}
\sisetup{locale=FR,output-decimal-marker={,},group-minimum-digits=4,parse-numbers=false}
\DeclareSIUnit\ohm{\ensuremath{\symup{\Omega}}} % U+2126 absent de la police
\DeclareSIUnit\division{div} % réglages d'oscilloscope (ms/div, V/div)
\DeclareSIUnit\tr{tr} % tours (vitesse de rotation, ex. \SI{1500}{\tr\per\minute})
\DeclareSIUnit\molar{M} % concentration molaire (ex. \SI{3}{\molar}, \SI{1}{\micro\molar})
\DeclareSIUnit\an{an} % année (le modèle confond parfois avec \year, un compteur TeX) — ex. \SI{5}{\kilo\meter\per\an}
\DeclareSIUnit\FCFA{FCFA} % franc CFA (ex. \SI{500}{\FCFA\per\kilo\gram})
\DeclareSIUnit\degreeBrix{\textdegree Brix} % teneur en sucre d'un sirop/jus (réfractométrie)
\DeclareSIUnit\month{mois} % le modèle utilise parfois \month, un compteur TeX, comme unité de durée
\usepackage{qrcode}
% \degree : commande fréquemment utilisée par les modèles pour un angle ou une
% température en dehors de \SI{}{} (non fournie par les paquets chargés).
\providecommand{\degree}{^{\circ}}
% \qed (fin de démonstration), utilisé par les modèles sans amsthm
\providecommand{\qed}{\hfill$\square$}
% \barre : double barre des valeurs interdites dans un tableau de variations (tkz-tab)
\providecommand{\barre}{\ensuremath{\|}}
% environnement proof (amsthm non chargé) utilisé par les modèles
\ifdefined\proof\else\newenvironment{proof}[1][Démonstration]{\par\noindent\textit{#1.}\ }{\qed\par}\fi
\usepackage{lastpage}
\usepackage{hyperref}
\hypersetup{hidelinks}

% ============================================================
% Palette « Pop » OnBuch+ — mêmes hex que la classe P (plus_ui.dart)
% ============================================================
\definecolor{poporange}{HTML}{F59321}
\definecolor{popdark}{HTML}{E07A0C}
\definecolor{popcream}{HTML}{FFF6E8}
\definecolor{popink}{HTML}{1C1714}
\definecolor{poppurple}{HTML}{7A5AE0}
\definecolor{popgreen}{HTML}{2E9E6A}
\definecolor{poppink}{HTML}{E0457B}
\definecolor{popblue}{HTML}{2F7DE1}
\definecolor{popgold}{HTML}{C98A12}
\definecolor{popmuted}{HTML}{8A7F76}
\definecolor{popline}{HTML}{EADFD2}
\definecolor{popgray}{HTML}{8A7F76}
\colorlet{popgrey}{popgray}
% Alias tolérants (le modèle nomme parfois les couleurs différemment)
\colorlet{popwhite}{white}
\colorlet{popblack}{popink}
\colorlet{popred}{poppink}
\colorlet{popyellow}{popgold}
% Teintes claires (fonds de tableaux/encadrés) — le modèle nomme parfois
% les couleurs avec un suffixe L (ex. popblueL) sans les avoir définies.
\colorlet{poporangeL}{poporange!20}
\colorlet{popdarkL}{popdark!20}
\colorlet{poppurpleL}{poppurple!20}
\colorlet{popgreenL}{popgreen!20}
\colorlet{poppinkL}{poppink!20}
\colorlet{popblueL}{popblue!20}
\colorlet{popgoldL}{popgold!20}
\colorlet{popredL}{popred!20}
\colorlet{popgrayL}{popgray!20}
% Teintes claires pour fonds de tableaux / zones
\colorlet{popblueL}{popblue!10!white}
\colorlet{poporangeL}{poporange!14!white}
\colorlet{popgreenL}{popgreen!12!white}
\colorlet{poppurpleL}{poppurple!12!white}
\colorlet{poppinkL}{poppink!12!white}
\colorlet{popgoldL}{popgold!14!white}

\pagecolor{popcream}

% ============================================================
% Typographie
% ============================================================
\defaultfontfeatures{Ligatures=TeX}
\setmainfont{PlusJakartaSans-Medium.ttf}[Path=../../../fonts/,BoldFont=PlusJakartaSans-Bold.ttf]
% Maths en sans-serif (Fira Math) avec chiffres et lettres droites de Plus
% Jakarta, pour que les formules se fondent dans le texte.
\usepackage[math-style=ISO,bold-style=ISO]{unicode-math}
\setmathfont{FiraMath-Regular.otf}[Path=../../../fonts/]
\setmathfont{PlusJakartaSans-Medium.ttf}[Path=../../../fonts/,range={up/num,up/{Latin,latin},\mathup}]
\usepackage{icomma} % virgule décimale sans espace en mode math
% Caractères mathématiques Unicode absents des polices de texte (Plus Jakarta,
% Archivo Black) : rendus par leur équivalent mathématique, en texte comme en
% formule — sinon ils s'affichent en carré vide (ex. « Arithmétique dans ℤ »).
\usepackage{newunicodechar}
\newunicodechar{ℝ}{\ensuremath{\mathbb{R}}}
\newunicodechar{ℤ}{\ensuremath{\mathbb{Z}}}
\newunicodechar{ℂ}{\ensuremath{\mathbb{C}}}
\newunicodechar{ℕ}{\ensuremath{\mathbb{N}}}
\newunicodechar{ℚ}{\ensuremath{\mathbb{Q}}}
\newunicodechar{𝔻}{\ensuremath{\mathbb{D}}}
\newunicodechar{α}{\ensuremath{\alpha}}
\newunicodechar{β}{\ensuremath{\beta}}
\newunicodechar{γ}{\ensuremath{\gamma}}
\newunicodechar{δ}{\ensuremath{\delta}}
\newunicodechar{ε}{\ensuremath{\varepsilon}}
\newunicodechar{θ}{\ensuremath{\theta}}
\newunicodechar{λ}{\ensuremath{\lambda}}
\newunicodechar{μ}{\ensuremath{\mu}}
\newunicodechar{σ}{\ensuremath{\sigma}}
\newunicodechar{τ}{\ensuremath{\tau}}
\newunicodechar{φ}{\ensuremath{\varphi}}
\newunicodechar{ψ}{\ensuremath{\psi}}
\newunicodechar{ω}{\ensuremath{\omega}}
\newunicodechar{Σ}{\ensuremath{\Sigma}}
% majuscules produites par \MakeUppercase dans les titres (α-aminés -> Α)
\newunicodechar{Α}{\ensuremath{\alpha}}
\newunicodechar{Β}{\ensuremath{\beta}}
\newunicodechar{Γ}{\ensuremath{\Gamma}}
\newunicodechar{Δ}{\ensuremath{\Delta}}
\newunicodechar{Ε}{\ensuremath{\varepsilon}}
\newunicodechar{Θ}{\ensuremath{\Theta}}
\newunicodechar{Λ}{\ensuremath{\Lambda}}
\newunicodechar{Μ}{\ensuremath{\mu}}
\newunicodechar{Τ}{\ensuremath{\tau}}
\newunicodechar{Φ}{\ensuremath{\Phi}}
\newunicodechar{Ψ}{\ensuremath{\Psi}}
\newunicodechar{Ω}{\ensuremath{\Omega}}
\newunicodechar{↦}{\ensuremath{\mapsto}}
\newunicodechar{∈}{\ensuremath{\in}}
\newunicodechar{∉}{\ensuremath{\notin}}
\newunicodechar{∧}{\ensuremath{\wedge}}
\newunicodechar{⇔}{\ensuremath{\Leftrightarrow}}
\newunicodechar{⟺}{\ensuremath{\Longleftrightarrow}}
\newunicodechar{⇒}{\ensuremath{\Rightarrow}}
\newunicodechar{⟹}{\ensuremath{\Longrightarrow}}
\newunicodechar{∘}{\ensuremath{\circ}}
\newunicodechar{∪}{\ensuremath{\cup}}
\newunicodechar{∩}{\ensuremath{\cap}}
\newunicodechar{≡}{\ensuremath{\equiv}}
\newunicodechar{⊂}{\ensuremath{\subset}}
\newunicodechar{⊥}{\ensuremath{\perp}}
\newunicodechar{∃}{\ensuremath{\exists}}
\newunicodechar{∀}{\ensuremath{\forall}}
\newunicodechar{∛}{\ensuremath{\sqrt[3]{\,}}}
\newunicodechar{∜}{\ensuremath{\sqrt[4]{\,}}}
\newunicodechar{⃗}{\ensuremath{{}^{\to}}}
\newunicodechar{ⁿ}{\ifmmode^{n}\else\textsuperscript{\ensuremath{n}}\fi}
\newunicodechar{ˣ}{\ifmmode^{x}\else\textsuperscript{\ensuremath{x}}\fi}
\newunicodechar{ᵇ}{\ifmmode^{b}\else\textsuperscript{\ensuremath{b}}\fi}
\newunicodechar{ʳ}{\ifmmode^{r}\else\textsuperscript{\ensuremath{r}}\fi}
\newunicodechar{ᵘ}{\ifmmode^{u}\else\textsuperscript{\ensuremath{u}}\fi}
\newunicodechar{ʸ}{\ifmmode^{y}\else\textsuperscript{\ensuremath{y}}\fi}
\newunicodechar{ᵃ}{\ifmmode^{a}\else\textsuperscript{\ensuremath{a}}\fi}
\newunicodechar{ᵉ}{\ifmmode^{e}\else\textsuperscript{\ensuremath{e}}\fi}
\newunicodechar{ᵏ}{\ifmmode^{k}\else\textsuperscript{\ensuremath{k}}\fi}
\newunicodechar{ᵖ}{\ifmmode^{p}\else\textsuperscript{\ensuremath{p}}\fi}
\newunicodechar{ᴺ}{\ifmmode^{N}\else\textsuperscript{\ensuremath{N}}\fi}
\newunicodechar{ᶜ}{\ifmmode^{c}\else\textsuperscript{\ensuremath{c}}\fi}
\newunicodechar{ʰ}{\ifmmode^{h}\else\textsuperscript{\ensuremath{h}}\fi}
\newunicodechar{ᵠ}{\ifmmode^{q}\else\textsuperscript{\ensuremath{q}}\fi}
\newunicodechar{ₙ}{\ifmmode_{n}\else\textsubscript{\ensuremath{n}}\fi}
\newunicodechar{ₐ}{\ifmmode_{a}\else\textsubscript{\ensuremath{a}}\fi}
\newunicodechar{ᵢ}{\ifmmode_{i}\else\textsubscript{\ensuremath{i}}\fi}
\newunicodechar{ᵣ}{\ifmmode_{r}\else\textsubscript{\ensuremath{r}}\fi}
\newunicodechar{ₖ}{\ifmmode_{k}\else\textsubscript{\ensuremath{k}}\fi}
\newunicodechar{ᵦ}{\ifmmode_{\beta}\else\textsubscript{\ensuremath{\beta}}\fi}
\newunicodechar{ₓ}{\ifmmode_{x}\else\textsubscript{\ensuremath{x}}\fi}
\newfontfamily\popdisplay{ArchivoBlack-Regular.ttf}[Path=../../../fonts/]
\newfontfamily\poplogo{SpaceGrotesk-Bold.ttf}[Path=../../../fonts/]
\newfontfamily\popsemi{PlusJakartaSans-SemiBold.ttf}[Path=../../../fonts/]
\newfontfamily\popxbold{PlusJakartaSans-ExtraBold.ttf}[Path=../../../fonts/]
\newfontfamily\popxboldls{PlusJakartaSans-ExtraBold.ttf}[Path=../../../fonts/,LetterSpace=3.0]

\setlist[enumerate]{leftmargin=1.7em,itemsep=5pt,topsep=4pt}
\setlist[itemize]{leftmargin=1.7em,itemsep=4pt,topsep=4pt,label={\color{poporange}\textbullet}}
\setlength{\parindent}{0pt}
\setlength{\parskip}{5pt}
\renewcommand{\arraystretch}{1.25}

% pgfplots : style Pop par défaut
\pgfplotsset{
  every axis/.append style={axis line style={popink,line width=1pt},tick style={popink},
    grid style={popline,line width=0.5pt},label style={font=\small},tick label style={font=\footnotesize}},
  popcurve/.style={poporange,line width=1.8pt,no markers,smooth},
}
% Même style disponible dans un \draw TikZ simple (hors pgfplots)
\tikzset{popcurve/.style={poporange,line width=1.8pt}}
\tikzset{popgrid/.style={popline,very thin}}
\colorlet{popcurve}{poporange} % le nom du style est parfois utilisé comme couleur

% ============================================================
% Pill / squiggle
% ============================================================
\newcommand{\poppill}[3]{%
  \tikz[baseline=-0.55ex]\node[draw=popink,line width=1.2pt,rounded corners=2.6pt,%
    fill=#2,inner xsep=6.5pt,inner ysep=3pt]{%
    \popxboldls\fontsize{7}{7}\selectfont\color{#3}\MakeUppercase{#1}};%
}
\newcommand{\popsquiggle}[1]{%
  \tikz[baseline=-0.4ex]{
    \draw[#1,line width=2.6pt,line cap=round]
      plot [smooth] coordinates {(0,0.09) (0.34,0.01) (0.68,0.21) (1.02,0.01) (1.36,0.10)};
  }%
}
% Mot-clé surligné (marqueur orange sous le texte)
\newcommand{\cle}[1]{{\popxbold\color{popdark}#1}}

% ============================================================
% En-tête / pied
% ============================================================
\pagestyle{fancy}
\fancyhf{}
\renewcommand{\headrulewidth}{0pt}
\renewcommand{\footrulewidth}{0pt}
\lhead{\raisebox{-0.15ex}{\poplogo\fontsize{13}{13}\selectfont\color{popink}OnBuch\color{poporange}+}}
\rhead{\poppill{\DOCMATIERE\ \textbullet\ \DOCNIVEAU\ \textbullet\ Leçon \DOCLECON}{white}{popink}}
\lfoot{\popxbold\fontsize{7.6}{7.6}\selectfont\color{popmuted}\MakeUppercase{Cours signé OnBuch+}\ \textbullet\ {\popsemi\DOCTITRE}}
\rfoot{\tikz[baseline=-0.6ex]\node[rounded corners=8.5pt,fill=popink,minimum height=17pt,minimum width=17pt,inner xsep=5pt,inner sep=0]{\popxbold\fontsize{8}{8}\selectfont\color{popcream}\thepage\,/\,\pageref*{LastPage}};}

% ============================================================
% Titres
% ============================================================
\newcommand{\chaptitre}[2]{%
  \begin{tcolorbox}[enhanced,colback=poporange,colframe=popink,boxrule=2.6pt,arc=11pt,
      drop shadow=popink,left=15pt,right=15pt,top=12pt,bottom=11pt,before skip=3pt,after skip=18pt]
    \poppill{#2}{popink}{popcream}\par\vspace{9pt}
    {\raggedright\hyphenpenalty=10000\exhyphenpenalty=10000\popdisplay\fontsize{22}{24}\selectfont\color{popcream}\MakeUppercase{#1}\par}
  \end{tcolorbox}
}
% Section numérotée : pastille orange avec le numéro + titre Archivo
\newcounter{popsec}
\newcommand{\coursec}[1]{%
  \stepcounter{popsec}\setcounter{popsub}{0}%
  \par\vspace{16pt}\noindent
  \begin{minipage}{\linewidth}
  \tikz[baseline=-0.7ex]\node[draw=popink,line width=1.6pt,fill=poporange,rounded corners=4pt,minimum size=22pt,inner sep=0]{\popdisplay\fontsize{12}{12}\selectfont\color{popcream}\thepopsec};%
  \hspace{8pt}{\popdisplay\fontsize{15}{17}\selectfont\color{popink}\MakeUppercase{#1}}\par
  \vspace{4pt}\noindent{\color{poporange}\rule{36pt}{3.6pt}}
  \end{minipage}\par\nopagebreak\vspace{8pt}
}
\newcounter{popsub}
\newcommand{\courssub}[1]{%
  \stepcounter{popsub}%
  \par\vspace{9pt}\noindent{\popxbold\fontsize{11.5}{13}\selectfont\color{popink}{\color{poporange}\thepopsec.\thepopsub}\hspace{6pt}#1}\par\nopagebreak\vspace{4pt}
}

% ============================================================
% Boîtes pédagogiques — même recette : fond blanc, bordure épaisse colorée,
% ombre dure encre, badge plein. Titre optionnel [#1].
% ============================================================
\tcbset{popbase/.style={breakable,enhanced,colback=white,boxrule=2.3pt,arc=9pt,
  shadow={3pt}{-3pt}{0pt}{popink},left=14pt,right=14pt,top=11pt,bottom=11pt,before skip=11pt,after skip=15pt,
  fontupper=\popsemi\fontsize{10.6}{14.5}\selectfont\color{popink},
  fontlower=\popsemi\fontsize{10.6}{14.5}\selectfont\color{popink}}}
% Coche dessinée (le glyphe ✓ n'existe pas dans les polices du document)
\newcommand{\popcheck}[1]{\tikz[baseline=-0.5ex]\draw[#1,line width=1.6pt,line cap=round,line join=round](-0.13,0.0)--(-0.04,-0.09)--(0.13,0.1);}
\providecommand{\coche}{\popcheck{popgreen}}
\providecommand{\resultat}[1]{\par\smallskip\noindent\fcolorbox{popgreen}{popgreenL}{\parbox{\dimexpr\linewidth-2\fboxsep-2\fboxrule\relax}{\textbf{Résultat.} #1}}\par\smallskip}
\newcommand{\popboxhead}[3]{\poppill{#1}{#2}{white}\ifx\relax#3\relax\else\hspace{6pt}{\popxbold\fontsize{10.6}{12}\selectfont\color{popink}#3}\fi\par\vspace{7pt}}

\newtcolorbox{aretenir}[1][]{popbase,colframe=popgreen,before upper={\popboxhead{À retenir}{popgreen}{#1}}}
% Texte en allemand (document de lecture, dialogue, extrait) : cadre
% d'étude. Un titre optionnel (source, auteur) s'affiche dans l'en-tête.
\newtcolorbox{textebox}[1][]{popbase,colframe=popblue,before upper={\popboxhead{Text}{popblue}{#1}\begin{otherlanguage}{ngerman}},after upper={\end{otherlanguage}}}
% Marques ✓ / ✗ (forme correcte / fautive) souvent utilisées par les modèles
\providecommand{\cross}{\textcolor{poppink}{\ensuremath{\times}}}
\providecommand{\xmark}{\cross}\providecommand{\crossmark}{\cross}\providecommand{\cmark}{\ensuremath{\checkmark}}
% Ligne de réponse / justification à compléter (inventée par les modèles)
\providecommand{\justification}[1]{\par\noindent\textit{Justification~:} \dotfill\par}
% \textipa (package tipa non chargé) : la police ne couvre pas l'API -> simple passage
\providecommand{\textipa}[1]{#1}
% Soulignement (ulem non chargé) : \uline, \ul
\providecommand{\uline}[1]{\underline{#1}}\providecommand{\ul}[1]{\underline{#1}}
% \alert (beamer) inventée par les modèles : texte mis en évidence
\providecommand{\alert}[1]{\textbf{\textcolor{poppink}{#1}}}
% variantes ASCII d'ordinaux écrites en macro
\providecommand{\iere}{\textsuperscript{re}}\providecommand{\ieme}{\textsuperscript{e}}\providecommand{\ere}{\textsuperscript{re}}\providecommand{\eme}{\textsuperscript{e}}\providecommand{\er}{\textsuperscript{er}}\providecommand{\ie}{\textsuperscript{e}}
% Ordinaux français écrits en macro par les modèles : 1\ière, 3\ième
\newcommand{\ière}{\textsuperscript{re}}\newcommand{\ième}{\textsuperscript{e}}
% \exemple (inventée par les modèles) : étiquette « Exemple : »
\providecommand{\exemple}{\textbf{Exemple~:}\ }
% \square utilisable aussi hors mode math (cases à cocher)
\AtBeginDocument{\let\squareMath\square\renewcommand{\square}{\ensuremath{\squareMath}}}
% Mot ou phrase en allemand à l'intérieur d'un paragraphe français
\newcommand{\all}[1]{\ifmmode\text{\textit{#1}}\else\begingroup\selectlanguage{ngerman}\itshape #1\endgroup\fi}
\let\esp\all % alias toléré
\newtcolorbox{exemplebox}[1][]{popbase,colframe=poporange,before upper={\popboxhead{Exemple}{poporange}{#1}}}
\newtcolorbox{definition}[1][]{popbase,colframe=popblue,before upper={\popboxhead{Définition}{popblue}{#1}}}
\newtcolorbox{propriete}[1][]{popbase,colframe=poppurple,before upper={\popboxhead{Propriété}{poppurple}{#1}}}
\newtcolorbox{methode}[1][]{popbase,colframe=popgold,before upper={\popboxhead{Méthode}{popgold}{#1}}}
\newtcolorbox{attention}[1][]{popbase,colframe=poppink,before upper={\popboxhead{Attention}{poppink}{#1}}}
\newtcolorbox{experience}[1][]{popbase,colframe=popblue,colback=popblueL,before upper={\popboxhead{Expérience}{popblue}{#1}}}
% « Le savais-tu ? » : carte violette pleine (contexte, culture, Cameroun)
\newtcolorbox{savaistu}[1][]{popbase,colback=poppurple,colframe=popink,
  fontupper=\popsemi\fontsize{10.4}{14.2}\selectfont\color{white},
  before upper={\poppill{Le savais-tu\,?}{white}{poppurple}\ifx\relax#1\relax\else\hspace{6pt}{\popxbold\color{white}#1}\fi\par\vspace{7pt}}}
% Exercice résolu : énoncé en haut, \tcblower puis la solution sur fond crème
\newcounter{popexo}\newcounter{popexr}
\newtcolorbox{exoresolu}[1][]{popbase,colframe=poporange,colbacklower=popcream,
  segmentation style={popink,line width=1.2pt,dashed},
  before upper={\stepcounter{popexr}\popboxhead{Exercice résolu \thepopexr}{poporange}{#1}},
  before lower={\poppill{Solution}{popink}{popcream}\par\vspace{6pt}}}
% Exercice (non résolu en place) — la correction est dans la partie Corrigés
\newtcolorbox{exercice}[1][]{popbase,colframe=popink,
  before upper={\stepcounter{popexo}\popboxhead{Exercice \thepopexo}{popink}{#1}}}
% Corrigé
\newtcolorbox{corrige}[1][]{popbase,colframe=popgreen,colback=popgreenL,
  before upper={\popboxhead{Corrigé}{popgreen}{#1}}}
% Objectifs / prérequis / bilan
\newtcolorbox{objectifs}{popbase,colframe=popink,colback=poporangeL,
  before upper={\popboxhead{Objectifs de la leçon}{popink}{}}}
\newtcolorbox{prerequis}{popbase,colframe=popink,colback=popblueL,
  before upper={\popboxhead{Prérequis}{popblue}{}}}
\newtcolorbox{bilan}{popbase,colframe=popink,colback=white,boxrule=2.8pt,
  before upper={\popboxhead{Fiche bilan}{poporange}{L'essentiel à connaître pour l'examen}}}
% Figure encadrée avec légende
\newtcolorbox{popfigure}[1][]{enhanced,breakable=false,colback=white,colframe=popline,boxrule=1.4pt,arc=8pt,
  left=10pt,right=10pt,top=10pt,bottom=8pt,before skip=10pt,after skip=12pt,halign=center,
  lower separated=false,halign lower=center,
  fontlower=\popsemi\fontsize{9}{11.5}\selectfont\color{popmuted},
  #1}
\newcommand{\legende}[1]{\par\vspace{6pt}{\popsemi\fontsize{9}{11.5}\selectfont\color{popmuted}#1}}

% ============================================================
% Signature OnBuch+
% ============================================================
\newcommand{\poplogoinline}[1]{{\poplogo\fontsize{#1}{#1}\selectfont\color{popink}OnBuch\color{poporange}+}}
\newcommand{\signatureonbuch}{%
  \begin{tcolorbox}[enhanced,colback=popink,colframe=popink,boxrule=2.2pt,arc=10pt,
      drop shadow=poporange,left=14pt,right=14pt,top=10pt,bottom=10pt,before skip=0pt,after skip=0pt]
    \tikz[baseline=-0.7ex]\node[circle,fill=popgreen,draw=popcream,line width=1.3pt,minimum size=22pt,inner sep=0]{\popcheck{white}};%
    \hspace{9pt}\begin{minipage}[c]{0.62\linewidth}
      {\popxbold\fontsize{12}{13}\selectfont\color{popcream}Signé\ \textbullet\ {\poplogo OnBuch\color{poporange}+}}\par\vspace{2pt}
      {\raggedright\popsemi\fontsize{8}{10.5}\selectfont\color{popline}\MakeUppercase{Équipe pédagogique OnBuch+}\\\MakeUppercase{Conforme au programme officiel MINESEC}\par}
    \end{minipage}\hfill
    {\poplogo\fontsize{22}{22}\selectfont\color{popcream}OnBuch\color{poporange}+}
  \end{tcolorbox}%
}

% ============================================================
% Page de garde
% #1 titre · #2 matière · #3 niveau · #4 module · #5 n° leçon · #6 durée ·
% #7 description / objectifs (texte ou itemize)
% ============================================================
\newcommand{\pagegarde}[7]{%
  \thispagestyle{empty}
  \newgeometry{a4paper,margin=1.7cm,top=1.6cm,bottom=1.4cm}%
  \begin{tikzpicture}[remember picture,overlay]
    \fill[poporange!18] ([xshift=-3.2cm,yshift=-3.6cm]current page.north east) circle (4.6cm);
    \fill[poppurple!14] ([xshift=2.4cm,yshift=7.6cm]current page.south west) circle (3.8cm);
  \end{tikzpicture}%
  {\poplogo\fontsize{30}{30}\selectfont\color{popink}OnBuch\color{poporange}+}\hfill
  \raisebox{0.6ex}{\poppill{Cours signé \textbullet\ Premium}{popink}{popcream}}\par
  \vspace{1pt}\popsquiggle{poppurple}\par
  \vspace{8pt}
  \begin{tcolorbox}[enhanced,colback=poporange,colframe=popink,boxrule=2.8pt,arc=13pt,
      drop shadow=popink,left=18pt,right=18pt,top=12pt,bottom=13pt,after skip=10pt]
    \poppill{#2 \textbullet\ #3}{popink}{popcream}\hspace{5pt}\poppill{#4}{white}{popink}\par\vspace{8pt}
    {\popdisplay\fontsize{14}{14}\selectfont\color{popink}LEÇON #5}\par\vspace{4pt}
    {\raggedright\hyphenpenalty=10000\exhyphenpenalty=10000\popdisplay\fontsize{25}{27}\selectfont\color{popcream}\MakeUppercase{#1}\par}
    \vspace{8pt}
    {\popxbold\fontsize{9.5}{9.5}\selectfont\color{popink}\MakeUppercase{Durée conseillée : #6}}
  \end{tcolorbox}
  \begin{tcolorbox}[enhanced,colback=white,colframe=popink,boxrule=2.2pt,arc=10pt,
      drop shadow=popink,left=16pt,right=16pt,top=10pt,bottom=10pt,after skip=8pt]
    \poppill{Dans cette leçon}{poppurple}{white}\par\vspace{5pt}
    \popsemi\fontsize{9.3}{12.6}\selectfont\color{popink}#7
  \end{tcolorbox}
  \noindent\begin{minipage}[t]{0.487\linewidth}
    \begin{tcolorbox}[enhanced,colback=popgreen,colframe=popink,boxrule=2.2pt,arc=10pt,
        drop shadow=popink,left=11pt,right=11pt,top=10pt,bottom=10pt,height=3.1cm,valign=center]
      \tcbox[colback=white,colframe=popink,boxrule=1.3pt,arc=5pt,boxsep=0pt,nobeforeafter,
        left=3pt,right=3pt,top=3pt,bottom=3pt,valign=center]{\qrcode[height=1.9cm]{https://play.google.com/store/apps/details?id=com.luvvix.onbuch&hl=fr}}%
      \hspace{8pt}\begin{minipage}[c]{0.5\linewidth}
        {\popdisplay\fontsize{11}{12.5}\selectfont\color{white}\MakeUppercase{Télécharge OnBuch}}\par\vspace{4pt}
        {\raggedright\popsemi\fontsize{8.4}{11}\selectfont\color{white}Gratuit \textbullet\ Google Play\\Tuteur IA, annales, concours, résultats\par}
      \end{minipage}
    \end{tcolorbox}
  \end{minipage}\hfill
  \begin{minipage}[t]{0.487\linewidth}
    \begin{tcolorbox}[enhanced,colback=poppurple,colframe=popink,boxrule=2.2pt,arc=10pt,
        drop shadow=popink,left=11pt,right=11pt,top=10pt,bottom=10pt,height=3.1cm,valign=center]
      \tikz[baseline=-0.7ex]\node[circle,fill=white,draw=popink,line width=1.3pt,minimum size=1.6cm,inner sep=0]{\popdisplay\fontsize{24}{24}\selectfont\color{poppurple}+};%
      \hspace{8pt}\begin{minipage}[c]{0.6\linewidth}
        {\popdisplay\fontsize{11}{12.5}\selectfont\color{white}\MakeUppercase{Passe à OnBuch+}}\par\vspace{4pt}
        {\raggedright\popsemi\fontsize{8.4}{11}\selectfont\color{white}Tous les cours signés, Tuteur IA illimité, fiches PDF — onglet OnBuch+ de l'app\par}
      \end{minipage}
    \end{tcolorbox}
  \end{minipage}
  \vspace{8pt}
  \popcontenu
  \vfill
  \signatureonbuch
  \restoregeometry
  \newpage
}

% Rangée « Ce cours contient » (page de garde)
\newcommand{\poptile}[3]{% #1 couleur #2 gros texte #3 légende
  \begin{tcolorbox}[enhanced,colback=#1,colframe=popink,boxrule=2pt,arc=9pt,shadow={2.5pt}{-2.5pt}{0pt}{popink},
      left=6pt,right=6pt,top=9pt,bottom=9pt,halign=center,valign=center,height=1.75cm,width=\linewidth,nobeforeafter]
    {\popdisplay\fontsize{12.5}{13}\selectfont\color{white}\MakeUppercase{#2}}\par\vspace{3pt}
    {\centering\popsemi\fontsize{7.8}{9.5}\selectfont\color{white}#3\par}
  \end{tcolorbox}}
\newcommand{\popcontenu}{%
  \par\noindent{\popxboldls\fontsize{7.5}{7.5}\selectfont\color{popmuted}CE COURS SIGNÉ CONTIENT}\par\vspace{7pt}
  \noindent\begin{minipage}[t]{0.235\linewidth}\poptile{popblue}{Cours}{complet, illustré, pas à pas}\end{minipage}\hfill
  \begin{minipage}[t]{0.235\linewidth}\poptile{poporange}{Méthodes}{et exercices résolus}\end{minipage}\hfill
  \begin{minipage}[t]{0.235\linewidth}\poptile{poppink}{Exercices}{type examen + corrigés}\end{minipage}\hfill
  \begin{minipage}[t]{0.235\linewidth}\poptile{popgreen}{Fiche bilan}{l'essentiel à retenir}\end{minipage}\par}

% Sommaire
\newtcolorbox{sommaire}{enhanced,colback=white,colframe=popink,boxrule=2.2pt,arc=10pt,
  drop shadow=popink,left=18pt,right=18pt,top=13pt,bottom=13pt,before skip=6pt,after skip=20pt,
  before upper={\poppill{Sommaire}{poppurple}{white}\par\vspace{9pt}\popsemi\fontsize{10.6}{14.5}\selectfont\color{popink}}}

% Signature de fin de cours — poussée tout en bas de la dernière page
\newcommand{\findecours}{%
  \par\vfill\nopagebreak
  \begin{center}\popsquiggle{poporange}\end{center}\vspace{4pt}
  \signatureonbuch
}
\AtBeginDocument{\def\llbracket{\mathopen{[\mkern-3.5mu[}}\def\rrbracket{\mathclose{]\mkern-3.5mu]}}} % intervalles d'entiers (glyphe ⟦ absent de la police)
% Glyphes absents de Fira Math : redéfinis à partir de glyphes disponibles
\AtBeginDocument{%
  \def\setminus{\mathbin{\backslash}}\let\smallsetminus\setminus
  \def\vdots{\mathord{\vbox{\baselineskip3.2pt\lineskiplimit0pt\kern4pt\hbox{.}\hbox{.}\hbox{.}}}}%
  \def\ddots{\mathinner{\mkern1mu\raise6.5pt\hbox{.}\mkern2mu\raise3.6pt\hbox{.}\mkern2mu\raise0.7pt\hbox{.}\mkern1mu}}%
  \def\triangle{\mathord{\scalebox{0.9}{$\symup{\Delta}$}}}%
  \def\nabla{\mathord{\raisebox{\depth}{\scalebox{1}[-1]{$\symup{\Delta}$}}}}%
  \def\checkmark{\popcheck{popdark}}%
  \def\star{\mathbin{\ast}}\def\bigstar{\mathord{\ast}}%
  \def\diamond{\mathbin{\scalebox{0.6}{$\Diamond$}}}%
  \def\bigcup{\mathop{\mathchoice{\raisebox{-0.3em}{\scalebox{1.7}{$\cup$}}}{\scalebox{1.25}{$\cup$}}{\cup}{\cup}}\displaylimits}%
  \def\bigcap{\mathop{\mathchoice{\raisebox{-0.3em}{\scalebox{1.7}{$\cap$}}}{\scalebox{1.25}{$\cap$}}{\cap}{\cap}}\displaylimits}%
  \def\lesseqqgtr{\mathrel{\lessgtr}}\def\bigoplus{\mathop{\oplus}\displaylimits}%
}
\providecommand{\ssi}{si et seulement si} % abréviation fréquente
\AtBeginDocument{\providecommand{\wideparen}{\overparen}} % arc de cercle

\begin{document}

\pagegarde{L'aide au développement, les échanges commerciaux et la bourse des valeurs}{Allemand}{Tle A}{Chapitre 2 : Vie économique}{AL09}{2 h}{Cette leçon propose l'étude d'un texte sur l'aide au développement, les échanges commerciaux et la bourse des valeurs. Elle permet d'acquérir le vocabulaire économique essentiel, de maîtriser les prépositions avec le génitif et les participes en fonction d'épithète, et de produire un résumé ou un court commentaire.
\vspace{4pt}
\begin{multicols}{2}\small
\begin{itemize}
\item À la fin de cette leçon, je sais comprendre globalement puis en détail un texte allemand sur l'aide au développement, le commerce et la bourse.
\item Je sais employer correctement le vocabulaire thématique (die Entwicklungshilfe, der Handel, die Börse, die Aktie, der Export, der Import) avec articles et pluriels.
\item Je sais reconnaître et utiliser les prépositions suivies du génitif (wegen, trotz, während, innerhalb, anlässlich...).
\item Je sais former et employer les participes I et II en fonction d'épithète avec la bonne déclinaison.
\item Je sais conjuguer et reconnaître les temps composés (Perfekt, Plusquamperfekt, Futur II) dans un texte économique.
\item Je sais répondre par écrit à des questions de compréhension en phrases complètes.
\end{itemize}\end{multicols}}

\begin{sommaire}
\begin{multicols}{2}
\begin{enumerate}
\item Texte support et première compréhension
\item Lexique thématique : trois champs
\item Grammaire 1 : les prépositions avec le génitif
\item Grammaire 2 : les participes en fonction d'épithète
\item Conjugaison : révision des temps composés
\item Méthode du commentaire et production guidée
\item Activité d'intégration
\item Méthodes et exercices résolus
\item Exercices
\item Corrigés des exercices
\item Fiche bilan
\end{enumerate}
\end{multicols}
\end{sommaire}

\begin{objectifs}
\begin{itemize}
\item À la fin de cette leçon, je sais comprendre globalement puis en détail un texte allemand sur l'aide au développement, le commerce et la bourse.
\item Je sais employer correctement le vocabulaire thématique (die Entwicklungshilfe, der Handel, die Börse, die Aktie, der Export, der Import) avec articles et pluriels.
\item Je sais reconnaître et utiliser les prépositions suivies du génitif (wegen, trotz, während, innerhalb, anlässlich...).
\item Je sais former et employer les participes I et II en fonction d'épithète avec la bonne déclinaison.
\item Je sais conjuguer et reconnaître les temps composés (Perfekt, Plusquamperfekt, Futur II) dans un texte économique.
\item Je sais répondre par écrit à des questions de compréhension en phrases complètes.
\item Je sais rédiger un résumé structuré et un court paragraphe d'opinion sur un sujet économique.
\item Je sais appliquer la méthode du commentaire de texte exigée au Bac.
\end{itemize}
\end{objectifs}

\begin{prerequis}
\begin{itemize}
\item La déclinaison de l'adjectif après article défini et indéfini (Première)
\item Le Perfekt et le Plusquamperfekt avec haben/sein (Seconde/Première)
\item La formation du Partizip I et du Partizip II (Première)
\item Le génitif des articles et des noms (der, die, das + -s/-es)
\item Le vocabulaire général de l'économie vu en Première (die Wirtschaft, das Geld, die Firma)
\end{itemize}
\end{prerequis}

\newpage

\coursec{Texte support et première compréhension}

\courssub{Situation d'accroche et problématique}

Imaginez la scène dans votre classe de Terminale à Yaoundé. Pendant la récréation, un élève lit à voix haute un article de presse : \all{„Deutschland erhöht seine Entwicklungshilfe für Kamerun“} (« L'Allemagne augmente son aide au développement pour le Cameroun »). Un autre raconte que son oncle exporte du cacao vers le port de Hambourg. Un troisième a entendu à la radio que le DAX a gagné 2 \% à la Bourse de Francfort. Trois informations, trois réalités économiques différentes : \cle{l'aide au développement}, \cle{les échanges commerciaux} et \cle{la bourse des valeurs}. Pourtant, ces trois réalités relient concrètement le Cameroun aux pays germanophones.

\textbf{Problématique :} comment comprendre un texte économique en allemand et comment parler correctement de ces trois réalités ? C'est l'objet de cette leçon. Nous allons lire un texte de presse, le comprendre globalement puis en détail, et apprendre à répondre à des questions comme à l'examen du Baccalauréat.

\courssub{Lecture du texte support}

\begin{methode}[Lire un texte économique en quatre étapes]
\begin{enumerate}
\item \textbf{Lire le titre et prévoir le contenu.} Le titre \all{„Wirtschaftliche Brücken“} (« ponts économiques ») annonce un texte sur les relations économiques entre deux pays.
\item \textbf{Première lecture rapide} (silencieuse, 2 minutes) pour saisir l'idée générale, sans s'arrêter sur les mots inconnus.
\item \textbf{Deuxième lecture attentive} avec surlignage : souligner les chiffres, les acteurs (pays, entreprises, investisseurs) et les mots-clés du thème.
\item \textbf{Répondre aux questions en reformulant} avec ses propres mots, jamais en recopiant des phrases entières du texte.
\end{enumerate}
\end{methode}

Lisez maintenant le texte d'abord silencieusement, puis à voix haute après le professeur. Faites attention à la prononciation : \all{Entwicklungshilfe} se prononce « ènt-vi-kloungs-hil-fé », \all{Zusammenarbeit} « tsou-za-men-ar-baït », \all{Börse} « beur-zé ».

\begin{textebox}[Wirtschaftliche Brücken zwischen Deutschland und Kamerun]
Deutschland gehört zu den wichtigsten Handelspartnern Kameruns. Trotz der großen Entfernung ist die wirtschaftliche Zusammenarbeit zwischen den beiden Ländern eng. Während der letzten Jahre hat Deutschland seine Entwicklungshilfe für Kamerun deutlich erhöht: Wegen des schnellen Bevölkerungswachstums werden in vielen Regionen Schulen, Krankenhäuser und Straßen gebaut. Diese Projekte schaffen Arbeitsplätze und verbessern das Leben der Menschen, besonders auf dem Land.

Auch der Handel zwischen den beiden Staaten ist lebendig. Kamerun exportiert vor allem Rohstoffe wie Kakao, Kaffee und Holz nach Europa; im Gegenzug importiert es Maschinen, Autos und Medikamente aus Deutschland. Der Hafen von Douala spielt dabei eine zentrale Rolle, denn von hier aus reisen die kamerunischen Waren nach Hamburg.

An der Börse in Frankfurt spielen diese Waren ebenfalls eine wichtige Rolle: Die steigenden Kurse vieler Aktien zeigen das wachsende Interesse der Anleger an Afrika. Viele deutsche Unternehmen investieren wegen der günstigen Preise in kamerunische Projekte, zum Beispiel in die Landwirtschaft und in erneuerbare Energien. Die von beiden Ländern unterzeichneten Verträge sollen die Zusammenarbeit in den kommenden Jahren noch verstärken. Experten sind sich einig: Wenn der Handel fair bleibt, profitieren beide Seiten von dieser wirtschaftlichen Brücke zwischen Europa und Afrika.
\end{textebox}

\textbf{Glossaire marginal} (à recopier dans le cahier de vocabulaire) :

\begin{center}
\begin{tabularx}{\linewidth}{|X|X|}\hline
\rowcolor{popblueL} \textbf{Deutsch} & \textbf{Français} \\ \hline
die Zusammenarbeit (nur Sg.) & la coopération, la collaboration \\ \hline
der Handelspartner, - & le partenaire commercial \\ \hline
die Entfernung, -en & la distance \\ \hline
die Entwicklungshilfe (nur Sg.) & l'aide au développement \\ \hline
das Bevölkerungswachstum (nur Sg.) & la croissance démographique \\ \hline
der Rohstoff, -e & la matière première \\ \hline
der Hafen, “- & le port \\ \hline
die Ware, -n & la marchandise \\ \hline
der Kurs, -e & le cours (boursier) \\ \hline
steigen / sinken (stieg, ist gestiegen) & monter / baisser \\ \hline
das Unternehmen, - & l'entreprise \\ \hline
die Investition, -en & l'investissement \\ \hline
der Anleger, - & l'investisseur (en bourse) \\ \hline
unterzeichnen & signer (un contrat) \\ \hline
\end{tabularx}
\end{center}

\courssub{Repérage de la structure du texte}

Un bon lecteur identifie d'abord le plan du texte. Ici, le texte est construit en trois parties, qui correspondent exactement aux trois thèmes de la leçon. Observez l'organigramme suivant :

\begin{popfigure}
\begin{tikzpicture}[
node distance=0.6cm and 1.2cm,
main/.style={rectangle, rounded corners, draw=popdark, fill=popblueL, align=center, font=\bfseries, minimum width=4.2cm, minimum height=1cm},
part/.style={rectangle, rounded corners, draw=popdark, fill=poporangeL, align=center, minimum width=3.6cm, minimum height=1.4cm},
key/.style={rectangle, rounded corners, draw=popgreen, fill=popgreenL, align=center, font=\small, minimum width=3.6cm}
]
\node[main, align=center] (t){Texte : „Wirtschaftliche Brücken“\\ Deutschland -- Kamerun};
\node[part, below left=1.1cm and -1.5cm of t, align=center] (p1){1. Entwicklungshilfe\\ (l'aide au développement)};
\node[part, below=1.1cm of t, align=center] (p2){2. Handel\\ (les échanges commerciaux)};
\node[part, below right=1.1cm and -1.5cm of t, align=center] (p3){3. Börse\\ (la bourse des valeurs)};
\node[key, below=0.35cm of p1, align=center] (k1){Mot-clé :\\ \all{Schulen, Krankenhäuser bauen}};
\node[key, below=0.35cm of p2, align=center] (k2){Mot-clé :\\ \all{exportieren / importieren}};
\node[key, below=0.35cm of p3, align=center] (k3){Mot-clé :\\ \all{steigende Kurse}};
\draw[-{Stealth}, thick, popblue] (t) -- (p1);
\draw[-{Stealth}, thick, popblue] (t) -- (p2);
\draw[-{Stealth}, thick, popblue] (t) -- (p3);
\draw[-{Stealth}, thick, popgreen] (p1) -- (k1);
\draw[-{Stealth}, thick, popgreen] (p2) -- (k2);
\draw[-{Stealth}, thick, popgreen] (p3) -- (k3);
\end{tikzpicture}
\legende{Organigramme : les trois parties du texte et leur mot-clé}
\end{popfigure}

\begin{aretenir}[Le plan du texte]
\begin{itemize}
\item \textbf{Partie 1} (phrases 1 à 5) : \all{die Entwicklungshilfe} — l'Allemagne finance des écoles, des hôpitaux et des routes.
\item \textbf{Partie 2} (phrases 6 à 8) : \all{der Handel} — le Cameroun exporte des matières premières et importe des produits manufacturés.
\item \textbf{Partie 3} (phrases 9 à 12) : \all{die Börse} — les cours montent, les entreprises allemandes investissent, des contrats sont signés.
\end{itemize}
\end{aretenir}

\courssub{Compréhension globale}

Avant les détails, répondez aux trois questions fondamentales de tout commentaire de texte.

\textbf{1. Thème (Worum geht es ?)} Le texte parle des relations économiques entre l'Allemagne et le Cameroun : aide au développement, commerce et investissements en bourse.

\textbf{2. Source probable (Woher stammt der Text ?)} Le ton informatif, les faits présentés objectivement et le vocabulaire économique indiquent un \cle{article de presse}, par exemple un quotidien allemand comme la \all{Frankfurter Allgemeine Zeitung}, ou un magazine économique.

\textbf{3. Intention de l'auteur (Was will der Autor ?)} L'auteur veut \cle{informer} le lecteur sur la coopération germano-camerounaise et montrer, avec une conclusion optimiste (\all{„profitieren beide Seiten“}), que cette relation peut être profitable aux deux pays.

\begin{exemplebox}[Deux phrases clés du texte, traduites]
\all{Trotz der großen Entfernung ist die Zusammenarbeit eng.} = Malgré la grande distance, la coopération est étroite.\par
\all{Die steigenden Kurse zeigen das Interesse der Anleger.} = Les cours en hausse montrent l'intérêt des investisseurs.
\end{exemplebox}

\courssub{Compréhension détaillée : chiffres, acteurs, causes et conséquences}

Repérez dans le texte les informations précises. Un tableau vous aide à les classer :

\begin{center}
\begin{tabularx}{\linewidth}{|X|X|}\hline
\rowcolor{popblueL} \textbf{Question} & \textbf{Réponse trouvée dans le texte} \\ \hline
Acteurs (Wer ?) & l'Allemagne, le Cameroun, les entreprises allemandes, les investisseurs (\all{die Anleger}) \\ \hline
Exportations du Cameroun & cacao, café, bois (\all{Kakao, Kaffee, Holz}) \\ \hline
Importations du Cameroun & machines, voitures, médicaments (\all{Maschinen, Autos, Medikamente}) \\ \hline
Cause de l'aide allemande & la croissance démographique rapide (\all{wegen des schnellen Bevölkerungswachstums}) \\ \hline
Conséquence des projets & des emplois créés et une vie améliorée (\all{schaffen Arbeitsplätze}) \\ \hline
Cause des investissements & les prix avantageux (\all{wegen der günstigen Preise}) \\ \hline
Lieu stratégique & le port de Douala, vers Hambourg \\ \hline
\end{tabularx}
\end{center}

Remarquez la logique \cle{cause $\rightarrow$ conséquence} exprimée par les prépositions \all{wegen} et \all{trotz} suivies du génitif : \all{wegen des Bevölkerungswachstums} (à cause de la croissance démographique), \all{trotz der großen Entfernung} (malgré la grande distance). Nous étudierons ces prépositions en détail dans la section grammaire.

\courssub{Questions de compréhension}

\begin{exercice}[Fragen zum Text]
Répondez en allemand, par des \textbf{phrases complètes}, en reformulant.
\begin{enumerate}
\item Wer gehört zu den wichtigsten Handelspartnern Kameruns ?
\item Was baut Deutschland in Kamerun, und warum ?
\item Welche Rohstoffe exportiert Kamerun nach Europa ?
\item Wie reisen die kamerunischen Waren nach Hamburg ?
\item Warum investieren viele deutsche Unternehmen in Kamerun ?
\item Was zeigen die steigenden Kurse an der Frankfurter Börse ?
\end{enumerate}
\end{exercice}

\begin{corrige}[Exercice : Fragen zum Text]
\begin{enumerate}
\item \all{Deutschland gehört zu den wichtigsten Handelspartnern Kameruns.} (L'Allemagne fait partie des partenaires commerciaux les plus importants du Cameroun.)
\item \all{Deutschland baut Schulen, Krankenhäuser und Straßen, weil die Bevölkerung schnell wächst.} (L'Allemagne construit des écoles, des hôpitaux et des routes parce que la population augmente vite.)
\item \all{Kamerun exportiert Kakao, Kaffee und Holz.} (Le Cameroun exporte du cacao, du café et du bois.)
\item \all{Sie reisen mit dem Schiff vom Hafen in Douala nach Hamburg.} (Elles voyagent par bateau du port de Douala à Hambourg.)
\item \all{Sie investieren in Kamerun, weil die Preise dort günstig sind.} (Elles investissent au Cameroun parce que les prix y sont avantageux.)
\item \all{Sie zeigen, dass sich die Anleger immer mehr für Afrika interessieren.} (Ils montrent que les investisseurs s'intéressent de plus en plus à l'Afrique.)
\end{enumerate}
\end{corrige}

\courssub{Entraînement Bac : Vrai ou Faux justifié}

Au Baccalauréat, l'exercice \emph{Vrai/Faux} exige une \textbf{justification par une citation exacte du texte}, entre guillemets allemands. Voici un modèle résolu.

\begin{exoresolu}[Richtig oder falsch ? Begründen Sie mit einem Zitat]
Énoncé : \all{Kamerun importiert Holz aus Deutschland.} — Vrai ou faux ? Justifiez par une citation.
\tcblower
\textbf{Solution :} \emph{Falsch.} Le texte dit le contraire : le bois fait partie des exportations du Cameroun, pas de ses importations.\par
Justification (citation exacte) : \all{„Kamerun exportiert vor allem Rohstoffe wie Kakao, Kaffee und Holz nach Europa; im Gegenzug importiert es Maschinen, Autos und Medikamente aus Deutschland.“}\par
Méthode : on cite la phrase complète qui contient la preuve, sans la modifier, et l'on souligne mentalement le mot décisif (\all{exportiert}).
\end{exoresolu}

\begin{exercice}[Richtig oder falsch ? — À vous de jouer]
Justifiez chaque réponse par une citation du texte.
\begin{enumerate}
\item \all{Die wirtschaftliche Zusammenarbeit zwischen Deutschland und Kamerun ist schwach.}
\item \all{Die deutsche Entwicklungshilfe schafft Arbeitsplätze.}
\item \all{Die Anleger interessieren sich nicht für Afrika.}
\item \all{Beide Länder haben Verträge unterzeichnet.}
\end{enumerate}
\end{exercice}

\begin{corrige}[Exercice : Richtig oder falsch ?]
\begin{enumerate}
\item \emph{Falsch.} Citation : \all{„Trotz der großen Entfernung ist die wirtschaftliche Zusammenarbeit zwischen den beiden Ländern eng.“}
\item \emph{Richtig.} Citation : \all{„Diese Projekte schaffen Arbeitsplätze und verbessern das Leben der Menschen.“}
\item \emph{Falsch.} Citation : \all{„Die steigenden Kurse vieler Aktien zeigen das wachsende Interesse der Anleger an Afrika.“}
\item \emph{Richtig.} Citation : \all{„Die von beiden Ländern unterzeichneten Verträge sollen die Zusammenarbeit in den kommenden Jahren noch verstärken.“}
\end{enumerate}
\end{corrige}

\begin{attention}[Pièges à éviter au Bac]
\begin{itemize}
\item \textbf{Ne traduisez jamais mot à mot.} \all{„Die Zusammenarbeit ist eng“} ne se traduit pas par « la coopération est étroite » au sens physique : \all{eng} signifie ici « étroite, intense » au sens figuré.
\item \textbf{Une réponse recopiée du texte sans reformulation est pénalisée.} Aux questions de compréhension, reformulez avec vos propres mots ; la citation mot à mot n'est acceptée que dans l'exercice Vrai/Faux.
\item \textbf{Faux amis :} \all{die Aktie} = l'action (en bourse), pas « l'acte » ; \all{der Kurs} = le cours boursier, pas seulement « le cours » scolaire ; \all{die Ware} = la marchandise, pas « la guerre » (\all{der Krieg}).
\item \textbf{Majuscules :} tous les noms prennent la majuscule : \all{die Börse, der Handel, das Unternehmen}.
\end{itemize}
\end{attention}

\begin{savaistu}[Le DAX et le port de Hambourg]
Le \all{DAX} (\all{Deutscher Aktienindex}) est l'indice principal de la Bourse de Francfort : il regroupe les trente plus grandes entreprises allemandes, puis quarante depuis 2021. Quand la radio annonce « le DAX a gagné 2 \% », cela signifie que la valeur moyenne de ces entreprises a augmenté. Le port de \all{Hamburg}, sur l'Elbe, est l'une des premières portes d'entrée du cacao africain en Europe : une partie du cacao camerounais y transite avant d'être transformé en chocolat. Voilà pourquoi l'oncle de notre élève expédie ses fèves vers Hambourg !
\end{savaistu}

\begin{experience}[À l'oral : présentez le texte en une minute]
Par groupes de trois : l'élève A résume la partie \all{Entwicklungshilfe}, l'élève B la partie \all{Handel}, l'élève C la partie \all{Börse}, chacun en deux phrases maximum, sans regarder le texte. Commencez par : \all{„Der Text handelt von …“} ou \all{„Im ersten Teil geht es um …“}. La classe vérifie ensuite que les trois mots-clés de l'organigramme ont bien été cités.
\end{experience}

\coursec{Lexique thématique : trois champs}

Dans la section précédente, nous avons lu un texte sur l'aide au développement, les échanges commerciaux et la bourse, et nous en avons dégagé le sens global. Mais pour parler d'économie en allemand — et pour réussir les épreuves de compréhension et d'expression du baccalauréat — il ne suffit pas de comprendre : il faut posséder un vocabulaire précis, avec ses articles, ses pluriels et ses verbes. C'est l'objet de cette section : nous allons organiser le lexique du texte en trois champs sémantiques — \all{die Entwicklungshilfe}, \all{der Handel}, \all{die Börse} — puis nous apprendrons à construire des familles de mots, technique redoutable pour enrichir rapidement son vocabulaire.

\begin{aretenir}[Règle d'or du vocabulaire allemand]
On ne mémorise jamais un nom allemand seul. On l'apprend toujours avec \cle{son article} (der, die, das) et \cle{son pluriel}. Exemple : \all{die Aktie, -n} signifie « l'action », féminin, pluriel \all{die Aktien}. Sans l'article, impossible de décliner correctement ; sans le pluriel, impossible de parler au pluriel.
\end{aretenir}

\courssub{Champ 1 : die Entwicklungshilfe (l'aide au développement)}

\begin{definition}[Vocabulaire de l'aide au développement]
\begin{itemize}
\item \all{die Entwicklungshilfe, -n} : l'aide au développement (souvent employée au singulier)
\item \all{das Entwicklungsland, \textsuperscript{¨}-er} : le pays en voie de développement (pluriel \all{die Entwicklungsländer})
\item \all{das Projekt, -e} : le projet
\item \all{die Zusammenarbeit} (Sg.) : la coopération, la collaboration (pas de pluriel, comme en français « la coopération »)
\item \all{die Investition, -en} : l'investissement
\item \all{der Kredit, -e} : le crédit, le prêt
\item \all{die Schuld, -en} : la dette (aussi : la faute, la culpabilité)
\item \all{helfen — hilft — half — hat geholfen} (+ datif) : aider
\item \all{unterstützen — hat unterstützt} : soutenir (préfixe inséparable \all{unter-}, pas de \all{ge-} au participe)
\item \all{fördern — hat gefördert} : promouvoir, encourager
\end{itemize}
\end{definition}

\begin{exemplebox}[Le champ lexical en contexte]
\all{Deutschland unterstützt viele Projekte in den Entwicklungsländern.} « L'Allemagne soutient de nombreux projets dans les pays en voie de développement. »

\all{Die Regierung hat dem Land einen Kredit gegeben.} « Le gouvernement a accordé un crédit au pays. » (Attention : \all{helfen} se construit avec le datif : \all{Ich helfe dem Land.} « J'aide le pays. »)

\all{Diese Investition fördert die Zusammenarbeit zwischen Kamerun und Deutschland.} « Cet investissement favorise la coopération entre le Cameroun et l'Allemagne. »

\all{Viele Länder haben hohe Schulden.} « Beaucoup de pays ont des dettes élevées. »
\end{exemplebox}

\begin{attention}[Faux ami : die Schuld]
\all{die Schuld} signifie à la fois « la dette » (sens économique : \all{Schulden haben} « avoir des dettes », souvent au pluriel) et « la faute, la culpabilité » (\all{Es ist meine Schuld.} « C'est ma faute. »). Le contexte économique lève l'ambiguïté. Ne traduis pas « la faute » par \all{die Schuld} dans un texte de bourse sans vérifier le sens !
\end{attention}

\courssub{Champ 2 : der Handel (le commerce)}

\begin{definition}[Vocabulaire des échanges commerciaux]
\begin{itemize}
\item \all{der Handel} (Sg.) : le commerce
\item \all{der Handelspartner, -} : le partenaire commercial (pluriel identique : \all{die Handelspartner})
\item \all{der Export, -e} : l'exportation ; \all{der Import, -e} : l'importation
\item \all{exportieren — hat exportiert} : exporter ; \all{importieren — hat importiert} : importer
\item \all{der Rohstoff, -e} : la matière première
\item \all{die Ware, -n} : la marchandise
\item \all{der Zoll, \textsuperscript{¨}-e} : la douane (pluriel \all{die Zölle} : les droits de douane)
\item \all{der Vertrag, \textsuperscript{¨}-e} : le contrat, le traité (pluriel \all{die Verträge})
\item \all{unterzeichnen — hat unterzeichnet} : signer (un document officiel ; préfixe inséparable \all{unter-}, pas de \all{ge-})
\end{itemize}
\end{definition}

\begin{exemplebox}[Exemples traduits]
\all{Kamerun exportiert Kakao und importiert Maschinen.} « Le Cameroun exporte du cacao et importe des machines. »

\all{Deutschland ist ein wichtiger Handelspartner Kameruns.} « L'Allemagne est un partenaire commercial important du Cameroun. » (Génitif \all{Kameruns} après le nom composé).

\all{Die beiden Länder haben einen Vertrag unterzeichnet.} « Les deux pays ont signé un contrat. »

\all{Rohstoffe wie Öl und Kaffee spielen im Handel eine große Rolle.} « Les matières premières comme le pétrole et le café jouent un grand rôle dans le commerce. »

\all{Die Waren müssen durch den Zoll.} « Les marchandises doivent passer par la douane. » (Accusatif après \all{durch}).
\end{exemplebox}

\begin{attention}[Verbes en \all{-ieren} et préfixes inséparables : pas de \all{ge-} au participe]
Les verbes en \all{-ieren} (\all{exportieren}, \all{importieren}, \all{investieren}, \all{studieren}, \all{reparieren}...) forment leur participe passé \cle{sans ge-} : \all{exportiert}, \all{importiert}, \all{investiert}. Erreur fréquente des francophones : *\all{geexportiert}* est faux. \\
De même, les verbes à préfixe inséparable comme \all{unterzeichnen}, \all{unterstützen}, \all{fördern} (attention, \all{fördern} n'a pas de préfixe, mais \all{unterstützen} a \all{unter-} inséparable) n'ont pas de \all{ge-} : \all{unterzeichnet}, \all{unterstützt}, \all{gefördert}.
\end{attention}

\courssub{Champ 3 : die Börse (la bourse)}

\begin{definition}[Vocabulaire de la bourse des valeurs]
\begin{itemize}
\item \all{die Börse, -n} : la bourse (des valeurs)
\item \all{die Aktie, -n} : l'action (titre boursier)
\item \all{der Kurs, -e} : le cours (d'une action, d'une devise)
\item \all{der Anleger, -} : l'investisseur, le particulier qui place son argent (pluriel \all{die Anleger})
\item \all{der Gewinn, -e} : le gain, le bénéfice ; \all{der Verlust, -e} : la perte
\item \all{steigen — stieg — ist gestiegen} : monter, augmenter (verbe fort, auxiliaire \all{sein})
\item \all{sinken — sank — ist gesunken} : baisser, diminuer (verbe fort, auxiliaire \all{sein})
\item \all{fallen — fällt — fiel — ist gefallen} : tomber, chuter (verbe fort, auxiliaire \all{sein})
\item \all{investieren — hat investiert} : investir (verbe faible en \all{-ieren}, auxiliaire \all{haben})
\end{itemize}
\end{definition}

\begin{propriete}[Les verbes de variation économique se conjuguent avec \all{sein}]
\all{steigen}, \all{sinken} et \all{fallen} expriment un changement d'état (une variation de niveau, de cours, de prix). Comme tous les verbes de déplacement ou de transformation d'état, ils forment leurs temps composés avec l'auxiliaire \cle{sein} :

\all{Die Kurse sind gestiegen.} « Les cours ont monté. »

\all{Der Gewinn ist gesunken.} « Le bénéfice a baissé. »

\all{Die Aktie ist um fünf Prozent gefallen.} « L'action a chuté de cinq pour cent. »

En revanche, \all{investieren} (action volontaire, pas changement d'état automatique) se conjugue avec \cle{haben} : \all{Der Anleger hat investiert.} « L'investisseur a investi. »
\end{propriete}

\begin{exemplebox}[Exemples traduits]
\all{Das Unternehmen hat in ein Projekt investiert.} « L'entreprise a investi dans un projet. » (\all{investieren in} + accusatif)

\all{Viele Anleger kaufen Aktien an der Börse.} « Beaucoup d'investisseurs achètent des actions à la bourse. »

\all{Wer Aktien kauft, kann Gewinne, aber auch Verluste machen.} « Celui qui achète des actions peut faire des gains, mais aussi des pertes. »
\end{exemplebox}

\begin{attention}[Faux amis : \all{der Kurs} et \all{die Aktie}]
\all{der Kurs} désigne le cours \emph{boursier} (le prix d'une action) ou un cours de langue (\all{ein Deutschkurs}) ; le cours scolaire, la leçon, se dit \all{der Unterricht}. \\
\all{die Aktie} est l'action au sens \emph{financier} ; l'action au sens général (ce qu'on fait, acte) se dit \all{die Handlung}. Ne dis jamais *\all{die Aktie des Helden}* pour « l'action du héros » !
\end{attention}

\begin{savaistu}[La Bourse de Francfort et le DAX]
La \all{Frankfurter Wertpapierbörse}, la Bourse de Francfort, est la plus grande place boursière d'Allemagne et l'une des plus importantes d'Europe. Son indice vedette s'appelle le \all{DAX} (\all{Deutscher Aktienindex}, « indice allemand des actions ») : il regroupe les grandes entreprises cotées allemandes. Quand un journaliste allemand dit \all{Der DAX ist gestiegen}, cela signifie : « Le DAX a monté », c'est-à-dire que la valeur moyenne de ces entreprises a augmenté.
\end{savaistu}

\courssub{Familles de mots : un mot appris, trois mots gagnés}

L'allemand construit son vocabulaire par dérivation et composition. Observer les familles de mots permet de deviner le sens de mots inconnus dans un texte.

\begin{center}
\begin{tabularx}{\linewidth}{|X|X|X|}
\hline
\rowcolor{popblueL} Verbe / base & Mots dérivés & Sens \\
\hline
\all{handeln} (commercer, agir) & \all{der Handel}, \all{der Händler}, \all{der Handelspartner} & le commerce, le commerçant, le partenaire commercial \\
\hline
\all{entwickeln} (développer) & \all{die Entwicklung}, \all{die Entwicklungshilfe}, \all{das Entwicklungsland} & le développement, l'aide au développement, le pays en voie de développement \\
\hline
\all{steigen} (monter) & \all{der Anstieg}, \all{steigend} & la hausse, croissant / en hausse \\
\hline
\all{investieren} (investir) & \all{die Investition}, \all{der Investor} & l'investissement, l'investisseur \\
\hline
\all{gewinnen} (gagner) & \all{der Gewinn}, \all{der Gewinner} & le gain, le gagnant \\
\hline
\end{tabularx}
\end{center}

Remarque la logique des mots composés : \all{Entwicklung} + \all{Hilfe} = \all{Entwicklungshilfe} (l'aide au développement) ; \all{Handel} + \all{Partner} = \all{Handelspartner}. En allemand, c'est toujours le \cle{dernier élément} qui porte le genre et le sens principal : \all{die Hilfe} est féminin, donc \all{die Entwicklungshilfe} est féminin.

\begin{popfigure}
\begin{tikzpicture}[mindmap, grow cyclic, every node/.style={concept, align=center, font=\small},
root concept/.append style={concept color=popblue, font=\bfseries},
level 1/.append style={level distance=3.4cm, sibling angle=120, concept color=poporange},
level 2/.append style={level distance=2.6cm, concept color=popgreenL, font=\footnotesize}]
\node[root concept]{WIRTSCHAFT}
 child { node {die Entwicklungshilfe}
   child { node {das Projekt} }
   child { node {der Kredit} }
   child { node {die Schuld} }
   child { node {fördern} } }
 child { node {der Handel}
   child { node {der Export} }
   child { node {der Import} }
   child { node {die Ware} }
   child { node {der Zoll} } }
 child { node {die Börse}
   child { node {die Aktie} }
   child { node {der Kurs} }
   child { node {der Anleger} }
   child { node {der Gewinn} } };
\end{tikzpicture}
\legende{Carte mentale du lexique économique : trois champs, quatre mots-clés par champ.}
\end{popfigure}

\courssub{Expressions utiles à réemployer}

\begin{definition}[Trois expressions à connaître par cœur]
\begin{itemize}
\item \all{im Gegenzug} : en contrepartie. \all{Kamerun liefert Rohstoffe; im Gegenzug bekommt es Maschinen.} « Le Cameroun fournit des matières premières ; en contrepartie, il reçoit des machines. »
\item \all{eine Rolle spielen} : jouer un rôle. \all{Der Handel spielt eine wichtige Rolle in der Wirtschaft.} « Le commerce joue un rôle important dans l'économie. »
\item \all{in etwas investieren} (+ accusatif) : investir dans quelque chose. \all{Die Firma investiert in neue Maschinen.} « L'entreprise investit dans de nouvelles machines. »
\end{itemize}
\end{definition}

\begin{exoresolu}[Complète avec le mot qui convient]
Énoncé — Complète chaque phrase avec le mot correct de la liste : \all{die Aktie}, \all{der Kredit}, \all{der Rohstoff}, \all{der Vertrag}, \all{der Gewinn}.

1. \all{Kakao ist ein wichtiger \ldots\ für Kamerun.}
2. \all{Die Bank gibt dem Bauern einen \ldots.}
3. \all{Die beiden Firmen haben einen \ldots\ unterzeichnet.}
4. \all{Der Anleger kauft eine \ldots\ an der Börse.}
5. \all{Am Ende des Jahres macht das Unternehmen einen großen \ldots.}
\tcblower
Solution détaillée :

1. \all{Kakao ist ein wichtiger \textbf{Rohstoff} für Kamerun.} « Le cacao est une matière première importante pour le Cameroun. » (\all{ein} exige un masculin ou un neutre ; le sens « matière première » impose \all{der Rohstoff}.)

2. \all{Die Bank gibt dem Bauern einen \textbf{Kredit}.} « La banque accorde un crédit au paysan. » (C'est la banque qui prête de l'argent ; \all{einen} + accusatif masculin.)

3. \all{Die beiden Firmen haben einen \textbf{Vertrag} unterzeichnet.} « Les deux entreprises ont signé un contrat. » (Le verbe \all{unterzeichnen} appelle un document : contrat ou traité.)

4. \all{Der Anleger kauft eine \textbf{Aktie} an der Börse.} « L'investisseur achète une action à la bourse. » (À la bourse, on achète des actions ; \all{eine} + féminin.)

5. \all{Am Ende des Jahres macht das Unternehmen einen großen \textbf{Gewinn}.} « À la fin de l'année, l'entreprise réalise un grand bénéfice. » (Expression figée : \all{einen Gewinn machen} « faire un bénéfice ».)
\end{exoresolu}

\begin{exercice}[À toi de jouer]
1. Traduis en allemand : « Le Cameroun exporte du café et importe des voitures. » / « Les cours sont montés. » / « L'Allemagne soutient un projet au Cameroun. »
2. Trouve la famille de mots : donne deux dérivés de \all{handeln} et deux de \all{entwickeln}.
3. Conjugue au Perfekt : \all{die Kurse (steigen)}, \all{das Unternehmen (investieren)}, \all{der Preis (fallen)}.
\end{exercice}

\begin{corrige}[Exercice « À toi de jouer »]
1. \all{Kamerun exportiert Kaffee und importiert Autos.} / \all{Die Kurse sind gestiegen.} (auxiliaire \all{sein} !) / \all{Deutschland unterstützt ein Projekt in Kamerun.}
2. \all{handeln} : \all{der Handel}, \all{der Händler} (ou \all{der Handelspartner}) ; \all{entwickeln} : \all{die Entwicklung}, \all{das Entwicklungsland} (ou \all{die Entwicklungshilfe}).
3. \all{Die Kurse sind gestiegen.} / \all{Das Unternehmen hat investiert.} / \all{Der Preis ist gefallen.}
\end{corrige}

\begin{aretenir}[L'essentiel du lexique]
Trois champs, trois verbes-images : \all{helfen/unterstützen} pour l'aide au développement, \all{exportieren/importieren} pour le commerce, \all{steigen/sinken/fallen} pour la bourse — ces trois derniers avec \cle{sein} au Perfekt. Chaque nom s'apprend avec article et pluriel ; chaque famille de mots multiplie ton vocabulaire. Tu es maintenant prêt à aborder la grammaire : les prépositions avec le génitif, très fréquentes dans les textes économiques.
\end{aretenir}

\coursec{Grammaire 1 : les prépositions avec le génitif}

Dans la section précédente, nous avons constitué notre lexique autour de trois champs : \all{die Entwicklungshilfe}, \all{der Handel} et \all{die Börse}. En relisant le texte support, vous avez peut-être remarqué que ces mots sont souvent introduits par de petits mots comme \all{wegen}, \all{trotz} ou \all{während}, et que le nom qui les suit change de forme : \all{des Bevölkerungswachstums}, \all{der letzten Jahre}... C'est exactement le phénomène que nous allons étudier maintenant : les \cle{prépositions suivies du génitif}, un outil indispensable pour lire et rédiger des textes économiques.

\courssub{Observation : que remarque-t-on dans le texte ?}

\begin{textebox}[Mini-texte d'observation : extrait recomposé]
\all{Trotz der großen Entfernung zwischen Europa und Afrika wächst der Handel ständig. Wegen des Bevölkerungswachstums brauchen viele Länder mehr Nahrungsmittel und mehr Schulen. Während der letzten Jahre hat Deutschland seine Entwicklungshilfe erhöht. Wegen der günstigen Preise kaufen viele Händler Kakao und Kaffee in Kamerun. Aufgrund der steigenden Kurse an der Börse investieren immer mehr Anleger ihr Geld in Aktien.}

\emph{Malgré la grande distance entre l'Europe et l'Afrique, le commerce ne cesse de croître. À cause de la croissance démographique, beaucoup de pays ont besoin de plus de denrées alimentaires et de plus d'écoles. Pendant ces dernières années, l'Allemagne a augmenté son aide au développement. En raison des prix avantageux, beaucoup de commerçants achètent du cacao et du café au Cameroun. En raison de la hausse des cours à la bourse, de plus en plus d'investisseurs placent leur argent dans des actions.}

\textbf{Glossaire :} die Entfernung, -en (la distance) ; das Bevölkerungswachstum (la croissance démographique) ; die Entwicklungshilfe, -n (l'aide au développement) ; günstig (avantageux) ; der Händler, - (le commerçant) ; der Kurs, -e (le cours, la cotation) ; der Anleger, - (l'investisseur) ; die Aktie, -n (l'action).
\end{textebox}

Observons attentivement les groupes soulignés par la structure :

\begin{center}
\begin{tabularx}{\linewidth}{|X|X|X|}
\hline
\rowcolor{popblueL} Préposition & Nom qui suit (forme) & Traduction française \\
\hline
\all{trotz} & \all{der großen Entfernung} (fém. sing.) & malgré la grande distance \\
\hline
\all{wegen} & \all{des Bevölkerungswachstums} (neutre sing.) & à cause de la croissance démographique \\
\hline
\all{während} & \all{der letzten Jahre} (pluriel) & pendant les dernières années \\
\hline
\all{wegen} & \all{der günstigen Preise} (pluriel) & à cause des prix avantageux \\
\hline
\all{aufgrund} & \all{der steigenden Kurse} (pluriel) & en raison de la hausse des cours \\
\hline
\end{tabularx}
\end{center}

\textbf{Deux constatations essentielles :}
\begin{enumerate}
\item Après ces prépositions, l'article n'est jamais \all{der/die/das} (nominatif) ni \all{den/die/das} (accusatif), mais \all{des} ou \all{der} : c'est le \cle{génitif}.
\item Les noms masculins et neutres au singulier prennent en plus une terminaison \all{-(e)s} : \all{das Wachstum} devient \all{des Wachstums}.
\end{enumerate}

\courssub{La règle : forme et liste des prépositions}

\begin{propriete}[Prépositions + génitif : la règle fondamentale]
Après les prépositions \all{wegen} (à cause de), \all{trotz} (malgré), \all{während} (pendant), \all{innerhalb} (à l'intérieur de / en l'espace de), \all{außerhalb} (en dehors de), \all{aufgrund} (en raison de), \all{anlässlich} (à l'occasion de) et \all{statt} / \all{anstelle} (au lieu de), le nom se met au \cle{GÉNITIF}.

\all{trotz des hohen Preises} = malgré le prix élevé ; \all{während des Krieges} = pendant la guerre ; \all{wegen der Krise} = à cause de la crise.
\end{propriete}

\begin{aretenir}[Rappel : les formes du génitif]
\begin{center}
\begin{tabular}{|l|l|l|l|}
\hline
\rowcolor{popblueL} & Masculin & Neutre & Féminin / Pluriel \\
\hline
Article défini & \all{des} + nom en \all{-(e)s} & \all{des} + nom en \all{-(e)s} & \all{der} + nom inchangé \\
\hline
Exemple & \all{des Marktes} & \all{des Landes} & \all{der Hilfe} / \all{der Länder} \\
\hline
Article indéfini & \all{eines} + nom en \all{-(e)s} & \all{eines} + nom en \all{-(e)s} & \all{einer} + nom inchangé \\
\hline
Exemple & \all{eines Vertrags} & \all{eines Unternehmens} & \all{einer Bank} \\
\hline
Déclinaison forte (sans article) & \all{hohen Preises} & \all{starken Wachstums} & \all{guter Qualität} \\
\hline
\end{tabular}
\end{center}

Terminaison du nom masculin/neutre au singulier : \all{-es} pour les monosyllabes (\all{des Marktes}, \all{des Wachstums}, \all{des Jahres}) ; \all{-s} pour les polysyllabes terminés par une syllabe atone (\all{des Anlegers}, \all{des Handels}, \all{des Unternehmens}). Les noms propres prennent \all{-s} sans article : \all{wegen Kameruns}, \all{trotz Deutschlands}.
\end{aretenir}

\courssub{Les huit prépositions au programme : sens et exemples}

\begin{popfigure}
\begin{center}
\begin{tikzpicture}[
  prep/.style={rectangle, rounded corners=3pt, draw=popblue, fill=popblueL, align=center, font=\small\bfseries, minimum width=2.6cm, minimum height=0.9cm},
  sens/.style={rectangle, rounded corners=3pt, draw=poporange, fill=poporangeL, align=center, font=\small, minimum width=2.8cm, minimum height=0.9cm},
  ex/.style={rectangle, rounded corners=3pt, draw=popgreen, fill=popgreenL, align=left, font=\small, minimum width=5.2cm, minimum height=0.9cm},
  trad/.style={rectangle, rounded corners=3pt, draw=poppurple, fill=poppurpleL, align=left, font=\small\itshape, minimum width=4.6cm, minimum height=0.9cm},
  head/.style={font=\small\bfseries, text=popdark}
]
\node[head] at (0,0.6) {Préposition};
\node[head] at (3.4,0.6) {Sens};
\node[head] at (7.6,0.6) {Exemple};
\node[head] at (12.6,0.6) {Traduction};
\node[prep] at (0,0) {wegen};
\node[sens] at (3.4,0) {à cause de};
\node[ex] at (7.6,0) {wegen des hohen Preises};
\node[trad] at (12.6,0) {à cause du prix élevé};
\node[prep] at (0,-1.1) {trotz};
\node[sens] at (3.4,-1.1) {malgré};
\node[ex] at (7.6,-1.1) {trotz der Krise};
\node[trad] at (12.6,-1.1) {malgré la crise};
\node[prep] at (0,-2.2) {während};
\node[sens] at (3.4,-2.2) {pendant};
\node[ex] at (7.6,-2.2) {während des Jahres};
\node[trad] at (12.6,-2.2) {pendant l'année};
\node[prep] at (0,-3.3) {aufgrund};
\node[sens] at (3.4,-3.3) {en raison de};
\node[ex] at (7.6,-3.3) {aufgrund der steigenden Kurse};
\node[trad] at (12.6,-3.3) {en raison de la hausse des cours};
\end{tikzpicture}
\end{center}
\legende{Quatre prépositions à génitif fréquentes dans les textes économiques}
\end{popfigure}

\begin{exemplebox}[Exemples traduits, classés par fonction]
\textbf{La cause} (\all{wegen}, \all{aufgrund}) :

\all{Wegen des Bevölkerungswachstums werden Schulen gebaut.} « À cause de la croissance démographique, on construit des écoles. »

\all{Aufgrund der steigenden Kurse investieren viele Anleger.} « En raison de la hausse des cours, beaucoup d'investisseurs investissent. »

\textbf{La concession} (\all{trotz}) :

\all{Trotz der großen Entfernung exportiert Kamerun seinen Kakao nach Europa.} « Malgré la grande distance, le Cameroun exporte son cacao vers l'Europe. »

\textbf{Le temps} (\all{während}, \all{innerhalb}) :

\all{Während der letzten Jahre hat Deutschland seine Hilfe erhöht.} « Pendant ces dernières années, l'Allemagne a augmenté son aide. »

\all{Innerhalb eines Monats ist der Kurs der Aktie um zehn Prozent gestiegen.} « En l'espace d'un mois, le cours de l'action a augmenté de dix pour cent. »

\textbf{Le lieu} (\all{innerhalb}, \all{außerhalb}) :

\all{Außerhalb der Europäischen Union gelten andere Zölle.} « En dehors de l'Union européenne, d'autres droits de douane sont en vigueur. »

\textbf{L'occasion et le remplacement} (\all{anlässlich}, \all{statt}) :

\all{Anlässlich des Besuchs der Kanzlerin wurde ein neues Abkommen unterzeichnet.} « À l'occasion de la visite de la chancelière, un nouvel accord a été signé. »

\all{Statt eines Kredits bekam das Land eine direkte Hilfe.} « Au lieu d'un crédit, le pays a reçu une aide directe. »
\end{exemplebox}

\courssub{Comparaison français / allemand}

En français, « à cause de », « malgré », « pendant » sont simplement suivis de l'article normal : \emph{malgré le prix}, \emph{pendant la guerre}. En allemand, il faut \textbf{transformer} l'article au génitif et, pour les masculins et neutres, ajouter \all{-(e)s} au nom. C'est une étape supplémentaire que le francophone oublie presque toujours.

\begin{center}
\begin{tabularx}{\linewidth}{|X|X|X|}
\hline
\rowcolor{popblueL} Français & Deutsch & Remarque \\
\hline
à cause de la crise & \all{wegen der Krise} & fém. : \all{der}, nom inchangé \\
\hline
malgré le prix élevé & \all{trotz des hohen Preises} & masc. : \all{des} + \all{-es} au nom \\
\hline
pendant la guerre & \all{während des Krieges} & masc. : \all{des} + \all{-es} \\
\hline
en l'espace d'un an & \all{innerhalb eines Jahres} & neutre : \all{eines} + \all{-s} \\
\hline
au lieu d'un crédit & \all{statt eines Kredits} & masc. : \all{eines} + \all{-s} \\
\hline
en dehors des villes & \all{außerhalb der Städte} & pluriel : \all{der}, nom inchangé \\
\hline
\end{tabularx}
\end{center}

\begin{attention}[Pièges fréquents des francophones]
\begin{enumerate}
\item \textbf{Oublier le \all{-s}/\all{-es} du génitif} aux noms masculins et neutres : \all{des Vertrags}, \all{des Landes}, \all{des Wachstums}. Écrire \all{trotz der Preis} est une double faute (mauvais article + nom sans terminaison) ; la bonne forme est \all{trotz des Preises}.
\item \textbf{Le datif à l'oral} : dans la langue parlée, on entend parfois \all{wegen dem Wetter} ou \all{trotz dem Regen}. C'est toléré à l'oral familier, mais c'est une \textbf{faute à l'écrit}, surtout au Bac : écrivez toujours \all{wegen des Wetters}.
\item \textbf{Confondre les cas} : après \all{wegen/trotz/während}, jamais d'accusatif. \all{während den Krieg} est faux ; dites \all{während des Krieges}.
\item \textbf{Faux ami} : \all{statt} signifie « au lieu de », pas « à la place » au sens de « ville » ; et \all{anlässlich} exige toujours le génitif : \all{anlässlich des Festes}.
\end{enumerate}
\end{attention}

\courssub{Méthode : comment choisir la bonne préposition}

\begin{methode}[Choisir la préposition selon le sens]
\begin{enumerate}
\item \textbf{Identifier la relation logique} entre les deux idées de la phrase.
\item \textbf{Cause} (pourquoi ?) → \all{wegen} (neutre, courant) ou \all{aufgrund} (plus soutenu, typique de la presse économique).
\item \textbf{Concession} (opposition, « et pourtant ») → \all{trotz}.
\item \textbf{Temps} (durée, simultanéité) → \all{während} ; durée limitée → \all{innerhalb}.
\item \textbf{Lieu} (espace) → \all{innerhalb} (dedans) / \all{außerhalb} (dehors).
\item \textbf{Vérifier le génitif} : article \all{des/der}, terminaison \all{-(e)s} pour les masculins et neutres, adjectif en \all{-en} : \all{trotz des hohen Preises}.
\end{enumerate}
\end{methode}

\courssub{Entraînement}

\begin{exoresolu}[Mise au génitif : exemple corrigé]
\textbf{Énoncé.} Complétez avec la forme correcte du génitif : \all{Trotz \_\_\_\_\_\_\_\_\_\_ (das hohe Risiko) kauft der Anleger die Aktie.}
\tcblower
\textbf{Solution détaillée.} Étape 1 : la préposition \all{trotz} exige le génitif. Étape 2 : \all{das Risiko} est un neutre au singulier → article \all{des}, nom avec \all{-s} : \all{des Risikos}. Étape 3 : l'adjectif \all{hoch} après article défini au génitif prend la désinence \all{-en} (déclinaison faible) : \all{hohen}. Réponse : \all{Trotz des hohen Risikos kauft der Anleger die Aktie.} « Malgré le risque élevé, l'investisseur achète l'action. »
\end{exoresolu}

\begin{exercice}[À vous : le génitif dans l'économie]
Complétez les phrases suivantes avec la forme correcte du génitif.

\begin{enumerate}
\item \all{Wegen \_\_\_\_\_\_\_\_\_\_ (das Bevölkerungswachstum) importiert das Land mehr Reis.}
\item \all{Trotz \_\_\_\_\_\_\_\_\_\_ (die große Entfernung) kommen viele Waren aus Europa nach Douala.}
\item \all{Während \_\_\_\_\_\_\_\_\_\_ (das letzte Jahr) sind die Preise gestiegen.}
\item \all{Aufgrund \_\_\_\_\_\_\_\_\_\_ (die günstigen Preise) kaufen die Händler viel Kakao.}
\item \all{Innerhalb \_\_\_\_\_\_\_\_\_\_ (eine Woche) hat der Kurs der Aktie stark reagiert.}
\item \all{Statt \_\_\_\_\_\_\_\_\_\_ (ein Kredit) erhält das Land eine Entwicklungshilfe.}
\end{enumerate}
\end{exercice}

\begin{corrige}[Exercice : le génitif dans l'économie]
\begin{enumerate}
\item \all{Wegen des Bevölkerungswachstums importiert das Land mehr Reis.} (neutre : \all{des} + \all{-s})
\item \all{Trotz der großen Entfernung kommen viele Waren aus Europa nach Douala.} (féminin : \all{der}, nom inchangé)
\item \all{Während des letzten Jahres sind die Preise gestiegen.} (neutre : \all{des Jahres})
\item \all{Aufgrund der günstigen Preise kaufen die Händler viel Kakao.} (pluriel : \all{der})
\item \all{Innerhalb einer Woche hat der Kurs der Aktie stark reagiert.} (féminin indéfini : \all{einer})
\item \all{Statt eines Kredits erhält das Land eine Entwicklungshilfe.} (masculin indéfini : \all{eines} + \all{-s})
\end{enumerate}
\end{corrige}

\begin{aretenir}[L'essentiel de la section]
Les prépositions \all{wegen}, \all{trotz}, \all{während}, \all{innerhalb}, \all{außerhalb}, \all{aufgrund}, \all{anlässlich} et \all{statt} sont toujours suivies du \cle{génitif} : \all{des} + nom en \all{-(e)s} pour les masculins et neutres, \all{der} pour les féminins et les pluriels. Elles expriment la cause, la concession, le temps ou le lieu et sont omniprésentes dans les textes économiques. À l'écrit du Bac, jamais de datif après ces prépositions.
\end{aretenir}

\coursec{Grammaire 2 : les participes en fonction d'épithète}

Après avoir étudié les prépositions avec le génitif (\all{wegen der steigenden Preise}, \all{trotz des Handels}), nous restons dans le même texte économique pour observer une autre construction typique de la langue journalistique et administrative allemande : le \cle{participe employé comme épithète}, c'est-à-dire placé devant un nom, comme un adjectif. Cette structure est un des marqueurs les plus visibles de l'allemand écrit soutenu, et sa maîtrise est indispensable pour l'épreuve de compréhension de l'écrit au baccalauréat.

\courssub{Observation : que remarque-t-on dans le texte ?}

Reprenons deux groupes nominaux extraits de notre texte support sur l'aide au développement et les échanges commerciaux :

\begin{exemplebox}[Deux groupes nominaux à observer]
\all{Die steigenden Kurse zeigen das wachsende Interesse der Anleger.}\\
« Les cours en hausse montrent l'intérêt croissant des investisseurs. »

\all{Die von beiden Ländern unterzeichneten Verträge sollen die Zusammenarbeit verstärken.}\\
« Les contrats signés par les deux pays doivent renforcer la coopération. »
\end{exemplebox}

Observons attentivement :

\begin{itemize}
\item \all{steigend} vient du verbe \all{steigen} (monter) : on a ajouté \textbf{-d} à l'infinitif, puis une désinence \textbf{-en} devant le nom pluriel \all{Kurse}.
\item \all{wachsend} vient de \all{wachsen} (croître) : même mécanisme, infinitif + \textbf{-d} + désinence \textbf{-e} devant le nom neutre \all{Interesse}.
\item \all{unterzeichnet} est le participe passé de \all{unterzeichnen} (signer), suivi lui aussi d'une désinence \textbf{-en} devant \all{Verträge}.
\item En français, on traduit naturellement par une \textbf{proposition relative} : « les cours \emph{qui montent} », « les contrats \emph{qui ont été signés} ». L'allemand, lui, condense tout cela en un seul mot placé devant le nom.
\end{itemize}

\begin{aretenir}[L'idée essentielle]
En allemand, un participe (présent ou passé) peut fonctionner comme un \textbf{adjectif épithète} : il se place \textbf{devant le nom} et prend la \textbf{désinence adjectivale} (genre, nombre, cas). C'est ce qu'on appelle la \emph{Partizipialattribut} ou épithète participiale.
\end{aretenir}

\courssub{Formation du Partizip I et du Partizip II épithètes}

\begin{propriete}[Formation des participes épithètes]
\textbf{1. Partizip I épithète} = infinitif + \textbf{-d} + désinence adjectivale.\\
\all{steigen} $\rightarrow$ \all{steigend} $\rightarrow$ \all{das steigende Interesse} « l'intérêt croissant » ; \all{die steigenden Kurse} « les cours en hausse ».

\textbf{2. Partizip II épithète} = participe passé + désinence adjectivale.\\
\all{unterzeichnen} $\rightarrow$ \all{unterzeichnet} $\rightarrow$ \all{das unterzeichnete Dokument} « le document signé » ; \all{die unterzeichneten Verträge} « les contrats signés ».
\end{propriete}

Rappel sur la formation du participe passé (Partizip II) :

\begin{center}
\begin{tabularx}{\linewidth}{|X|X|X|}
\hline
\rowcolor{popblueL} Type de verbe & Formation & Exemple épithète \\
\hline
Verbe faible & ge + radical + t & \all{die exportierten Waren} (les marchandises exportées) \\
\hline
Verbe fort & ge + (radical modifié) + en & \all{die gestiegenen Preise} (les prix qui ont augmenté) \\
\hline
Verbe en -ieren & radical + t, \textbf{sans ge-} & \all{das investierte Geld} (l'argent investi) \\
\hline
Préfixe inséparable (be-, ver-, ent-, er-, unter-, über-...) & \textbf{sans ge-} & \all{der unterzeichnete Vertrag} (le contrat signé) \\
\hline
Préfixe séparable & ge entre préfixe et radical & \all{die eingeführten Produkte} (les produits importés) \\
\hline
\end{tabularx}
\end{center}

\begin{attention}[Pas de ge- avec -ieren et les préfixes inséparables]
Les verbes en \textbf{-ieren} et les verbes à \textbf{préfixe inséparable} (\all{be-, ver-, ent-, er-, unter-, über-}...) forment leur Partizip II \textbf{sans ge-} : \all{investiert} (et non « geinvestiert »), \all{unterzeichnet} (et non « geunterzeichnet »), \all{exportiert}, \all{organisiert}, \all{vereinbart}. Erreur très fréquente chez les francophones !
\end{attention}

\courssub{Sens : actif et en cours (Partizip I) ou passif et achevé (Partizip II)}

Le choix entre Partizip I et Partizip II n'est pas arbitraire : il obéit à une logique de sens qu'il faut bien comprendre.

\begin{definition}[Valeur des deux participes]
\textbf{Partizip I épithète} : l'action est \textbf{active} et \textbf{en cours} ; le nom \emph{fait} l'action.\\
\all{die steigenden Kurse} = les cours \emph{qui montent} (les cours montent, action en cours).

\textbf{Partizip II épithète} : l'action est \textbf{passive} ou \textbf{achevée} ; le nom \emph{subit} l'action.\\
\all{die unterzeichneten Verträge} = les contrats \emph{qui ont été signés} (on a signé les contrats, action achevée et passive).
\end{definition}

\begin{popfigure}
\begin{tikzpicture}[
  part/.style={rectangle, rounded corners, draw=popblue, fill=popblueL, align=center, minimum width=4.2cm, minimum height=1.6cm, font=\small},
  partb/.style={rectangle, rounded corners, draw=poporange, fill=poporangeL, align=center, minimum width=4.2cm, minimum height=1.6cm, font=\small},
  nom/.style={rectangle, rounded corners, draw=popgreen, fill=popgreenL, align=center, minimum width=3.6cm, minimum height=1.2cm, font=\small\bfseries},
  fleche/.style={-{Stealth[length=3mm]}, thick}
]
\node[part, align=center] (p1) at (0,2.2){\textbf{Partizip I}\\ actif, action en cours\\ \all{steigend}, \all{wachsend}};
\node[partb, align=center] (p2) at (0,-2.2){\textbf{Partizip II}\\ passif, action achevée\\ \all{unterzeichnet}, \all{exportiert}};
\node[nom, align=center] (n) at (8,0){le NOM\\ \all{die Kurse}\\ \all{die Verträge}};
\draw[fleche, popblue] (p1.east) -- node[above, font=\small, align=center]{\all{die steigenden Kurse}\\ = qui montent} (n.160);
\draw[fleche, poporange] (p2.east) -- node[below, font=\small, align=center]{\all{die unterzeichneten Verträge}\\ = qui ont été signés} (n.200);
\end{tikzpicture}
\legende{Les deux participes épithètes déterminent le nom : Partizip I (actif, en cours) et Partizip II (passif, achevé).}
\end{popfigure}

\courssub{L'accord : le participe épithète se décline comme un adjectif}

Point capital : devant le nom, le participe se comporte \textbf{exactement comme un adjectif} et suit les mêmes tableaux de désinences (après article défini, indéfini ou sans article). Le tableau ci-dessous rappelle la \textbf{déclinaison faible} (après article défini \all{der/die/das} ou pronom démonstratif \all{dieser/diese/dieses}) pour le Partizip I \all{steigend} :

\begin{center}
\begin{tabular}{|l|l|l|l|l|}
\hline
\rowcolor{popblueL} & \textbf{Masculin} & \textbf{Féminin} & \textbf{Neutre} & \textbf{Pluriel} \\
\hline
Nominatif & \all{der steigende Kurs} & \all{die steigende Nachfrage} & \all{das steigende Interesse} & \all{die steigenden Kurse} \\
\hline
Accusatif & \all{den steigenden Kurs} & \all{die steigende Nachfrage} & \all{das steigende Interesse} & \all{die steigenden Kurse} \\
\hline
Datif & \all{dem steigenden Kurs} & \all{der steigenden Nachfrage} & \all{dem steigenden Interesse} & \all{den steigenden Kursen} \\
\hline
Génitif & \all{des steigenden Kurses} & \all{der steigenden Nachfrage} & \all{des steigenden Interesses} & \all{der steigenden Kurse} \\
\hline
\end{tabular}
\end{center}

\begin{aretenir}[Rappel : les trois déclinaisons de l'adjectif/participe]
\begin{itemize}
\item \textbf{Déclinaison faible} (après \all{der, die, das, dieser, jener...}) : désinence \textbf{-e} (N sg. M/N/F), \textbf{-en} (partout ailleurs).
\item \textbf{Déclinaison mixte} (après \all{ein, kein, mein, dein...}) : comme faible, sauf N sg. M (\textbf{-er}), N sg. N (\textbf{-es}), A sg. N (\textbf{-es}).
\item \textbf{Déclinaison forte} (sans article) : l'adjectif/participe porte la marque de cas du nom (\textbf{-er, -e, -es, -en}).
\end{itemize}
\end{aretenir}

Avec article indéfini (déclinaison mixte) : \all{ein steigender Kurs} (un cours en hausse), \all{ein unterzeichneter Vertrag} (un contrat signé). Au pluriel après article défini : toujours \textbf{-en} : \all{die steigenden Kurse}, \all{die unterzeichneten Verträge}.

\begin{exemplebox}[Exemples variés, traduits]
\all{Die wachsende Wirtschaft Kameruns braucht ausländische Investitionen.}\\
« L'économie croissante du Cameroun a besoin d'investissements étrangers. »

\all{Der exportierte Kakao wird in Europa verarbeitet.}\\
« Le cacao exporté est transformé en Europe. »

\all{Wegen der fallenden Preise verkaufen viele Anleger ihre Aktien.}\\
« À cause des prix en baisse, beaucoup d'investisseurs vendent leurs actions. »

\all{Die im Hafen von Douala geladenen Waren kommen aus der ganzen Region.}\\
« Les marchandises chargées dans le port de Douala viennent de toute la région. »

\all{Das entwickelte Land hilft den ärmeren Ländern.}\\
« Le pays développé aide les pays plus pauvres. »
\end{exemplebox}

Remarquez dans le quatrième exemple une \textbf{épithète participiale longue} : tout le groupe \all{im Hafen von Douala} vient s'insérer entre l'article et le participe. En français, il faut développer : « les marchandises \emph{qui sont chargées dans le port de Douala} ».

\courssub{Comparaison français / allemand}

\begin{center}
\begin{tabularx}{\linewidth}{|X|X|X|}
\hline
\rowcolor{popblueL} Français & Deutsch & Remarque \\
\hline
les cours qui montent & \all{die steigenden Kurse} & le français préfère la relative \\
\hline
l'intérêt qui grandit & \all{das wachsende Interesse} & Partizip I = actif, en cours \\
\hline
les contrats qui ont été signés & \all{die unterzeichneten Verträge} & Partizip II = passif, achevé \\
\hline
les produits importés & \all{die importierten Produkte} & le français a aussi l'épithète courte \\
\hline
les contrats signés par les deux pays & \all{die von beiden Ländern unterzeichneten Verträge} & épithète longue : typique de la presse allemande \\
\hline
\end{tabularx}
\end{center}

\begin{attention}[Épithète ou verbe du Perfekt ? Ne pas confondre !]
Comparez bien :\\
\all{Die Kurse sind gestiegen.} = « Les cours ont monté. » Ici \all{gestiegen} est le \textbf{participe du Perfekt}, placé en fin de phrase, \textbf{sans désinence}.\\
\all{Die gestiegenen Kurse...} = « Les cours qui ont monté... » Ici \all{gestiegenen} est une \textbf{épithète}, placée devant le nom, \textbf{avec désinence} (-en).\\
Même participe de base, deux fonctions différentes : la place dans la phrase et la présence de la désinence permettent de les distinguer.
\end{attention}

\begin{savaistu}[L'épithète participiale, marque de la presse économique]
L'épithète participiale longue, comme \all{die von beiden Ländern unterzeichneten Verträge}, est typique de la presse économique et politique allemande (\emph{Frankfurter Allgemeine Zeitung}, \emph{Der Spiegel}, \emph{Die Zeit}). Elle permet de dire beaucoup en peu de mots. Savoir la découper — repérer l'article, le participe final et le nom — est un atout décisif pour l'épreuve de compréhension de l'écrit au baccalauréat : les textes économiques sur la bourse, le commerce ou l'aide au développement en regorgent.
\end{savaistu}

\courssub{Texte d'application : la bourse et le commerce vus par la presse}

\begin{textebox}[Börse und Handel in Kamerun]
Die steigenden Kurse an den internationalen Märkten freuen die kamerunischen Kakaobauern. Der exportierte Kakao bringt dem Land wichtige Devisen. Die von der Regierung unterzeichneten Verträge sollen den Handel mit den europäischen Partnern verstärken.

An der Börse beobachten die Anleger die wachsende Nachfrage nach afrikanischen Rohstoffen. Die investierten Summen sind in den letzten Jahren deutlich gestiegen. Viele im Hafen von Douala gegründete junge Unternehmen profitieren von den fallenden Transportkosten.

Aber es gibt auch Probleme: die importierten Produkte sind oft billiger als die lokal produzierten Waren. Die von Experten diskutierte Frage lautet: Wie kann die Entwicklungshilfe die einheimische Wirtschaft wirklich stärken? Die geplanten Projekte müssen den Menschen vor Ort helfen, nicht nur den ausländischen Investoren.

Ein bekanntes Sprichwort sagt: „Handel verbindet die Völker.“ Die zwischen Afrika und Europa wachsenden Beziehungen zeigen, dass dieser Satz auch heute noch stimmt.
\end{textebox}

\textbf{Glossaire :} \all{der Kurs, -e} (le cours, la cotation) ; \all{der Kakaobauer, -n} (le planteur de cacao) ; \all{die Devise, -n} (la devise) ; \all{die Nachfrage} (la demande) ; \all{der Rohstoff, -e} (la matière première) ; \all{die Summe, -n} (la somme) ; \all{die Ware, -n} (la marchandise) ; \all{die Entwicklungshilfe} (l'aide au développement) ; \all{der Investor, -en} (l'investisseur) ; \all{das Sprichwort, -er} (le proverbe) ; \all{die Beziehung, -en} (la relation).

\textbf{Questions de compréhension :}
\begin{enumerate}
\item Pourquoi les planteurs de cacao camerounais sont-ils contents ?
\item Que doivent faire les contrats signés par le gouvernement ?
\item Quel problème posent les produits importés, selon le texte ?
\item Relevez dans le texte trois participes épithètes et précisez s'il s'agit du Partizip I ou du Partizip II.
\end{enumerate}

\courssub{Exercices d'application}

\begin{exoresolu}[Version : traduisez en français]
Énoncé : traduisez les trois phrases suivantes en français correct.\\
1. \all{Die steigenden Preise machen vielen Familien Sorgen.}\\
2. \all{Die von beiden Ländern unterzeichneten Verträge treten nächsten Monat in Kraft.}\\
3. \all{Das in Kamerun investierte Geld schafft neue Arbeitsplätze.}
\tcblower
Solution détaillée :\\
1. « Les prix qui augmentent (en hausse) inquiètent beaucoup de familles. » — \all{steigenden} : Partizip I épithète, nominatif pluriel (sujet) ; action active et en cours. \all{jmdm. Sorgen machen} = inquiéter quelqu'un (datif).\\
2. « Les contrats signés par les deux pays entrent en vigueur le mois prochain. » — épithète longue : \all{von beiden Ländern} se loge entre l'article et le participe ; \all{in Kraft treten} = entrer en vigueur.\\
3. « L'argent investi au Cameroun crée de nouveaux emplois. » — \all{investiert} : Partizip II sans ge- (verbe en -ieren) ; \all{der Arbeitsplatz, -\"e} = l'emploi, le poste de travail.
\end{exoresolu}

\begin{exercice}[Thème : traduisez en allemand]
1. « les prix augmentant rapidement »\\
2. « les contrats signés hier »\\
3. « Les marchandises exportées passent par le port de Douala. »\\
4. « La demande croissante fait monter les cours. »
\end{exercice}

\begin{corrige}[Exercice de thème]
1. \all{die schnell steigenden Preise} — Partizip I de \all{steigen} + désinence -en (pluriel) ; l'adverbe \all{schnell} se place devant le participe.\\
2. \all{die gestern unterzeichneten Verträge} — Partizip II de \all{unterzeichnen}, sans ge- (préfixe inséparable) ; \all{gestern} devant le participe.\\
3. \all{Die exportierten Waren gehen über den Hafen von Douala.} — \all{exportiert} sans ge- (verbe en -ieren).\\
4. \all{Die wachsende Nachfrage lässt die Kurse steigen.} — \all{wachsende} : Partizip I de \all{wachsen}, féminin singulier après article défini (déclinaison faible : -e).
\end{corrige}

\begin{attention}[Les trois pièges à éviter]
\begin{enumerate}
\item \textbf{Oublier la désinence} : on écrit \all{die steigenden Kurse}, jamais « die steigend Kurse ». Le participe épithète s'accorde en genre, nombre et cas, comme tout adjectif.
\item \textbf{Ajouter un ge- fautif} : \all{investiert}, \all{unterzeichnet}, \all{exportiert} — jamais « geinvestiert » ni « geunterzeichnet ».
\item \textbf{Traduire mot à mot le français} : « les cours qui montent » ne se traduit pas par une relative lourde en allemand soutenu ; préférez l'épithète : \all{die steigenden Kurse}. À l'inverse, en version, développez l'épithète allemande en relative française naturelle.
\end{enumerate}
\end{attention}

\begin{aretenir}[Bilan de la section]
\begin{itemize}
\item Partizip I épithète = infinitif + \textbf{-d} + désinence : action \textbf{active, en cours} (\all{die steigenden Kurse}).
\item Partizip II épithète = participe passé + désinence : action \textbf{passive, achevée} (\all{die unterzeichneten Verträge}).
\item Verbes en \textbf{-ieren} et préfixes \textbf{inséparables} : Partizip II \textbf{sans ge-}.
\item Le participe épithète se \textbf{décline comme un adjectif} devant le nom (faible, mixte, forte selon le déterminant).
\item Le français préfère la relative (« qui monte », « qui a été signé ») ; l'allemand journalistique préfère l'épithète participiale, parfois longue.
\end{itemize}
\end{aretenir}

\coursec{Conjugaison : révision des temps composés}

Après avoir étudié l'emploi des participes comme épithètes (\emph{steigende Kurse}, \emph{unterzeichnete Verträge}), nous passons maintenant à leur fonction principale : la formation des \textbf{temps composés}. En allemand, comme en français, ces temps permettent de situer une action par rapport à un autre moment (passé achevé, antériorité, achèvement futur). Le texte support de cette leçon nous offre des occurrences parfaites pour réviser le \textbf{Perfekt}, le \textbf{Plusquamperfekt} et découvrir le \textbf{Futur II}, indispensable pour l'examen du Baccalauréat.

\courssub{Repérage dans le texte support}

Relisons deux phrases extraites de notre texte sur l'aide au développement et la bourse :

\begin{textebox}[Extraits du texte support]
„Deutschland hat seine Entwicklungshilfe \textbf{erhöht}.“
„Nachdem die Kurse \textbf{gestiegen} \textbf{waren}, investierten die Anleger.“
\end{textebox}

\noindent
\textbf{Observation :}
\begin{itemize}
    \item \all{hat ... erhöht} : auxiliaire \cle{haben} au présent + \cle{Partizip II} (ge-h-oht) en position finale $\rightarrow$ \textbf{Perfekt}.
    \item \all{gestiegen ... waren} : auxiliaire \cle{sein} à l'imparfait (Präteritum) + \cle{Partizip II} devant l'auxiliaire (ordre de la subordonnée \all{nachdem}) $\rightarrow$ \textbf{Plusquamperfekt}.
\end{itemize}
Ces deux temps expriment le passé, mais avec des valeurs aspectuelles différentes : le Perfekt pour le passé \emph{achevé} (souvent à valeur de passé simple en narration ou de passé composé à l'oral), le Plusquamperfekt pour l'\emph{antériorité} (un passé antérieur à un autre passé).

\courssub{Le Perfekt : haben / sein + Partizip II}

Le Perfekt est le temps du passé le plus fréquent à l'oral et dans la presse écrite pour rendre compte d'un fait accompli.

\begin{propriete}[Formation du Perfekt]
\textbf{Forme :} \textbf{Présent de l'auxiliaire (haben / sein) + Partizip II en fin de phrase (phrase principale) ou devant l'auxiliaire (subordonnée).}

\textbf{Choix de l'auxiliaire :}
\begin{enumerate}
    \item \textbf{Haben} : verbes transitifs (avec accusatif), verbes intransitifs sans changement d'état/lieu, verbes réfléchis, verbes modaux.
    \item \textbf{Sein} : verbes intransitifs indiquant un \textbf{déplacement} (gehen, fahren, fliegen), un \textbf{changement d'état} (sterben, einschlafen, \cle{steigen}, \cle{sinken}, \cle{fallen}, werden, bleiben) ou l'équivalence (sein, werden, bleiben).
\end{enumerate}

\textbf{Partizip II (régulier / faible) :} \cle{ge-} + radical + \cle{-t} (ex. \all{exportieren} $\rightarrow$ \all{exportiert}, \all{investieren} $\rightarrow$ \all{investiert}, \all{unterzeichnen} $\rightarrow$ \all{unterzeichnet}).
\textbf{Partizip II (irrégulier / fort) :} \cle{ge-} + radical modifié + \cle{-en} (ex. \all{steigen} $\rightarrow$ \all{gestiegen}, \all{sinken} $\rightarrow$ \all{gesunken}, \all{helfen} $\rightarrow$ \all{geholfen}).
\end{propriete}

\begin{exemplebox}[Exemples au Perfekt — Contexte économique]
\all{Deutschland \textbf{hat} die Exporte \textbf{erhöht}.} — L'Allemagne a augmenté ses exportations. (transitif $\rightarrow$ haben)
\all{Die Firma \textbf{hat} einen Vertrag \textbf{unterzeichnet}.} — L'entreprise a signé un contrat.
\all{Die Aktienkurse \textbf{sind} stark \textbf{gestiegen}.} — Les cours des actions ont fortement monté. (changement d'état $\rightarrow$ sein)
\all{Der Euro \textbf{ist} gegenüber dem Dollar \textbf{gesunken}.} — L'euro a baissé par rapport au dollar.
\all{Wir \textbf{haben} in neue Projekte \textbf{investiert}.} — Nous avons investi dans de nouveaux projets.
\all{Die Bundesregierung \textbf{hat} den Entwicklungsländern \textbf{geholfen}.} — Le gouvernement fédéral a aidé les pays en développement. (datif $\rightarrow$ haben)
\end{exemplebox}

\begin{attention}[Piège majeur : le choix de l'auxiliaire avec \emph{steigen}, \emph{sinken}, \emph{fallen}]
Les francophones ont tendance à utiliser \textbf{haben} par analogie avec « \emph{avoir} monté / baissé ».
\alert{En allemand, ces verbes marquent un changement d'état $\rightarrow$ ils s'emploient OBLIGATOIREMENT avec \textbf{sein}.}
\begin{center}
\begin{tabular}{|l|l|l|}\hline
\textbf{Verbe} & \textbf{Correct (sein)} & \textbf{Incorrect (haben)} \\ \hline
steigen (monter) & Der Preis \textbf{ist} gestiegen. & *Der Preis \textbf{hat} gestiegen. \\ \hline
sinken (baisser) & Die Temperatur \textbf{ist} gesunken. & *Die Temperatur \textbf{hat} gesunken. \\ \hline
fallen (tomber) & Die Aktie \textbf{ist} gefallen. & *Die Aktie \textbf{hat} gefallen. \\ \hline
\end{tabular}
\end{center}
\emph{Astuce mnémotechnique :} « \textbf{Sein} für den \textbf{Wandel} » (Sein pour le changement).
\end{attention}

\courssub{Le Plusquamperfekt : hatte / war + Partizip II}

Le Plusquamperfekt (Vorvergangenheit) exprime l'\textbf{antériorité} : une action passée \emph{avant} une autre action passée (souvent introduite par \all{nachdem}, \all{bevor}, \all{als}).

\begin{propriete}[Formation et emploi du Plusquamperfekt]
\textbf{Forme :} \textbf{Präteritum de l'auxiliaire (hatte / war) + Partizip II.}
Même choix d'auxiliaire qu'au Perfekt (haben / sein).
\textbf{Ordre des mots :} En phrase principale : auxiliaire (2e position) ... Partizip II (finale). En subordonnée (\all{nachdem}, \all{bevor}) : Partizip II ... auxiliaire (finale).

\textbf{Emploi typique :}
\begin{itemize}
    \item Subordonnée temporelle avec \all{nachdem} (après que) + Plusquamperfekt $\rightarrow$ principale au Präteritum/Perfekt.
    \item Subordonnée avec \all{bevor} / \all{ehe} (avant que) + Plusquamperfekt (ou Präteritum).
\end{itemize}
\end{propriete}

\begin{exemplebox}[Exemples au Plusquamperfekt]
\all{Nachdem die Kurse \textbf{gestiegen} \textbf{waren}, verkauften die Anleger.} — Après que les cours \textbf{eurent monté} / \textbf{avaient monté}, les investisseurs vendirent.
\all{Die Regierung \textbf{hatte} das Abkommen \textbf{unterzeichnet}, bevor der Gipfel begann.} — Le gouvernement \textbf{avait signé} l'accord avant que le sommet ne commence.
\all{Wir \textbf{hatten} bereits \textbf{investiert}, als die Krise ausbrach.} — Nous \textbf{avions déjà investi} quand la crise a éclaté.
\all{Der Export \textbf{war} vor 2020 stark \textbf{gesunken}.} — L'exportation \textbf{avait fortement baissé} avant 2020.
\end{exemplebox}

\courssub{Le Futur II (Futur antérieur) — Complément indispensable pour le Bac}

Le Futur II (Vollendete Zukunft) est peu utilisé à l'oral mais \textbf{exigé à l'écrit (Bac)} pour exprimer une hypothèse sur le passé ou un achèvement certain dans le futur.

\begin{propriete}[Formation et emploi du Futur II]
\textbf{Forme :} \textbf{werden (Présent) + Partizip II + haben / sein (Infinitif) en fin de phrase.}
L'auxiliaire final (haben/sein) suit la même logique que pour le Perfekt.

\textbf{Deux emplois principaux :}
\begin{enumerate}
    \item \textbf{Supposition / Hypothèse sur le passé} (souvent avec \all{wohl}, \all{schon}, \all{bereits}) : « Il a probablement... »
    \item \textbf{Achèvement futur} (avec indication de temps \all{bis} + date) : « Il aura... d'ici là. »
\end{enumerate}
\end{propriete}

\begin{exemplebox}[Exemples au Futur II]
\all{Bis 2030 \textbf{wird} Deutschland die Hilfe \textbf{erhöht haben}.} — D'ici 2030, l'Allemagne \textbf{aura augmenté} son aide.
\all{Die Anleger \textbf{werden} ihre Aktien \textbf{verkauft haben}.} — Les investisseurs \textbf{auront vendu} leurs actions. (Achèvement futur)
\all{Der Kurs \textbf{wird} wohl \textbf{gestiegen sein}.} — Le cours \textbf{a probablement monté}. (Hypothèse sur le passé $\rightarrow$ sein)
\all{Die Verträge \textbf{werden} bis morgen \textbf{unterzeichnet sein}.} — Les contrats \textbf{auront été signés} d'ici demain. (Passif futur, Partizip II + sein)
\end{exemplebox}

\courssub{Tableau récapitulatif des trois temps composés}

\begin{center}
\begin{tabularx}{\linewidth}{|X|X|X|X|}\hline
\rowcolor{popblueL}
\textbf{Temps} & \textbf{Structure} & \textbf{Emploi principal} & \textbf{Exemple (steigen / sinken / exportieren)} \\ \hline
\textbf{Perfekt} &
Présent (haben/sein) + Partizip II &
Passé achevé (oral / presse) &
\all{Die Kurse \textbf{sind} \textbf{gestiegen}.} (Les cours ont monté.) \\
& & & \all{Wir \textbf{haben} \textbf{exportiert}.} (Nous avons exporté.) \\ \hline
\textbf{Plusquamperfekt} &
Präteritum (hatte/war) + Partizip II &
Antériorité (après \emph{nachdem}, \emph{bevor}) &
\all{Nachdem die Kurse \textbf{gestiegen} \textbf{waren}...} (Après que les cours eurent monté...) \\
& & & \all{Die Firma \textbf{hatte} \textbf{investiert}.} (L'entreprise avait investi.) \\ \hline
\textbf{Futur II} &
werden (Prés.) + Partizip II + haben/sein (Inf.) &
Hypothèse passé / Achèvement futur (\emph{bis}) &
\all{Die Kurse \textbf{werden} \textbf{gestiegen sein}.} (Les cours auront monté / ont dû monter.) \\
& & & \all{Bis 2025 \textbf{wird} die Firma \textbf{exportiert haben}.} (D'ici 2025, la firme aura exporté.) \\ \hline
\end{tabularx}
\end{center}

\courssub{Frise chronologique des temps composés}

La figure ci-dessous visualise l'axe temporel et la relation d'antériorité entre les temps.

\begin{popfigure}
\begin{tikzpicture}[
    node distance=0.8cm and 1.5cm,
    timebox/.style={rectangle, rounded corners, draw=popblue, thick, fill=popblueL, text width=3.2cm, align=center, minimum height=1.8cm, font=\small},
    arrow/.style={-Stealth, thick, popdark},
    labelnode/.style={font=\footnotesize, text=popmuted, align=center}
]
% Noeuds principaux
\node[timebox, align=center] (pqp){\textbf{Plusquamperfekt}\\\textit{(Antériorité passée)}\\\all{war gestiegen / hatte exportiert}};
\node[timebox, right=of pqp, align=center] (perf){\textbf{Perfekt}\\\textit{(Passé achevé)}\\\all{ist gestiegen / hat exportiert}};
\node[timebox, right=of perf, align=center] (pres){\textbf{Présent}\\\textit{(Repère « Maintenant »)}\\\all{steigt / exportiert}};
\node[timebox, right=of pres, align=center] (fut2){\textbf{Futur II}\\\textit{(Achèvement futur / Hypothèse)}\\\all{werden gestiegen sein / werden exportiert haben}};

% Fleches
\draw[arrow] (pqp.east) -- (perf.west) node[midway, above, labelnode] {antérieur à};
\draw[arrow] (perf.east) -- (pres.west) node[midway, above, labelnode] {achevé avant};
\draw[arrow] (pres.east) -- (fut2.west) node[midway, above, labelnode] {achèvement futur};

% Exemples contextuels sous chaque noeud
\node[labelnode, below=0.4cm of pqp] {\all{„Nachdem die Preise \textbf{gestiegen} \textbf{waren}...“}};
\node[labelnode, below=0.4cm of perf] {\all{„Die Preise \textbf{sind} \textbf{gestiegen}.“}};
\node[labelnode, below=0.4cm of pres] {\all{„Die Preise \textbf{steigen} heute.“}};
\node[labelnode, below=0.4cm of fut2] {\all{„Bis 2030 \textbf{werden} sie \textbf{gestiegen sein}.“}};

% Ligne de temps decoratif
\draw[popline, very thick] ($(pqp.west)+(-0.5,0)$) -- ($(fut2.east)+(0.5,0)$);
\foreach \x/\txt in {pqp/« Passé antérieur », perf/« Passé », pres/« Présent », fut2/« Futur »} {
    \draw[popline, thick] (\x.south) -- ++(0,-0.3) node[below=2pt, labelnode] {\txt};
}
\end{tikzpicture}
\legende{Repérage temporel : Plusquamperfekt $\leftarrow$ Perfekt $\leftarrow$ Présent $\rightarrow$ Futur II. Les verbes de changement d'état (steigen, sinken) prennent \textbf{sein} à tous les temps composés.}
\end{popfigure}

\courssub{Exercice d'application : conjugaison thématique}

Conjuguez les six verbes thématiques ci-dessous aux trois temps composés étudiés (Perfekt, Plusquamperfekt, Futur II) pour la personne indiquée. Respectez le choix de l'auxiliaire (haben / sein).

\begin{exoresolu}[Conjugaison guidée : 6 verbes $\times$ 3 temps]
\textbf{Verbes :} \all{exportieren} (régulier, haben), \all{steigen} (fort, sein), \all{investieren} (régulier, haben), \all{unterzeichnen} (régulier, haben), \all{helfen} (fort, haben + datif), \all{sinken} (fort, sein).

\textbf{Sujet :} \all{die Firmen} (3\textsuperscript{e} pers. plur.) / \all{der Kurs} (3\textsuperscript{e} pers. sing.) selon le sens.

\tcblower
\textbf{Solution détaillée :}

\begin{center}
\begin{tabular}{|l|c|c|c|}\hline
\rowcolor{popblueL}
\textbf{Verbe (Sujet)} & \textbf{Perfekt} & \textbf{Plusquamperfekt} & \textbf{Futur II} \\ \hline
\all{exportieren} (die Firmen) &
\all{die Firmen \textbf{haben exportiert}} &
\all{die Firmen \textbf{hatten exportiert}} &
\all{die Firmen \textbf{werden exportiert haben}} \\ \hline
\all{steigen} (der Kurs) &
\all{der Kurs \textbf{ist gestiegen}} &
\all{der Kurs \textbf{war gestiegen}} &
\all{der Kurs \textbf{werden gestiegen sein}} \\ \hline
\all{investieren} (die Firmen) &
\all{die Firmen \textbf{haben investiert}} &
\all{die Firmen \textbf{hatten investiert}} &
\all{die Firmen \textbf{werden investiert haben}} \\ \hline
\all{unterzeichnen} (die Firmen) &
\all{die Firmen \textbf{haben unterzeichnet}} &
\all{die Firmen \textbf{hatten unterzeichnet}} &
\all{die Firmen \textbf{werden unterzeichnet haben}} \\ \hline
\all{helfen} (die Firmen / den Ländern) &
\all{die Firmen \textbf{haben geholfen}} &
\all{die Firmen \textbf{hatten geholfen}} &
\all{die Firmen \textbf{werden geholfen haben}} \\ \hline
\all{sinken} (der Kurs) &
\all{der Kurs \textbf{ist gesunken}} &
\all{der Kurs \textbf{war gesunken}} &
\all{der Kurs \textbf{werden gesunken sein}} \\ \hline
\end{tabular}
\end{center}

\textbf{Points de vigilance corrigés :}
\begin{itemize}
    \item \all{steigen / sinken} $\rightarrow$ auxiliaire \textbf{sein} à tous les temps (Partizip II : \all{gestiegen}, \all{gesunken}).
    \item \all{helfen} $\rightarrow$ verbe fort (voyelle \emph{e} $\rightarrow$ \emph{a} au Präteritum \all{half}, Partizip \all{geholfen}) mais auxiliaire \textbf{haben} (verbe transitif indirect / datif).
    \item Futur II : \all{werden} (conjugué) + \textbf{Partizip II} + \textbf{haben/sein} (infinitif final). Ordre strict : \all{werden ... Partizip II ... haben/sein}.
\end{itemize}
\end{exoresolu}

\courssub{Production écrite guidée : mini-texte économique}

Pour consolider ces temps, rédigez un court paragraphe (6--8 lignes) en réinvestissant le vocabulaire de la leçon (\cle{Entwicklungshilfe}, \cle{Handel}, \cle{Börse}, \cle{Aktie}, \cle{Export}, \cle{Import}) et \textbf{au moins une occurrence de chaque temps composé} (Perfekt, Plusquamperfekt, Futur II).

\begin{textebox}[Modèle de production — « Bilan et perspectives 2024--2030 »]
„Im Jahr 2024 \textbf{hat} Deutschland seine \textbf{Entwicklungshilfe} für Afrika \textbf{erhöht}. Besonders der \textbf{Handel} mit Kakao und Kaffee \textbf{ist} gewachsen. Nachdem die \textbf{Export}-Zahlen \textbf{gestiegen} \textbf{waren}, \textbf{haben} viele deutsche Firmen neue Fabriken in Kamerun \textbf{investiert}. Auch an der \textbf{Börse} \textbf{sind} die \textbf{Aktien}-Kurse dieser Unternehmen \textbf{gestiegen}. Experten prognostizieren, dass bis 2030 der \textbf{Import} afrikanischer Rohstoffe \textbf{weiter gestiegen sein wird} und die Partnerschaft \textbf{vertieft haben wird}.“
\end{textebox}

\begin{center}
\begin{tabularx}{\linewidth}{|X|X|}\hline
\rowcolor{popblueL}
\textbf{Allemand (Glossaire difficile)} & \textbf{Français} \\ \hline
\all{die Entwicklungshilfe, -n} & l'aide au développement \\ \hline
\all{der Handel (kein Plural / die Händel)} & le commerce, les échanges \\ \hline
\all{die Börse, -n} & la bourse \\ \hline
\all{die Aktie, -n} & l'action (bourse) \\ \hline
\all{der Export, -e} & l'exportation \\ \hline
\all{der Import, -e} & l'importation \\ \hline
\all{erhöhen} (reg., haben) & augmenter (qc) \\ \hline
\all{wachsen} (fort, sein: gewachsen) & croître, augmenter (intransitif) \\ \hline
\all{investieren} (reg., haben) & investir \\ \hline
\all{prognostizieren} (reg., haben) & prévoir, pronostiquer \\ \hline
\all{der Rohstoff, -e} & la matière première \\ \hline
\all{vertiefen} (reg., haben) & approfondir, renforcer \\ \hline
\end{tabularx}
\end{center}

\begin{exercice}[Votre tour : Production écrite]
Rédigez en allemand un paragraphe de 6 à 8 lignes sur le thème : \emph{„Die wirtschaftliche Zusammenarbeit zwischen Deutschland und Kamerun“}.
\textbf{Contraintes obligatoires :}
\begin{enumerate}
    \item Utilisez au moins 5 mots du glossaire ci-dessus.
    \item Insérez \textbf{au moins 1 Perfekt}, \textbf{1 Plusquamperfekt} (avec \all{nachdem} ou \all{bevor}) et \textbf{1 Futur II} (avec \all{bis 2030} ou hypothèse \all{wohl}).
    \item Soulignez les formes verbales composées dans votre copie.
\end{enumerate}
\end{exercice}

\begin{corrige}[Exercice Production écrite — Proposition de correction]
„In den letzten Jahren \textbf{hat} der \textbf{Handel} zwischen Deutschland und Kamerun \textbf{zugenommen}. Nachdem deutsche Unternehmen neue Verträge \textbf{unterzeichnet hatten}, \textbf{sind} die \textbf{Exporte} von Maschinen \textbf{gestiegen}. Auch die \textbf{Entwicklungshilfe} \textbf{hat} geholfen, die Infrastruktur zu verbessern. An der \textbf{Börse} \textbf{sind} die \textbf{Aktien} entsprechender Firmen \textbf{gestiegen}. Bis 2030 \textbf{wird} das Handelsvolumen \textbf{sicherlich verdoppelt haben}.“
\end{corrige}

\begin{methode}[Comment réviser efficacement les temps composés pour le Bac]
\begin{enumerate}
    \item \textbf{Mémorisez les 3 structures} (Tableau récapitulatif) : Auxiliaire conjugué + Partizip II (+ Infinitif auxiliaire pour Futur II).
    \item \textbf{Classez vos verbes} par auxiliaire : Liste \textbf{SEIN} (mouvement, changement d'état : \all{gehen, fahren, steigen, sinken, fallen, werden, bleiben, sein}) vs Liste \textbf{HABEN} (tout le reste).
    \item \textbf{Entraînez-vous à la transformation} : Phrase au Présent $\rightarrow$ Perfekt $\rightarrow$ Plusquamperfekt $\rightarrow$ Futur II.
    \item \textbf{Repérez les marqueurs temporels} : \all{gestern/2024} (Perfekt/Präteritum), \all{nachdem/bevor} (Plusquamperfekt), \all{bis 2030/wohl} (Futur II).
\end{enumerate}
\end{methode}

\coursec{Méthode du commentaire et production guidée}

Après avoir révisé les temps composés, nous disposons de tous les outils linguistiques de la leçon : le vocabulaire de l'économie, les prépositions avec le génitif, les participes épithètes et les temps. Il reste à les mettre au service de l'\cle{expression écrite et orale}, compétence décisive au baccalauréat. Cette dernière partie vous donne une méthode complète : commenter un texte, le résumer, donner son opinion et débattre.

\courssub{Le commentaire de texte au Bac : la méthode}

\begin{definition}[Der Kommentar]
Le \cle{commentaire} (allemand : \all{der Kommentar, -e} ou \all{die Stellungnahme, -n}, prise de position) est un texte argumentatif dans lequel vous présentez un document, analysez son contenu et donnez votre avis personnel de manière structurée. Il se compose de trois parties : \all{die Einleitung, -en} (introduction), \all{der Hauptteil, -e} (développement), \all{der Schluss, -sse} (conclusion).
\end{definition}

\begin{methode}[Rédiger un commentaire en trois étapes]
\begin{enumerate}
\item \textbf{\all{Die Einleitung} (l'introduction).} Présentez le texte : le thème (\all{das Thema}), la source (\all{die Quelle} : journal, discours, article) et la problématique (\all{die Leitfrage}). Formule utile : \all{Der Text behandelt das Thema … und wirft die Frage auf: …} (Le texte traite du thème … et pose la question : …).
\item \textbf{\all{Der Hauptteil} (le développement).} Développez 2 ou 3 axes. Chaque axe = un argument + un exemple cité du texte ou de votre expérience (par exemple le Cameroun). Reliez les arguments avec des connecteurs : \all{zuerst} (d'abord), \all{außerdem} (de plus), \all{schließlich} (enfin).
\item \textbf{\all{Der Schluss} (la conclusion).} Faites le bilan (\all{zusammenfassend kann man sagen} : pour résumer, on peut dire) et ouvrez sur une question plus large (\all{Es bleibt die Frage, ob …} : reste la question de savoir si …).
\end{enumerate}
\end{methode}

\begin{popfigure}
\begin{center}
\begin{tikzpicture}[
node distance=0.55cm and 1.2cm,
box/.style={draw=popblue, fill=popblueL, rounded corners=3pt, align=center, inner sep=5pt, font=\small},
mid/.style={draw=poporange, fill=poporangeL, rounded corners=3pt, align=center, inner sep=5pt, font=\small},
end/.style={draw=popgreen, fill=popgreenL, rounded corners=3pt, align=center, inner sep=5pt, font=\small},
arr/.style={-{Stealth[length=2.5mm]}, thick, popdark}]
\node[box, align=center] (ein){\textbf{Einleitung}\\ Thème + source + problématique};
\node[mid, below left=0.7cm and -0.3cm of ein, align=center] (a1){\textbf{Argument 1}\\ + Beispiel};
\node[mid, below right=0.7cm and -0.3cm of ein, align=center] (a2){\textbf{Argument 2}\\ + Beispiel};
\node[end, below=2.6cm of ein, align=center] (schl){\textbf{Schluss}\\ Bilan + ouverture};
\draw[arr] (ein) -- (a1);
\draw[arr] (ein) -- (a2);
\draw[arr] (a1.south) |- ($(schl.north)+(0,0.25)$) -| (a2.south);
\draw[arr] ($(schl.north)+(0,0.25)$) -- (schl.north);
\node[font=\small\itshape, popmuted, above=0.1cm of ein] {Der Kommentar : plan en trois parties};
\end{tikzpicture}
\end{center}
\legende{Organigramme du commentaire : Einleitung, Hauptteil (arguments + exemples), Schluss.}
\end{popfigure}

\courssub{Commentaire modèle commenté}

Voici un commentaire modèle répondant à la problématique : \all{Wie wichtig ist die wirtschaftliche Zusammenarbeit zwischen Deutschland und Kamerun?} (Quelle est l'importance de la coopération économique entre l'Allemagne et le Cameroun ?).

\begin{textebox}[Modèle de commentaire]
Der Text behandelt die wirtschaftliche Zusammenarbeit zwischen Deutschland und Kamerun und wirft die Frage auf, wie wichtig sie für beide Länder ist. Zuerst muss man sagen, dass Deutschland ein wichtiger Handelspartner Kameruns ist: Es importiert Kakao und exportiert Maschinen. Außerdem finanziert die deutsche Entwicklungshilfe Schulen, Krankenhäuser und Straßen, zum Beispiel in der Region um Yaoundé. Meiner Meinung nach ist diese Zusammenarbeit sehr wichtig, weil sie Arbeitsplätze schafft und die wachsende Wirtschaft Kameruns unterstützt. Schließlich profitiert auch Deutschland, denn es findet neue Märkte für seine Produkte. Zusammenfassend kann man sagen, dass beide Länder von dieser Partnerschaft profitieren. Es bleibt die Frage, ob die Hilfe in Zukunft erhöht werden sollte.
\end{textebox}

\textbf{Traduction.} Le texte traite de la coopération économique entre l'Allemagne et le Cameroun et pose la question de son importance pour les deux pays. D'abord, il faut dire que l'Allemagne est un partenaire commercial important du Cameroun : elle importe du cacao et exporte des machines. De plus, l'aide allemande au développement finance des écoles, des hôpitaux et des routes, par exemple dans la région de Yaoundé. À mon avis, cette coopération est très importante parce qu'elle crée des emplois et soutient l'économie croissante du Cameroun. Enfin, l'Allemagne en profite aussi, car elle trouve de nouveaux marchés pour ses produits. Pour résumer, on peut dire que les deux pays profitent de ce partenariat. Reste la question de savoir si l'aide devrait être augmentée à l'avenir.

\textbf{Repérage en classe.} Soulignez les connecteurs : \all{zuerst} (argument 1), \all{außerdem} (argument 2), \all{meiner Meinung nach} (opinion personnelle), \all{schließlich} (dernier argument), \all{zusammenfassend kann man sagen} (bilan), \all{es bleibt die Frage, ob} (ouverture). Remarquez aussi le participe épithète \all{die wachsende Wirtschaft} (l'économie croissante) et le génitif \all{ein Handelspartner Kameruns} (un partenaire du Cameroun) : réutiliser la grammaire de la leçon rapporte des points !

\begin{attention}[Place du verbe après « meiner Meinung nach »]
\all{Meiner Meinung nach} occupe la première place de la phrase : le verbe conjugué vient donc \textbf{en deuxième position}, avant le sujet (inversion). On écrit \all{Meiner Meinung nach ist die Hilfe wichtig.} et jamais \emph{Meiner Meinung nach die Hilfe ist wichtig.} Même règle avec \all{zuerst}, \all{außerdem}, \all{schließlich} en tête de phrase : \all{Außerdem finanziert die Entwicklungshilfe Schulen.}
\end{attention}

\begin{exemplebox}[Connecteurs et expressions d'opinion]
\all{Meiner Meinung nach ist die Entwicklungshilfe wichtig, weil sie Schulen und Krankenhäuser finanziert.} « À mon avis, l'aide au développement est importante parce qu'elle finance des écoles et des hôpitaux. »\par
\all{Ich bin der Meinung, dass der Handel beiden Ländern hilft.} « Je suis d'avis que le commerce aide les deux pays. »\par
\all{Ein Vorteil ist, dass die Zusammenarbeit Arbeitsplätze schafft.} « Un avantage est que la coopération crée des emplois. »
\end{exemplebox}

\begin{center}
\begin{tabularx}{\linewidth}{|X|X|X|}
\hline
\rowcolor{popblueL} Français & Deutsch & Fonction \\
\hline
d'abord & zuerst / zunächst & argument 1 \\
\hline
de plus, en outre & außerdem / darüber hinaus & ajout \\
\hline
enfin, finalement & schließlich / zum Schluss & dernier argument \\
\hline
à mon avis & meiner Meinung nach & opinion \\
\hline
parce que / car & weil / denn & justification \\
\hline
pour résumer & zusammenfassend & bilan \\
\hline
reste la question de savoir si & es bleibt die Frage, ob & ouverture \\
\hline
\end{tabularx}
\end{center}

\courssub{Le résumé guidé (die Zusammenfassung)}

\begin{definition}[Die Zusammenfassung]
Le \cle{résumé} (\all{die Zusammenfassung, -en}) restitue les idées principales d'un texte en 5 lignes maximum, avec vos propres mots, au présent, et \textbf{sans aucune opinion personnelle}.
\end{definition}

\begin{methode}[Rédiger un résumé en 4 étapes]
\begin{enumerate}
\item \textbf{Souligner} les idées principales du texte (une idée par paragraphe).
\item \textbf{Reformuler} avec des synonymes : \all{wichtig} $\rightarrow$ \all{bedeutend} ; \all{helfen} $\rightarrow$ \all{unterstützen} ; \all{geben} $\rightarrow$ \all{finanzieren}.
\item \textbf{Relier} les idées avec des connecteurs : \all{zuerst}, \all{außerdem}, \all{schließlich}.
\item \textbf{Vérifier} la longueur (5 lignes maximum) et l'absence totale d'opinion personnelle (pas de \all{ich finde}, pas de \all{meiner Meinung nach}).
\end{enumerate}
\end{methode}

\begin{attention}[Résumé ou commentaire : ne pas confondre !]
Dans un résumé (\all{Zusammenfassung}), on ne donne \textbf{jamais} son avis : on rapporte objectivement ce que dit le texte. L'opinion personnelle n'apparaît que dans le commentaire (\all{Kommentar} / \all{Stellungnahme}). Erreur fréquente des francophones : écrire \all{Ich finde den Text gut} dans un résumé — cela coûte des points. Autre piège : copier des phrases entières du texte au lieu de reformuler.
\end{attention}

\begin{exoresolu}[Résumer le commentaire modèle]
Énoncé : Résumez le commentaire modèle ci-dessus en 5 lignes maximum, au présent, sans opinion personnelle.
\tcblower
Solution détaillée. Étape 1 : idées principales = (a) l'Allemagne est un partenaire commercial du Cameroun (cacao contre machines) ; (b) l'aide au développement finance des infrastructures ; (c) les deux pays profitent de la coopération. Étapes 2 et 3 : reformulation avec connecteurs.\par
\all{Der Text beschreibt die wirtschaftliche Zusammenarbeit zwischen Deutschland und Kamerun. Zuerst zeigt er, dass Deutschland Kakao importiert und Maschinen exportiert. Außerdem unterstützt die Entwicklungshilfe den Bau von Schulen und Krankenhäusern. Schließlich betont der Text, dass beide Länder von dieser Partnerschaft profitieren.}\par
« Le texte décrit la coopération économique entre l'Allemagne et le Cameroun. D'abord, il montre que l'Allemagne importe du cacao et exporte des machines. De plus, l'aide au développement soutient la construction d'écoles et d'hôpitaux. Enfin, le texte souligne que les deux pays profitent de ce partenariat. »\par
Étape 4 : 4 lignes, présent partout, aucun \all{ich} — le résumé est correct.
\end{exoresolu}

\courssub{Production écrite guidée : le paragraphe d'opinion}

\begin{exercice}[Soll Deutschland seine Entwicklungshilfe für Kamerun erhöhen ?]
Rédigez un court paragraphe d'opinion (8 à 10 lignes) sur la question : « L'Allemagne doit-elle augmenter son aide au développement pour le Cameroun ? » Respectez le plan imposé :
\begin{enumerate}
\item \textbf{Thèse} : votre position (\all{Meiner Meinung nach soll Deutschland …}).
\item \textbf{Argument 1} + exemple (écoles, hôpitaux, routes).
\item \textbf{Argument 2} + exemple (emplois, commerce, cacao).
\item \textbf{Conclusion} : bilan en une phrase (\all{Zusammenfassend …}).
\end{enumerate}
Contraintes : utilisez au moins 3 mots du vocabulaire de la leçon, 1 préposition avec le génitif (\all{wegen der Armut}, \all{trotz der Fortschritte}), 1 participe épithète (\all{die wachsende Wirtschaft}) et 1 temps composé (\all{Deutschland hat viele Projekte finanziert}).
\end{exercice}

\begin{corrige}[Production écrite guidée]
Modèle de production :\par
\all{Meiner Meinung nach soll Deutschland seine Entwicklungshilfe für Kamerun erhöhen. Zuerst ist die Hilfe wegen der Armut in vielen Regionen sehr wichtig: Deutschland hat schon viele Schulen und Krankenhäuser finanziert. Außerdem schafft die Zusammenarbeit Arbeitsplätze in der wachsenden Wirtschaft Kameruns, zum Beispiel im Export von Kakao. Trotz der Fortschritte braucht das Land noch Unterstützung. Zusammenfassend kann man sagen, dass eine erhöhte Hilfe beiden Ländern Vorteile bringt.}\par
« À mon avis, l'Allemagne devrait augmenter son aide au développement pour le Cameroun. D'abord, l'aide est très importante à cause de la pauvreté dans de nombreuses régions : l'Allemagne a déjà financé beaucoup d'écoles et d'hôpitaux. De plus, la coopération crée des emplois dans l'économie croissante du Cameroun, par exemple dans l'exportation du cacao. Malgré les progrès, le pays a encore besoin de soutien. Pour résumer, on peut dire qu'une aide augmentée apporte des avantages aux deux pays. »\par
Vérification : 3 mots du thème (\all{die Entwicklungshilfe}, \all{die Arbeitsplätze}, \all{der Export}), 2 prépositions + génitif (\all{wegen der Armut}, \all{trotz der Fortschritte}), 1 participe épithète (\all{die wachsende Wirtschaft}), 1 temps composé (\all{hat … finanziert}).
\end{corrige}

\courssub{Production orale : le débat en binômes}

\begin{experience}[Débat : pour ou contre l'augmentation de l'aide ?]
Par binômes, préparez un mini-débat de 4 minutes. L'élève A défend l'augmentation de l'aide au développement (\all{dafür}), l'élève B la critique (\all{dagegen}). Déroulement :
\begin{enumerate}
\item Chacun prépare 2 arguments avec le vocabulaire de la leçon (\all{die Entwicklungshilfe}, \all{der Handel}, \all{der Export}, \all{die Armut}, \all{die Arbeitsplätze}).
\item A commence : \all{Ich bin dafür, weil …} ; B répond : \all{Ich bin anderer Meinung, denn …} (Je suis d'un autre avis, car …).
\item Échangez les rôles, puis concluez ensemble : \all{Wir sind der Meinung, dass …}.
\end{enumerate}
Le reste de la classe note les arguments et vote à la fin.
\end{experience}

\begin{exemplebox}[Expressions utiles pour le débat]
\all{Ich bin dafür / dagegen.} « Je suis pour / contre. »\par
\all{Ich bin anderer Meinung.} « Je ne suis pas de ton avis. »\par
\all{Das stimmt, aber …} « C'est vrai, mais … »\par
\all{Ein Gegenargument ist, dass die Hilfe abhängig machen kann.} « Un contre-argument est que l'aide peut rendre dépendant. »
\end{exemplebox}

\courssub{Grille d'auto-évaluation}

Avant de rendre votre production, vérifiez chaque critère et cochez :

\begin{center}
\begin{tabularx}{\linewidth}{|X|l|}
\hline
\rowcolor{popblueL} Critère & Oui / Non \\
\hline
J'ai utilisé au moins 3 mots du vocabulaire du thème (\all{die Entwicklungshilfe}, \all{der Handel}, \all{die Börse}, \all{der Export} …) & \\
\hline
J'ai employé au moins 1 préposition avec le génitif (\all{wegen}, \all{trotz}, \all{während} + génitif) & \\
\hline
J'ai employé au moins 1 participe en fonction d'épithète (\all{die wachsende Wirtschaft}) & \\
\hline
J'ai employé au moins 1 temps composé (\all{hat finanziert}, \all{ist gewachsen}) & \\
\hline
J'ai respecté le plan (thèse, 2 arguments, conclusion) et les connecteurs & \\
\hline
Dans un résumé : aucune opinion personnelle ; dans un commentaire : opinion claire & \\
\hline
\end{tabularx}
\end{center}

\begin{savaistu}[Comment l'expression écrite est-elle notée au Bac ?]
Au baccalauréat, l'expression écrite est évaluée sur trois critères : la \cle{pertinence du contenu} (vous répondez vraiment à la question posée), la \cle{correction grammaticale} (place du verbe, cas, temps) et la \cle{richesse du lexique}. Conséquence pratique : réutiliser le vocabulaire de la leçon (\all{die Entwicklungshilfe}, \all{die Zusammenarbeit}, \all{die Arbeitsplätze}) et les structures grammaticales étudiées (génitif, participe épithète, temps composés) rapporte directement des points. Un texte simple mais correct et riche en vocabulaire du thème vaut toujours mieux qu'un texte ambitieux plein de fautes !
\end{savaistu}

\begin{aretenir}[L'essentiel de la section]
\begin{itemize}
\item Commentaire (\all{Kommentar}) = Einleitung (thème + source + problématique) $\rightarrow$ Hauptteil (2 ou 3 arguments + exemples) $\rightarrow$ Schluss (bilan + ouverture). L'opinion personnelle y est attendue.
\item Résumé (\all{Zusammenfassung}) = 5 lignes maximum, présent, ses propres mots, jamais d'opinion.
\item Connecteurs clés : \all{zuerst}, \all{außerdem}, \all{schließlich}, \all{meiner Meinung nach}, \all{zusammenfassend}. Après un connecteur en tête de phrase, le verbe reste en deuxième position.
\item Pour maximiser la note : 3 mots du thème + 1 préposition avec génitif + 1 participe épithète + 1 temps composé.
\end{itemize}
\end{aretenir}

\newpage

\coursec{Activité d'intégration}

\courssub{Objectifs de l'activité}

À l'issue de cette activité d'intégration, l'élève sera capable de :

\begin{itemize}
\item réinvestir dans une production écrite authentique le vocabulaire économique de la leçon (\all{die Entwicklungshilfe}, \all{der Handel}, \all{die Börse}, \all{die Aktie}, \all{der Export}, \all{der Import} et leurs familles de mots) ;
\item employer correctement au moins deux prépositions avec le génitif (\all{wegen}, \all{trotz}, \all{während}, \all{dank}, \all{anlässlich}) ;
\item utiliser un participe I et un participe II en fonction d'épithète (\all{die wachsende Wirtschaft}, \all{das finanzierte Projekt}) ;
\item construire une phrase au Plusquamperfekt pour situer un événement dans le passé ;
\item rédiger en groupe un court article de presse scolaire (10 à 12 lignes) et le présenter à l'oral devant la classe ;
\item évaluer la production des autres groupes à l'aide de la grille d'évaluation vue en section 6.
\end{itemize}

\courssub{Situation}

Le journal bilingue de ton lycée, \all{„Die Brücke / Le Pont“}, prépare son prochain numéro consacré à la coopération entre le Cameroun et l'Allemagne. La rédaction recherche des articles annonçant un nouveau projet de partenariat économique. Ton groupe a été choisi pour écrire l'article de une. Pour vous inspirer, la rédaction vous transmet la dépêche suivante, publiée par un portail d'information camerounais :

\begin{textebox}[Dépêche d'actualité — modèle de départ]
\textbf{Neues Schulprojekt in Yaoundé}\par
Die deutsche Botschaft in Yaoundé hat gestern ein neues Kooperationsprojekt angekündigt : Trotz der wirtschaftlichen Schwierigkeiten wird Deutschland den Bau einer modernen Berufsschule im Stadtteil Nkolbisson finanzieren. Das von der deutschen Regierung geförderte Projekt soll im nächsten Jahr beginnen. „Wir investieren in die Zukunft der kamerunischen Jugend“, erklärte die Botschafterin. Gleichzeitig wächst der Handel zwischen beiden Ländern : Der Export von Kakao nach Deutschland steigt, und mehrere kamerunische Firmen interessieren sich für die Börse in Frankfurt. Während der letzten Jahre hatte Deutschland bereits zahlreiche Projekte in den Bereichen Bildung und Gesundheit unterstützt.
\end{textebox}

\textbf{Glossaire :} \all{die Botschaft, -en} (l'ambassade) ; \all{ankündigen} (annoncer) ; \all{die Berufsschule, -n} (le lycée professionnel) ; \all{fördern} (promouvoir, financer) ; \all{der Bereich, -e} (le domaine) ; \all{unterstützen} (soutenir).

\courssub{Consignes}

\textbf{Travail en groupes de 3 élèves. Durée : 2 heures (préparation + présentation orale).}

\begin{enumerate}
\item \textbf{Comprendre le modèle.} Lis la dépêche ci-dessus avec ton groupe. Souligne en rouge tous les mots du lexique économique de la leçon, en bleu les prépositions avec le génitif, en vert les participes épithètes. Identifie la phrase au Plusquamperfekt et justifie l'emploi de ce temps.
\item \textbf{Choisir un sujet.} Choisissez ensemble UN des trois projets suivants : (a) la construction d'une école à Douala financée par l'Allemagne ; (b) un nouveau contrat d'exportation de cacao camerounais vers Hambourg ; (c) l'arrivée d'investisseurs camerounais à la Bourse de Francfort.
\item \textbf{Préparer le plan.} Rédigez un plan en trois parties : titre accrocheur + annonce du projet (qui ? quoi ? où ? quand ?), détails et chiffres imaginaires mais vraisemblables, perspectives d'avenir pour les deux pays.
\item \textbf{Rédiger l'article} (10 à 12 lignes) en respectant les contraintes imposées : au moins \textbf{6 mots du lexique} de la leçon, \textbf{2 prépositions avec le génitif}, \textbf{1 participe I épithète}, \textbf{1 participe II épithète}, \textbf{1 phrase au Plusquamperfekt}.
\item \textbf{Relire et corriger.} Vérifiez en groupe : la place du verbe (2\textsuperscript{e} position), la majuscule à tous les noms, les déclinaisons du génitif (\all{des}, \all{der} + adjectif en \all{-en}), l'accord des participes épithètes.
\item \textbf{Présenter à l'oral.} Chaque groupe lit son article devant la classe. Les autres élèves remplissent la grille d'évaluation de la section 6 (contenu, lexique, grammaire, prononciation) et votent pour le meilleur article, qui sera publié dans \all{„Die Brücke / Le Pont“}.
\end{enumerate}

\courssub{Ressources à mobiliser}

\begin{aretenir}[Boîte à outils de la leçon]
\textbf{Lexique :} \all{die Entwicklungshilfe, -n} ; \all{der Handel} ; \all{der Export, -e} / \all{der Import, -e} ; \all{exportieren / importieren} ; \all{die Börse, -n} ; \all{die Aktie, -n} ; \all{der Investor, -en} ; \all{die Zusammenarbeit} ; \all{finanzieren} ; \all{investieren}.\par
\textbf{Prépositions + génitif :} \all{wegen des Regens}, \all{trotz der Schwierigkeiten}, \all{während des Projekts}, \all{dank der Hilfe}, \all{anlässlich des Besuchs}.\par
\textbf{Participes épithètes :} participe I + terminaison d'adjectif = action en cours (\all{die steigenden Preise} — les prix qui montent) ; participe II + terminaison = action achevée ou passive (\all{das finanzierte Projekt} — le projet financé).\par
\textbf{Plusquamperfekt :} \all{hatte / war} + participe II : \all{Deutschland hatte das Projekt bereits unterstützt.} (L'Allemagne avait déjà soutenu le projet.)
\end{aretenir}

\begin{attention}[Pièges à éviter]
Ne dis jamais \all{*wegen die Schwierigkeiten} : après \all{wegen}, \all{trotz}, \all{während}, le génitif est obligatoire au lycée : \all{wegen der Schwierigkeiten}. Attention aussi à l'ordre des mots : le verbe conjugué reste en 2\textsuperscript{e} position même après un long complément : \all{Trotz der Krise \textbf{steigt} der Export.} Enfin, un participe épithète se décline comme un adjectif : \all{das von Deutschland finanziert\textbf{e} Projekt}.
\end{attention}

\courssub{Proposition de production et corrigé commenté}

\begin{exoresolu}[Article modèle : „Kakao für Hamburg“]
Rédige un article de 10 à 12 lignes sur le sujet (b) : un nouveau contrat d'exportation de cacao.
\tcblower
\textbf{Production modèle :}\par
\all{\textbf{Kamerunischer Kakao erobert Deutschland}}\par
\all{Große Neuigkeiten für unsere Wirtschaft : Trotz der weltweiten Krise hat eine kamerunische Firma aus Douala gestern einen wichtigen Vertrag mit deutschen Händlern unterschrieben. Der Export von Kakao nach Hamburg wird sich in den nächsten Jahren verdoppeln. Das von beiden Regierungen geförderte Projekt schafft viele Arbeitsplätze für die junge, dynamische Bevölkerung unseres Landes. Während der Verhandlungen hatte die deutsche Seite bereits Investitionen in Höhe von zwei Millionen Euro versprochen. Dank dieser Zusammenarbeit werden auch kamerunische Aktien an der Börse in Frankfurt gehandelt. „Diese wachsende Partnerschaft ist eine Chance für beide Länder“, erklärte der Handelsminister. Die Entwicklungshilfe wird künftig durch fairen Handel ergänzt — ein starkes Signal für die Zukunft !}\par
« Le cacao camerounais conquiert l'Allemagne — Grandes nouvelles pour notre économie : malgré la crise mondiale, une entreprise camerounaise de Douala a signé hier un contrat important avec des commerçants allemands. L'exportation de cacao vers Hambourg va doubler dans les prochaines années. Le projet financé par les deux gouvernements crée de nombreux emplois pour la population jeune et dynamique de notre pays. Pendant les négociations, la partie allemande avait déjà promis des investissements de deux millions d'euros. Grâce à cette coopération, des actions camerounaises seront aussi échangées à la Bourse de Francfort. "Ce partenariat croissant est une chance pour les deux pays", a déclaré le ministre du Commerce. L'aide au développement sera désormais complétée par un commerce équitable — un signal fort pour l'avenir ! »\par
\textbf{Commentaire des choix de langue :}
\begin{itemize}
\item \textbf{Lexique (8 mots)} : \all{die Wirtschaft}, \all{der Export}, \all{der Handel} (\all{Händler}, \all{Handelsminister}), \all{die Aktie}, \all{die Börse}, \all{die Zusammenarbeit}, \all{die Entwicklungshilfe}, \all{die Investition} — la contrainte des 6 mots est dépassée.
\item \textbf{Prépositions + génitif} : \all{Trotz der weltweiten Krise} ; \all{Während der Verhandlungen} ; \all{Dank dieser Zusammenarbeit} — trois occurrences, génitifs correctement marqués (\all{der}, adjectif en \all{-en}).
\item \textbf{Participe II épithète} : \all{das von beiden Regierungen geförderte Projekt} (action passive achevée) ; \textbf{participe I épithète} : \all{diese wachsende Partnerschaft} et \all{die junge, dynamische Bevölkerung} (action en cours).
\item \textbf{Plusquamperfekt} : \all{hatte ... versprochen} — antériorité par rapport à la signature du contrat.
\item \textbf{Structure journalistique} : titre accrocheur, annonce, détails, citation, ouverture — le modèle de la section 6 est respecté.
\end{itemize}
\end{exoresolu}

\courssub{Bilan de l'activité}

Cette activité a permis de mobiliser simultanément les quatre savoir-faire de la leçon : la \textbf{compréhension} d'un texte économique authentique, l'emploi du \textbf{lexique} de l'aide au développement, du commerce et de la bourse, l'application des règles de \textbf{grammaire} (génitif prépositionnel, participes épithètes) et de \textbf{conjugaison} (temps composés), et enfin la \textbf{production} écrite et orale d'un article de presse. Le vote de la classe, fondé sur la grille d'évaluation de la section 6, développe en outre l'esprit critique et la coopération. Pour aller plus loin, l'article gagnant sera affiché au tableau d'information du lycée et pourra servir de base à un débat : \all{„Ist Entwicklungshilfe noch wichtig, oder ist fairer Handel besser ?“} (L'aide au développement est-elle encore importante, ou le commerce équitable est-il meilleur ?).

\newpage

\coursec{Méthodes et exercices résolus}

\courssub{Méthode 1 : Comprendre un texte économique (global puis détaillé)}

\begin{methode}[Lire et comprendre un texte sur l'aide au développement et la bourse]
\begin{enumerate}
\item \textbf{Première lecture globale} : identifier le thème (\all{Entwicklungshilfe}, \all{Handel}, \all{Börse}), la source (article de presse, reportage, interview) et l'objectif de l'auteur (informer, critiquer, proposer).
\item \textbf{Repérer les mots-clés du champ lexical} : souligner \all{die Entwicklungshilfe}, \all{der Export}, \all{die Aktie}, \all{der Kurs}, \all{der Handelspartner}. Deviner les mots inconnus par les familles de mots (\all{handeln} $\rightarrow$ \all{der Handel} $\rightarrow$ \all{der Händler}).
\item \textbf{Lecture détaillée} : pour chaque question, localiser le passage du texte qui contient la réponse ; ne jamais répondre « de mémoire ».
\item \textbf{Répondre en phrases complètes} en allemand, en reprenant les mots de la question, sans recopier tout le texte.
\item \textbf{Vérifier} : sujet + verbe conjugué en 2\textsuperscript{e} position, majuscules aux noms, cas corrects après les prépositions.
\end{enumerate}
\end{methode}

\begin{exoresolu}[Compréhension de l'écrit : questions sur un texte]
\textbf{Texte (extrait)} : \all{Während der letzten Jahre hat die Entwicklungshilfe Deutschlands für Kamerun zugenommen. Trotz dieser Hilfe bleibt der Handel zwischen beiden Ländern schwach: Kamerun exportiert vor allem Kakao und Kaffee, importiert aber Maschinen und Autos. An der Börse in Frankfurt steigen die Kurse vieler Aktien, doch kleine afrikanische Firmen finden dort kaum Investoren.}

\textbf{Questions} : 1. \all{Was hat in den letzten Jahren zugenommen?} 2. \all{Was exportiert Kamerun?} 3. \all{Welches Problem haben kleine afrikanische Firmen?}
\tcblower
\textbf{Raisonnement} : on localise chaque information dans le texte, puis on reformule en phrase complète.

1. La réponse se trouve dans la première phrase : \all{die Entwicklungshilfe ... hat zugenommen}. Réponse rédigée : \all{In den letzten Jahren hat die Entwicklungshilfe Deutschlands für Kamerun zugenommen.} « Ces dernières années, l'aide au développement de l'Allemagne pour le Cameroun a augmenté. » Attention : \all{zunehmen} est un verbe à particule séparable ; au Perfekt : \all{hat zugenommen}.

2. Deuxième phrase : \all{Kamerun exportiert Kakao und Kaffee.} « Le Cameroun exporte du cacao et du café. »

3. Dernière phrase : \all{Kleine afrikanische Firmen finden an der Börse kaum Investoren.} « Les petites entreprises africaines ne trouvent guère d'investisseurs à la bourse. » \all{kaum} = « à peine, guère » ; ne pas le traduire par « pas ».

\textbf{Conclusion} : chaque réponse reprend les termes de la question, avec le verbe conjugué en 2\textsuperscript{e} position et les majuscules aux noms.
\end{exoresolu}

\courssub{Méthode 2 : Employer les prépositions avec le génitif}

\begin{methode}[Construire une phrase avec une préposition + génitif]
\begin{enumerate}
\item \textbf{Mémoriser la liste} des prépositions principales : \all{während} (pendant), \all{trotz} (malgré), \all{wegen} (à cause de), \all{statt/anstatt} (au lieu de), \all{innerhalb} (à l'intérieur de), \all{außerhalb} (en dehors de), \all{anlässlich} (à l'occasion de).
\item \textbf{Mettre le nom qui suit au génitif} : article \all{des} (masculin/neutre), \all{der} (féminin/pluriel) ; \all{-s} ou \all{-es} au nom masculin ou neutre singulier.
\item \textbf{Placer le groupe prépositionnel} au début (alors verbe en 2\textsuperscript{e} position : \all{Während der Krise steigen die Preise.}) ou après le verbe.
\item \textbf{Astuce familière} : à l'oral, \all{wegen} + datif est fréquent (\all{wegen dem Wetter}), mais à l'écrit et au Bac, exiger le génitif : \all{wegen des Wetters}.
\end{enumerate}
\end{methode}

\begin{exoresolu}[Thème : phrases avec prépositions + génitif]
Traduisez en allemand : 1. « Pendant la crise, beaucoup d'entreprises ont fermé. » 2. « Malgré l'aide au développement, la pauvreté reste un problème. » 3. « À cause des prix élevés, les familles achètent moins. »
\tcblower
\textbf{Raisonnement} : chaque phrase commence par une préposition à génitif ; on construit d'abord le groupe nominal au génitif, puis on place le verbe en 2\textsuperscript{e} position.

1. \all{Während der Krise haben viele Firmen geschlossen.} — \all{die Krise} est féminin : génitif \all{der Krise} (pas de \all{-s} au nom féminin). Groupe en tête $\Rightarrow$ inversion : \all{haben viele Firmen}.

2. \all{Trotz der Entwicklungshilfe bleibt die Armut ein Problem.} — \all{trotz} + génitif féminin \all{der Entwicklungshilfe}. Verbe \all{bleiben} en 2\textsuperscript{e} position. Ne pas traduire « reste » par \all{restet} : le verbe correct est \all{bleiben}.

3. \all{Wegen der hohen Preise kaufen die Familien weniger.} — \all{der Preis} est masculin, mais ici au pluriel : génitif pluriel \all{der Preise} ; adjectif épithète au génitif pluriel : \all{hohen}.

\textbf{Conclusion} : la préposition commande le génitif de tout le groupe nominal (article + adjectif + nom), et le verbe conjugué reste en 2\textsuperscript{e} position.
\end{exoresolu}

\courssub{Méthode 3 : Employer les participes en fonction d'épithète}

\begin{methode}[Transformer un participe en adjectif épithète]
\begin{enumerate}
\item \textbf{Choisir le participe} : participe I (\all{exportierend} = qui exporte, action active en cours) ou participe II (\all{exportiert} = exporté, action passive ou achevée).
\item \textbf{Ajouter les désinences adjectivales} comme pour un adjectif épithète normal : \all{die exportierten Waren}, \all{ein steigender Kurs}, \all{die steigenden Kurse}.
\item \textbf{Possibilité d'un groupe étendu} : le participe peut être précédé de compléments : \all{die von Kamerun exportierten Waren} = « les marchandises exportées par le Cameroun ». En français, on traduit par une relative.
\item \textbf{Vérifier l'accord} en genre, nombre et cas avec le nom : \all{mit den gestiegenen Preisen} (datif pluriel).
\end{enumerate}
\end{methode}

\begin{exoresolu}[Transformation : relatives et participes épithètes]
1. Transformez avec un participe épithète : \all{Die Waren, die von Kamerun exportiert werden, sind Kakao und Kaffee.} 2. Traduisez en français : \all{Die an der Börse gehandelten Aktien sind teuer.} 3. Traduisez en allemand : « Les prix qui montent inquiètent les consommateurs. »
\tcblower
\textbf{Raisonnement} : une relative passive au présent (\all{die ... exportiert werden}) se condense en participe II épithète ; une action active se rend par le participe I.

1. \all{Die von Kamerun exportierten Waren sind Kakao und Kaffee.} — participe II \all{exportiert} + désinence pluriel nominatif \all{-en} (après article défini \all{die}). « Les marchandises exportées par le Cameroun sont le cacao et le café. »

2. \all{Die an der Börse gehandelten Aktien sind teuer.} = « Les actions négociées à la bourse sont chères. » \all{handeln} = « négocier, commercer » ; \all{gehandelt} = participe II.

3. \all{Die steigenden Preise beunruhigen die Konsumenten.} — action active en cours $\Rightarrow$ participe I \all{steigend} + désinence \all{-en}. « Qui montent » = action du sujet « prix », donc participe I, pas II.

\textbf{Conclusion} : actif/en cours $\rightarrow$ participe I ; passif/achevé $\rightarrow$ participe II ; dans les deux cas, désinences adjectivales obligatoires.
\end{exoresolu}

\courssub{Méthode 4 : Maîtriser les temps composés (Perfekt, Plusquamperfekt, Futur II)}

\begin{methode}[Choisir et former un temps composé]
\begin{enumerate}
\item \textbf{Perfekt} : \all{haben/sein} + participe II. Auxiliaire \all{sein} pour les verbes de mouvement ou de changement d'état (\all{steigen, fallen, kommen, sterben}) : \all{Der Kurs ist gestiegen.}
\item \textbf{Plusquamperfekt} (antériorité par rapport au Präteritum) : \all{hatte/war} + participe II : \all{Nachdem die Preise gestiegen waren, kauften die Leute weniger.}
\item \textbf{Futur II} (hypothèse sur le passé, rare) : \all{werden ... haben/sein} + participe II : \all{Bis morgen wird die Firma die Waren exportiert haben.}
\item \textbf{Place} : auxiliaire conjugué en 2\textsuperscript{e} position, participe II en fin de phrase ; dans une subordonnée, l'auxiliaire va à la fin.
\end{enumerate}
\end{methode}

\begin{exoresolu}[Version : temps composés]
Traduisez en français : 1. \all{Die Aktienkurse sind gestern stark gefallen.} 2. \all{Nachdem die Regierung die Hilfe erhöht hatte, begannen neue Projekte.} 3. \all{Deutschland hat viele Maschinen nach Kamerun exportiert.}
\tcblower
\textbf{Raisonnement} : identifier l'auxiliaire (\all{sein/haben}) et le temps, puis choisir l'équivalent français.

1. \all{sind ... gefallen} : Perfekt avec \all{sein} car \all{fallen} exprime un changement d'état. « Les cours des actions ont fortement baissé hier. » \all{fallen} (fort) : \all{fällt, fiel, ist gefallen}.

2. \all{hatte ... erhöht} : Plusquamperfekt $\Rightarrow$ plus-que-parfait français ; \all{begannen} : Präteritum $\Rightarrow$ passé simple. « Après que le gouvernement eut augmenté l'aide, de nouveaux projets commencèrent. » (ou, en style courant : « ont commencé »).

3. \all{hat ... exportiert} : Perfekt avec \all{haben} (verbe transitif). « L'Allemagne a exporté beaucoup de machines vers le Cameroun. » Attention : \all{nach} + pays sans article = « vers ».

\textbf{Conclusion} : Perfekt $\rightarrow$ passé composé ; Plusquamperfekt $\rightarrow$ plus-que-parfait ; l'auxiliaire \all{sein} signale un mouvement ou un changement, jamais une nuance de sens à traduire.
\end{exoresolu}

\courssub{Méthode 5 : Rédiger un résumé et un court paragraphe d'opinion}

\begin{methode}[Résumer puis donner son avis sur un texte économique]
\begin{enumerate}
\item \textbf{Résumé} : relever les 3 ou 4 idées principales ; les reformuler avec ses propres mots ; utiliser des connecteurs : \all{zuerst} (d'abord), \all{dann} (ensuite), \all{außerdem} (de plus), \all{schließlich} (enfin).
\item \textbf{Ne pas recopier} de phrases entières du texte ; interdiction de la première personne dans le résumé.
\item \textbf{Avis personnel} : l'introduire clairement : \all{Meiner Meinung nach ...} (à mon avis), \all{Ich bin der Ansicht, dass ...} (j'estime que), \all{Ich finde es wichtig, dass ...}
\item \textbf{Argumenter} avec \all{denn} (car), \all{weil} + verbe en fin, \all{deshalb} (c'est pourquoi).
\item \textbf{Conclure} en une phrase : \all{Zusammenfassend kann man sagen, dass ...} (en résumé, on peut dire que...).
\end{enumerate}
\end{methode}

\begin{exoresolu}[Production écrite guidée : résumé + avis]
Rédigez en allemand un court texte (6 à 8 phrases) : résumez l'extrait suivant et donnez votre avis sur l'aide au développement. \all{Trotz der Entwicklungshilfe vieler Länder bleibt der Handel mit Afrika schwach. Kamerun exportiert Rohstoffe wie Kakao, importiert aber teure Maschinen. Viele junge Firmen suchen Investoren.}
\tcblower
\textbf{Raisonnement} : 3 idées à résumer (aide insuffisante / export de matières premières / manque d'investisseurs), puis avis argumenté.

\textbf{Modèle de production} : \all{Der Text handelt von der Entwicklungshilfe und dem Handel mit Afrika. Zuerst erklärt der Autor, dass der Handel trotz der Hilfe vieler Länder schwach bleibt. Dann zeigt er, dass Kamerun Rohstoffe wie Kakao exportiert, aber teure Maschinen importieren muss. Außerdem suchen viele junge Firmen Investoren. Meiner Meinung nach ist die Entwicklungshilfe wichtig, aber sie reicht nicht. Ich bin der Ansicht, dass afrikanische Länder mehr Produkte selbst herstellen müssen, denn so entstehen Arbeitsplätze. Deshalb sollten auch die Börsen mehr kleinen Firmen helfen. Zusammenfassend kann man sagen, dass fairer Handel besser ist als nur Hilfe.}

« Le texte traite de l'aide au développement et du commerce avec l'Afrique. D'abord, l'auteur explique que le commerce reste faible malgré l'aide de nombreux pays. Ensuite, il montre que le Cameroun exporte des matières premières comme le cacao mais doit importer des machines coûteuses. De plus, beaucoup de jeunes entreprises cherchent des investisseurs. À mon avis, l'aide au développement est importante, mais elle ne suffit pas. J'estime que les pays africains doivent fabriquer davantage de produits eux-mêmes, car cela crée des emplois. C'est pourquoi les bourses devraient aussi davantage aider les petites entreprises. En résumé, on peut dire qu'un commerce équitable vaut mieux que la seule aide. »

\textbf{Points grammaticaux vérifiés} : \all{trotz der Hilfe} (génitif), subordonnées avec \all{dass} et \all{weil/denn} correctement construites, verbe conjugué en 2\textsuperscript{e} position après \all{Meiner Meinung nach}, majuscules à tous les noms.
\end{exoresolu}

\begin{attention}[Les 5 erreurs qui coûtent des points au Bac]
\begin{enumerate}
\item \textbf{Datif au lieu du génitif} : écrire \all{trotz dem Handel} est une faute à l'écrit. Écrire \all{trotz des Handels}. Après \all{während, trotz, wegen, statt} : toujours le génitif au Bac.
\item \textbf{Oublier la désinence du participe épithète} : \all{die exportiert Waren} est faux ; il faut \all{die exportierten Waren}. Un participe devant un nom se décline comme un adjectif.
\item \textbf{Faux amis} : \all{die Aktie} = « l'action » (en bourse), pas « l'action » au sens d'acte ; \all{der Kurs} = « le cours » (de la bourse, d'une monnaie), pas « le cours » scolaire (\all{der Unterricht}) ; \all{die Fabrik} = « l'usine », pas « la fabrique » de papier.
\item \textbf{Mauvais auxiliaire au Perfekt} : \all{Der Kurs hat gestiegen} est faux ; les verbes de changement d'état prennent \all{sein} : \all{Der Kurs ist gestiegen.} Même piège avec \all{fallen} : \all{ist gefallen}.
\item \textbf{Place du verbe après un groupe en tête} : après \all{Trotz der Entwicklungshilfe}, le verbe vient immédiatement : \all{Trotz der Entwicklungshilfe bleibt die Armut ein Problem.} Ne jamais calquer l'ordre français « Malgré l'aide, la pauvreté reste... » qui placerait le sujet avant le verbe. Et majuscule à tous les noms : \all{die Hilfe, der Handel, die Börse}.
\end{enumerate}
\end{attention}

\newpage

\coursec{Exercices}

\courssub{Je vérifie mes connaissances}

\begin{exercice}[QCM : vocabulaire économique]
Pour chaque question, entoure la bonne réponse.

\begin{enumerate}
\item \all{Die Börse} est :
\begin{enumerate}
\item une banque centrale
\item un marché où l'on achète et vend des actions
\item un ministère des Finances
\end{enumerate}
\item \all{Die Entwicklungshilfe} signifie :
\begin{enumerate}
\item l'aide au développement
\item le commerce extérieur
\item la dette publique
\end{enumerate}
\item Le pluriel de \all{die Aktie} est :
\begin{enumerate}
\item die Aktien
\item die Akties
\item die Aktier
\end{enumerate}
\item \all{Exportieren} a pour contraire :
\begin{enumerate}
\item produzieren
\item importieren
\item investieren
\end{enumerate}
\item \all{Der Handel} désigne :
\begin{enumerate}
\item le marché aux puces
\item le commerce, les échanges
\item la monnaie
\end{enumerate}
\end{enumerate}
\end{exercice}

\begin{exercice}[Vrai ou faux ?]
Indique si les affirmations suivantes sont vraies (V) ou fausses (F) et justifie ta réponse en une phrase en français.

\begin{enumerate}
\item La préposition \all{wegen} est toujours suivie du datif.
\item \all{Trotz des Regens} signifie « malgré la pluie ».
\item Dans \all{die gestern unterzeichneten Verträge}, \all{unterzeichneten} est un participe passé employé comme épithète.
\item Le Plusquamperfekt exprime une action antérieure à une action passée.
\item \all{Während} se construit avec l'accusatif.
\end{enumerate}
\end{exercice}

\begin{exercice}[Familles de mots]
Complète le tableau avec le mot de la même famille qui manque (nom, verbe ou adjectif). Donne l'article et le pluriel des noms.

\begin{center}
\begin{tabularx}{\linewidth}{|X|X|X|}
\hline
\rowcolor{popblueL} Nom & Verbe & Adjectif \\
\hline
die Entwicklung, -en & ................ & entwickelt \\
\hline
................ & handeln & ................ \\
\hline
der Export, -e & ................ & ................ \\
\hline
................ & investieren & ................ \\
\hline
die Kooperation, -en & ................ & kooperativ \\
\hline
\end{tabularx}
\end{center}
\end{exercice}

\begin{exercice}[Conjugaison : temps composés]
Mets le verbe entre parenthèses au temps composé demandé.

\begin{enumerate}
\item Die Aktienkurse (steigen — Perfekt) gestern stark.
\item Die Regierung (unterzeichnen — Plusquamperfekt) den Vertrag schon vor der Krise.
\item Viele Investoren (verlieren — Perfekt) viel Geld an der Börse.
\item Kamerun (exportieren — Plusquamperfekt) schon vor 2000 viel Kakao nach Deutschland.
\item Die Entwicklungshilfe (helfen — Perfekt) vielen Dörfern in der Region.
\end{enumerate}
\end{exercice}

\courssub{Je m'entraîne}

\begin{exercice}[Prépositions avec le génitif]
Complète les huit phrases avec la préposition correcte (\all{wegen}, \all{trotz}, \all{während}, \all{aufgrund}) et la bonne désinence du génitif.

\begin{enumerate}
\item ................ d............ Krise sind die Aktienkurse gefallen.
\item ................ d............ starken Regens konnte die Ware nicht transportiert werden.
\item ................ d............ Verhandlungen haben die Minister nichts gesagt.
\item ................ d............ guten Ernte steigen die Exporte von Kakao.
\item ................ d............ hohen Preise kaufen viele Bauern neue Maschinen.
\item ................ d............ Konferenz in Yaoundé wurden drei Verträge unterzeichnet.
\item ................ d............ schlechten Straßen kommen die Lastwagen zu spät an.
\item ................ d............ neuen Abkommens wird der Handel leichter.
\end{enumerate}
\end{exercice}

\begin{exercice}[Participes épithètes : transformations]
Transforme la subordonnée relative en participe épithète, comme dans le modèle.

Modèle : \all{der Vertrag, der gestern unterzeichnet wurde} $\rightarrow$ \all{der gestern unterzeichnete Vertrag}

\begin{enumerate}
\item die Aktien, die an der Börse gehandelt werden
\item der Markt, der sich schnell entwickelt
\item die Waren, die aus Kamerun exportiert werden
\item die Zusammenarbeit, die seit Jahren besteht
\item die Investoren, die in Douala investieren
\end{enumerate}
\end{exercice}

\begin{exercice}[Texte à trous : la Bourse de Francfort]
Complète le texte avec les mots de la liste suivante : \all{Aktien — Börse — gefallen — gestiegen — Investoren — Kurs — Gewinn — gehandelt}.

\begin{textebox}[An der Frankfurter Börse]
An der Frankfurter ................ werden jeden Tag Millionen von ................ gehandelt. Gestern sind die Kurse vieler Unternehmen stark ................, aber heute sind sie wieder ................ Viele ................ hoffen auf einen großen ................ Der ................ der Aktie einer deutschen Bank hat in einer Woche zehn Euro gewonnen. An der Börse wird von Montag bis Freitag ................
\end{textebox}
\end{exercice}

\begin{exercice}[Appariements]
Relie chaque mot allemand à sa traduction française.

\begin{center}
\begin{tabularx}{\linewidth}{|X|X|}
\hline
\rowcolor{popblueL} Deutsch & Français \\
\hline
1. die Entwicklungshilfe & a. le taux de change \\
\hline
2. der Wechselkurs & b. la matière première \\
\hline
3. der Rohstoff & c. l'aide au développement \\
\hline
4. das Abkommen & d. la balance commerciale \\
\hline
5. die Handelsbilanz & e. l'accord, le traité \\
\hline
\end{tabularx}
\end{center}
\end{exercice}

\courssub{Je me prépare au Bac}

\begin{exercice}[Compréhension écrite type Bac]
Lis le texte suivant, puis réponds aux questions en allemand par des phrases complètes.

\begin{textebox}[Kamerun exportiert mehr nach Deutschland]
Der Handel zwischen Kamerun und Deutschland entwickelt sich positiv. Während der letzten Jahre hat Kamerun trotz einiger Schwierigkeiten immer mehr Rohstoffe nach Deutschland exportiert: Kakao, Kaffee, Holz und Baumwolle. Die gestern in Yaoundé unterzeichneten Verträge werden diese Zusammenarbeit noch verstärken. Aufgrund der neuen Abkommen sollen auch Maschinen und Autos aus Deutschland leichter nach Kamerun importiert werden. Viele Experten meinen, dass die wachsende Wirtschaft Kameruns neue Arbeitsplätze schaffen wird. Wegen der steigenden Preise an der Börse interessieren sich auch deutsche Investoren für den camerounais Markt. Die Entwicklungshilfe Deutschlands unterstützt außerdem Projekte in den Bereichen Bildung und Gesundheit. Die von beiden Regierungen geplanten Projekte sollen schon nächstes Jahr beginnen.
\end{textebox}

\textbf{Questions de compréhension :}
\begin{enumerate}
\item Welche Rohstoffe exportiert Kamerun nach Deutschland ? (Nenne vier.)
\item Was wurde gestern in Yaoundé unterzeichnet, und welche Folge hat das ?
\item Was soll dank der neuen Abkommen leichter importiert werden ?
\item Warum interessieren sich deutsche Investoren für Kamerun ?
\item In welchen Bereichen unterstützt die deutsche Entwicklungshilfe Projekte ?
\end{enumerate}

\textbf{Vrai ou faux ? Justifie ta réponse par une citation du texte.}
\begin{enumerate}
\item Kamerun exportiert nur Kakao nach Deutschland.
\item Die Projekte sollen erst in zehn Jahren beginnen.
\end{enumerate}

\textbf{Questions de langue :}
\begin{enumerate}
\item Relève dans le texte deux participes employés comme épithètes et traduis-les en français.
\item Relève deux prépositions suivies du génitif et recopie le groupe prépositionnel complet.
\end{enumerate}
\end{exercice}

\begin{exercice}[Thème et version type Bac]
\textbf{Thème.} Traduis en allemand les phrases suivantes.

\begin{enumerate}
\item Malgré la crise, les cours des actions ont monté et les contrats signés hier renforceront la coopération.
\item À cause des prix élevés, le Cameroun a exporté moins de café cette année.
\item Les produits importés d'Allemagne sont souvent plus chers que les produits locaux.
\item Pendant les négociations, les deux ministres ont signé un accord important.
\end{enumerate}

\textbf{Version.} Traduis en français le paragraphe suivant.

\begin{textebox}[Version]
Die an der Börse gehandelten Aktien hatten am Montag stark verloren. Viele Investoren hatten schon vor der Krise ihr Geld aus dem Land gebracht. Die gestern beschlossenen Maßnahmen der Regierung sollen den Markt jetzt stabilisieren. Trotz der schwierigen Lage hoffen alle auf bessere Zeiten.
\end{textebox}
\end{exercice}

\begin{exercice}[Expression écrite type Bac]
\textbf{Sujet :} \all{Soll Kamerun mehr Rohstoffe nach Deutschland exportieren ?}

Rédige en allemand un paragraphe d'opinion de \textbf{8 à 10 lignes} (environ 80 à 100 mots) en respectant le plan imposé suivant :

\begin{enumerate}
\item \textbf{Introduction :} présente le sujet en une ou deux phrases (le commerce entre le Cameroun et l'Allemagne).
\item \textbf{Arguments pour :} donne au moins deux avantages de l'exportation des matières premières (argent, emplois, coopération).
\item \textbf{Arguments contre ou limites :} donne au moins un inconvénient (dépendance, prix variables, transformation sur place).
\item \textbf{Conclusion :} exprime ton opinion personnelle avec \all{Ich meine, dass ...} ou \all{Meiner Meinung nach ...}.
\end{enumerate}

\textbf{Contraintes de langue :} emploie au moins une préposition avec le génitif (\all{trotz}, \all{wegen}, \all{aufgrund}, \all{während}), au moins un participe épithète et au moins un verbe au Perfekt.
\end{exercice}

\newpage

\coursec{Corrigés des exercices}

\coursec{Corrigés des exercices de la leçon AL09}

\begin{corrige}[Exercice 1 — QCM : vocabulaire économique]
\textbf{1. b} — \all{Die Börse} est un marché où l'on achète et vend des actions. \all{An der Börse werden Aktien gehandelt.} « À la Bourse, on négocie des actions. »

\textbf{2. a} — \all{Die Entwicklungshilfe} signifie « l'aide au développement ». Attention au faux ami : \all{Hilfe} = aide, pas « hiver » ni « fil ».

\textbf{3. a} — Le pluriel de \all{die Aktie} est \all{die Aktien} (féminin en \all{-ie} $\rightarrow$ pluriel en \all{-n}).

\textbf{4. b} — Le contraire de \all{exportieren} est \all{importieren}. \all{Kamerun exportiert Kakao und importiert Maschinen.} « Le Cameroun exporte du cacao et importe des machines. »

\textbf{5. b} — \all{Der Handel} désigne le commerce, les échanges. Famille : \all{handeln} (commercer), \all{der Händler} (le commerçant), \all{die Handelsbilanz} (la balance commerciale).
\end{corrige}

\begin{corrige}[Exercice 2 — Vrai ou faux ?]
\textbf{1. F} — Faux. \all{Wegen} est une préposition du \cle{génitif} : \all{wegen der Krise} (à cause de la crise). Le datif après \all{wegen} n'apparaît qu'à l'oral familier, jamais dans un texte écrit soigné.

\textbf{2. V} — Vrai. \all{Trotz des Regens} signifie bien « malgré la pluie » : \all{trotz} + génitif (\all{der Regen} $\rightarrow$ \all{des Regens}).

\textbf{3. V} — Vrai. \all{Unterzeichneten} est le participe passé de \all{unterzeichnen} (verbe à préfixe inséparable, donc sans \all{ge-}), employé comme épithète avec la désinence adjectivale \all{-en} (pluriel après article défini).

\textbf{4. V} — Vrai. Le Plusquamperfekt exprime l'antériorité par rapport à un moment passé : \all{Die Regierung hatte den Vertrag unterzeichnet, bevor die Krise begann.} « Le gouvernement avait signé le contrat avant que la crise ne commence. »

\textbf{5. F} — Faux. \all{Während} se construit avec le \cle{génitif} : \all{während der Verhandlungen} (pendant les négociations).
\end{corrige}

\begin{corrige}[Exercice 3 — Familles de mots]
\begin{center}
\begin{tabularx}{\linewidth}{|X|X|X|}
\hline
\rowcolor{popblueL} Nom & Verbe & Adjectif \\
\hline
die Entwicklung, -en & entwickeln & entwickelt \\
\hline
der Handel (pas de pluriel usuel) / die Handlung, -en & handeln & handelbar / kommerziell \\
\hline
der Export, -e & exportieren & exportiert \\
\hline
die Investition, -en & investieren & investiert \\
\hline
die Kooperation, -en & kooperieren & kooperativ \\
\hline
\end{tabularx}
\end{center}

\textbf{Remarques :} pour \all{handeln}, le nom attendu est \all{der Handel} (pas de pluriel usuel dans le sens de « commerce » ; \all{die Handlung} signifie « action, intrigue » et prend \all{-en} au pluriel). Pour \all{investieren}, le nom est \all{die Investition} (féminin en \all{-tion}, pluriel \all{-tionen}). Les verbes en \all{-ieren} forment leur participe passé sans \all{ge-} : \all{exportiert}, \all{investiert}, \all{kooperiert}.
\end{corrige}

\begin{corrige}[Exercice 4 — Conjugaison : temps composés]
\textbf{1.} \all{Die Aktienkurse \textbf{sind gestern stark gestiegen}.} — \all{Steigen} est un verbe de mouvement/changement d'état $\rightarrow$ auxiliaire \all{sein} ; participe passé fort \all{gestiegen}.

\textbf{2.} \all{Die Regierung \textbf{hatte den Vertrag schon vor der Krise unterzeichnet}.} — Plusquamperfekt : \all{hatte} + participe passé ; \all{unterzeichnen} est inséparable, donc \all{unterzeichnet} sans \all{ge-}.

\textbf{3.} \all{Viele Investoren \textbf{haben viel Geld an der Börse verloren}.} — Auxiliaire \all{haben} ; participe passé fort de \all{verlieren} : \all{verloren}.

\textbf{4.} \all{Kamerun \textbf{hatte schon vor 2000 viel Kakao nach Deutschland exportiert}.} — Verbe en \all{-ieren} : participe \all{exportiert} sans \all{ge-}.

\textbf{5.} \all{Die Entwicklungshilfe \textbf{hat vielen Dörfern in der Region geholfen}.} — \all{Helfen} est un verbe fort : participe \all{geholfen} ; attention, \all{helfen} se construit avec le datif : \all{vielen Dörfern}.

\textbf{Point de vigilance :} au Perfekt et au Plusquamperfekt, le participe passé se place en \cle{fin de phrase} ; seul l'auxiliaire est conjugué en deuxième position.
\end{corrige}

\begin{corrige}[Exercice 5 — Prépositions avec le génitif]
\textbf{1.} \all{\textbf{Wegen der} Krise sind die Aktienkurse gefallen.} « À cause de la crise, les cours des actions ont baissé. » (\all{die Krise} $\rightarrow$ génitif \all{der Krise})

\textbf{2.} \all{\textbf{Wegen des} starken Regens konnte die Ware nicht transportiert werden.} (\all{der Regen} $\rightarrow$ \all{des Regens} ; adjectif épithète : \all{starken})

\textbf{3.} \all{\textbf{Während der} Verhandlungen haben die Minister nichts gesagt.} « Pendant les négociations, les ministres n'ont rien dit. » (pluriel : \all{der Verhandlungen})

\textbf{4.} \all{\textbf{Aufgrund der} guten Ernte steigen die Exporte von Kakao.} « Grâce à la bonne récolte, les exportations de cacao augmentent. »

\textbf{5.} \all{\textbf{Aufgrund der} hohen Preise kaufen viele Bauern neue Maschinen.} (\all{die Preise} $\rightarrow$ pluriel : \all{der hohen Preise})

\textbf{6.} \all{\textbf{Während der} Konferenz in Yaoundé wurden drei Verträge unterzeichnet.} « Pendant la conférence de Yaoundé, trois contrats ont été signés. »

\textbf{7.} \all{\textbf{Wegen der} schlechten Straßen kommen die Lastwagen zu spät an.} (\all{die Straße}, pluriel \all{die Straßen} $\rightarrow$ génitif \all{der Straßen})

\textbf{8.} \all{\textbf{Aufgrund des} neuen Abkommens wird der Handel leichter.} (\all{das Abkommen} $\rightarrow$ \all{des Abkommens} ; adjectif : \all{neuen})

\textbf{Point de vigilance :} au génitif, les noms masculins et neutres prennent \all{-(e)s} (\all{des Regens}, \all{des Abkommens}) ; les féminins et les pluriels restent inchangés (\all{der Krise}, \all{der Verhandlungen}).
\end{corrige}

\begin{corrige}[Exercice 6 — Participes épithètes : transformations]
\textbf{1.} \all{die an der Börse \textbf{gehandelten} Aktien} « les actions négociées à la Bourse » (participe passé, passif).

\textbf{2.} \all{der sich schnell \textbf{entwickelnde} Markt} « le marché qui se développe rapidement » (participe présent \all{entwickelnd} + désinence \all{-e}, car le verbe est actif, masc. sing. nom. après article défini).

\textbf{3.} \all{die aus Kamerun \textbf{exportierten} Waren} « les marchandises exportées du Cameroun » (participe passé sans \all{ge-} car verbe en \all{-ieren}).

\textbf{4.} \all{die seit Jahren \textbf{bestehende} Zusammenarbeit} « la coopération qui existe depuis des années » (participe présent, action active, fém. sing. nom./acc. après article défini).

\textbf{5.} \all{die in Douala \textbf{investierenden} Investoren} « les investisseurs qui investissent à Douala » (participe présent, pluriel après article défini $\rightarrow$ \all{-en}).

\textbf{Règle :} la relative au \cle{passif} (\all{die ... gehandelt werden}) devient un \cle{participe passé} épithète ; la relative à l'\cle{actif} (\all{der sich entwickelt}) devient un \cle{participe présent} (\all{-end}). Dans les deux cas, le participe reçoit la désinence de l'adjectif épithète selon le genre, le nombre, le cas et le type d'article.
\end{corrige}

\begin{corrige}[Exercice 7 — Texte à trous : la Bourse de Francfort]
\begin{textebox}[Texte complété]
An der Frankfurter \textbf{Börse} werden jeden Tag Millionen von \textbf{Aktien} gehandelt. Gestern sind die Kurse vieler Unternehmen stark \textbf{gefallen}, aber heute sind sie wieder \textbf{gestiegen}. Viele \textbf{Investoren} hoffen auf einen großen \textbf{Gewinn}. Der \textbf{Kurs} der Aktie einer deutschen Bank hat in einer Woche zehn Euro gewonnen. An der Börse wird von Montag bis Freitag \textbf{gehandelt}.
\end{textebox}

\textbf{Traduction :} « À la Bourse de Francfort, des millions d'actions sont négociées chaque jour. Hier, les cours de nombreuses entreprises ont fortement baissé, mais aujourd'hui ils sont remontés. Beaucoup d'investisseurs espèrent un gros gain. Le cours de l'action d'une banque allemande a gagné dix euros en une semaine. À la Bourse, on négocie du lundi au vendredi. »

\textbf{Justifications :} \all{gefallen}/\all{gestiegen} : participes passés avec \all{sein} (changement d'état) ; \all{gehandelt} : participe passé du passif avec \all{werden} (\all{wird ... gehandelt}).
\end{corrige}

\begin{corrige}[Exercice 8 — Appariements]
\begin{center}
\begin{tabularx}{\linewidth}{|X|X|}
\hline
\rowcolor{popblueL} Deutsch & Français \\
\hline
1. die Entwicklungshilfe & \textbf{c.} l'aide au développement \\
\hline
2. der Wechselkurs & \textbf{a.} le taux de change \\
\hline
3. der Rohstoff & \textbf{b.} la matière première \\
\hline
4. das Abkommen & \textbf{e.} l'accord, le traité \\
\hline
5. die Handelsbilanz & \textbf{d.} la balance commerciale \\
\hline
\end{tabularx}
\end{center}

\textbf{Astuce de mémorisation :} \all{wechseln} = changer $\rightarrow$ \all{der Wechselkurs} = le cours (taux) du change ; \all{roh} = brut $\rightarrow$ \all{der Rohstoff} = la matière première (brute) ; \all{die Bilanz} = le bilan $\rightarrow$ \all{die Handelsbilanz} = le bilan (la balance) du commerce.
\end{corrige}

\begin{corrige}[Exercice 9 — Compréhension écrite type Bac]
\textbf{Barème indicatif : 10 points} — 5 questions de compréhension (5 pts), vrai/faux justifié (2 pts), questions de langue (3 pts).

\textbf{Questions de compréhension — réponses modèles :}

\textbf{1.} \all{Kamerun exportiert Kakao, Kaffee, Holz und Baumwolle nach Deutschland.} « Le Cameroun exporte du cacao, du café, du bois et du coton vers l'Allemagne. »

\textbf{2.} \all{Gestern wurden in Yaoundé Verträge unterzeichnet. Diese Verträge werden die Zusammenarbeit noch verstärken.} « Hier, des contrats ont été signés à Yaoundé. Ces contrats renforceront encore la coopération. »

\textbf{3.} \all{Dank der neuen Abkommen sollen Maschinen und Autos aus Deutschland leichter nach Kamerun importiert werden.} « Grâce aux nouveaux accords, des machines et des voitures d'Allemagne doivent être importées plus facilement au Cameroun. »

\textbf{4.} \all{Deutsche Investoren interessieren sich für Kamerun wegen der steigenden Preise an der Börse und wegen der wachsenden Wirtschaft.} « Les investisseurs allemands s'intéressent au Cameroun à cause de la hausse des prix à la Bourse et de la croissance de l'économie. »

\textbf{5.} \all{Die deutsche Entwicklungshilfe unterstützt Projekte in den Bereichen Bildung und Gesundheit.} « L'aide au développement allemande soutient des projets dans les domaines de l'éducation et de la santé. »

\textbf{Vrai ou faux :}

\textbf{1. Faux.} Citation : \all{„Kakao, Kaffee, Holz und Baumwolle“} — le Cameroun exporte quatre matières premières, pas seulement le cacao.

\textbf{2. Faux.} Citation : \all{„Die ... geplanten Projekte sollen schon nächstes Jahr beginnen.“} — les projets doivent commencer dès l'année prochaine, pas dans dix ans.

\textbf{Questions de langue :}

\textbf{1.} Participes épithètes : \all{die gestern in Yaoundé \textbf{unterzeichneten} Verträge} « les contrats signés hier à Yaoundé » ; \all{die von beiden Regierungen \textbf{geplanten} Projekte} « les projets planifiés par les deux gouvernements » ; on accepte aussi \all{die \textbf{wachsende} Wirtschaft} « l'économie croissante » (participe présent).

\textbf{2.} Prépositions + génitif : \all{\textbf{während der} letzten Jahre} (pendant les dernières années) ; \all{\textbf{trotz} einiger Schwierigkeiten} (malgré quelques difficultés) ; \all{\textbf{aufgrund der} neuen Abkommen} (en raison des nouveaux accords) ; \all{\textbf{wegen der} steigenden Preise} (à cause des prix en hausse).

\textbf{Points de vigilance :} répondre par des phrases complètes avec le verbe conjugué en deuxième position ; reprendre le vocabulaire du texte sans le recopier mot à mot en dehors des citations ; pour le vrai/faux, la citation entre guillemets allemands „…“ est obligatoire pour obtenir le point de justification.
\end{corrige}

\begin{corrige}[Exercice 10 — Thème et version type Bac]
\textbf{Barème indicatif : 10 points} — thème : 4 phrases à 1,5 pt (6 pts) ; version : 4 pts.

\textbf{Thème — traductions proposées :}

\textbf{1.} \all{Trotz der Krise sind die Aktienkurse gestiegen, und die gestern unterzeichneten Verträge werden die Zusammenarbeit verstärken.} — \all{trotz} + génitif ; \all{steigen} au Perfekt avec \all{sein} ; participe épithète \all{unterzeichneten}.

\textbf{2.} \all{Wegen der hohen Preise hat Kamerun dieses Jahr weniger Kaffee exportiert.} — \all{wegen} + génitif pluriel \all{der hohen Preise} ; \all{exportiert} sans \all{ge-}.

\textbf{3.} \all{Die aus Deutschland importierten Produkte sind oft teurer als die lokalen Produkte.} — participe passé épithète \all{importierten} ; comparatif \all{teurer als}.

\textbf{4.} \all{Während der Verhandlungen haben die beiden Minister ein wichtiges Abkommen unterzeichnet.} — \all{während} + génitif ; adjectif épithète neutre : \all{ein wichtiges Abkommen}.

\textbf{Version — traduction proposée :}

« Les actions négociées à la Bourse avaient fortement perdu lundi. De nombreux investisseurs avaient déjà sorti leur argent du pays avant la crise. Les mesures décidées hier par le gouvernement doivent maintenant stabiliser le marché. Malgré la situation difficile, tous espèrent des jours meilleurs. »

\textbf{Points de vigilance :} \all{hatten ... verloren} et \all{hatten ... gebracht} sont des Plusquamperfekt $\rightarrow$ plus-que-parfait français (« avaient perdu », « avaient sorti ») ; \all{die gestern beschlossenen Maßnahmen} : participe épithète à rendre par une relative ou un participe passé (« les mesures décidées hier ») ; \all{sollen} exprime ici la destination/le but (« doivent », « sont censées »), pas « devraient » (ce serait \all{sollten} = Konjunktiv II).
\end{corrige}

\begin{corrige}[Exercice 11 — Expression écrite type Bac]
\textbf{Barème indicatif : 10 points} — respect du plan en 4 parties (2 pts) ; contraintes de langue : génitif prépositionnel, participe épithète, Perfekt (3 pts) ; correction grammaticale et ordre des mots (3 pts) ; richesse du vocabulaire thématique (2 pts).

\textbf{Proposition de rédaction modèle :}

\begin{textebox}[Soll Kamerun mehr Rohstoffe exportieren ?]
Der Handel zwischen Kamerun und Deutschland hat sich in den letzten Jahren stark entwickelt. Kamerun exportiert vor allem Kakao, Kaffee und Holz. Der Export von Rohstoffen hat viele Vorteile : er bringt dem Land Geld und schafft neue Arbeitsplätze. Au\ss{}erdem stärkt er die Zusammenarbeit mit Deutschland. Aber es gibt auch Probleme : wegen der schwankenden Preise an der Börse ist das Einkommen der Bauern nicht sicher. Kamerun ist au\ss{}erdem von den importierten Produkten abhängig. Meiner Meinung nach \textbf{sollte} Kamerun seine Rohstoffe vor dem Export im Land verarbeiten. So würde die Wirtschaft noch stärker wachsen.
\end{textebox}

\textbf{Traduction :} « Le commerce entre le Cameroun et l'Allemagne s'est fortement développé ces dernières années. Le Cameroun exporte surtout du cacao, du café et du bois. L'exportation de matières premières a de nombreux avantages : elle rapporte de l'argent au pays et crée de nouveaux emplois. De plus, elle renforce la coopération avec l'Allemagne. Mais il y a aussi des problèmes : à cause des prix fluctuants à la Bourse, le revenu des paysans n'est pas sûr. Le Cameroun dépend en outre des produits importés. À mon avis, le Cameroun \textbf{devrait} transformer ses matières premières sur place avant l'exportation. Ainsi, l'économie croîtrait encore plus fortement. »

\textbf{Vérification des contraintes :} préposition + génitif : \all{wegen der schwankenden Preise} ; participe épithète : \all{die importierten Produkte} ; Perfekt : \all{hat sich ... entwickelt} ; Konjunktiv II pour le conseil : \all{sollte ... verarbeiten / würde ... wachsen}.

\textbf{Points de vigilance :} verbe conjugué toujours en deuxième position, y compris après \all{au\ss{}erdem} et \all{so} (\all{So würde die Wirtschaft ... wachsen}) ; majuscule à tous les noms (\all{Export}, \all{Arbeitsplätze}, \all{Wirtschaft}, \all{Rohstoffe}) ; \all{Meiner Meinung nach} + verbe en position 2 ; éviter le calque du français « selon mon opinion » ; \all{au\ss{}erdem} s'écrit avec \all{\ss} (après voyelle longue/diphtongue).
\end{corrige}

\newpage

\coursec{Fiche bilan}

\begin{popfigure}
\begin{center}
\begin{tikzpicture}[mindmap, grow cyclic, every node/.style={concept, font=\small}, text width=2.6cm, align=center, level 1/.append style={sibling angle=60, level distance=3.6cm}]
\node[concept color=popblue!70, text=white, font=\small\bfseries] {Aide au développement, échanges commerciaux, bourse}
  child[concept color=popgreen!70] { node {Lexique : \all{die Entwicklungshilfe}, \all{der Handel}} }
  child[concept color=poporange!80] { node {Bourse : \all{die Börse}, \all{die Aktie}, \all{der Kurs}} }
  child[concept color=poppurple!70, text=white] { node {Prépositions + génitif : \all{wegen}, \all{trotz}, \all{während}} }
  child[concept color=poppink!80] { node {Participes épithètes : \all{der exportierte Kaffee}} }
  child[concept color=popgold!85] { node {Temps composés : Perfekt, Plusquamperfekt, Futur II} }
  child[concept color=popblue!50] { node {Méthode : lire, résumer, commenter} };
\end{tikzpicture}
\end{center}
\legende{Carte mentale de la leçon AL09}
\end{popfigure}

\begin{bilan}
\textbf{1. Lexique clé (à connaître avec article et pluriel)}
\begin{itemize}
\item \all{die Entwicklungshilfe, -n} : l'aide au développement ; \all{das Entwicklungsland, -länder} : le pays en développement ; \all{das Industrieland, -länder} : le pays industrialisé.
\item \all{der Handel} : le commerce ; \all{der Export, -e} : l'exportation ; \all{der Import, -e} : l'importation ; \all{exportieren / importieren} : exporter / importer ; \all{der Handelspartner, -} : le partenaire commercial.
\item \all{die Börse, -n} : la bourse ; \all{die Aktie, -n} : l'action ; \all{der Kurs, -e} : le cours ; \all{die Aktiengesellschaft (AG), -en} : la société anonyme ; \all{steigen / fallen} (cours) : monter / baisser ; \all{investieren} : investir.
\item Familles de mots : \all{entwickeln} $\rightarrow$ \all{die Entwicklung} $\rightarrow$ \all{die Entwicklungshilfe} ; \all{handeln} $\rightarrow$ \all{der Handel} $\rightarrow$ \all{der Händler}.
\end{itemize}

\textbf{2. Grammaire à connaître par cœur}
\begin{center}
\begin{tabularx}{\linewidth}{|l|X|X|}
\hline
\rowcolor{popblueL} Point & Règle & Exemple \\
\hline
Prépositions + génitif & \all{während, wegen, trotz, statt, innerhalb, außerhalb} + génitif (\all{des, der}) & \all{wegen des hohen Preises} (à cause du prix élevé) \\
\hline
Participe I épithète & participe présent + \all{-d} + désinence d'adjectif & \all{die steigenden Kurse} (les cours qui montent) \\
\hline
Participe II épithète & participe passé + désinence d'adjectif (sens passif ou accompli) & \all{der exportierte Kakao} (le cacao exporté) \\
\hline
Perfekt & \all{haben/sein} + participe II (participe en fin de phrase) & \all{Die Kurse sind gestiegen.} \\
\hline
Plusquamperfekt & \all{hatte/war} + participe II (antériorité) & \all{Die Firma hatte investiert.} \\
\hline
Futur II & \all{werden} + participe II + \all{haben/sein} (supposition sur le passé) & \all{Er wird die Aktien verkauft haben.} \\
\hline
\end{tabularx}
\end{center}

\textbf{3. Méthodes express}
\begin{itemize}
\item Compréhension : lire le titre et la source, repérer le thème, faire une lecture globale puis détaillée, souligner les mots-clés du champ lexical économique.
\item Résumé : une phrase par idée principale, au présent, sans recopier le texte, avec des connecteurs (\all{zuerst, dann, außerdem, schließlich}).
\item Traduction : repérer d'abord le verbe conjugué (2\up{e} place), puis le sujet, puis les compléments ; attention aux participes épithètes placés avant le nom.
\end{itemize}
\end{bilan}

\begin{aretenir}[Checklist avant le Bac]
\begin{itemize}
\item[$\square$] Je connais l'article et le pluriel des noms clés : \all{die Entwicklungshilfe}, \all{der Handel}, \all{die Börse}, \all{die Aktie}, \all{der Export}, \all{der Import}.
\item[$\square$] Je cite sans hésiter les six prépositions suivies du génitif : \all{während, wegen, trotz, statt, innerhalb, außerhalb}.
\item[$\square$] Je forme correctement le génitif : \all{des Marktes}, \all{der Aktie}, \all{der Länder}.
\item[$\square$] Je sais que \all{trotz} et \all{wegen} prennent le génitif à l'écrit (le datif est familier).
\item[$\square$] Je transforme une subordonnée en participe épithète : \all{die Kurse, die steigen} $\rightarrow$ \all{die steigenden Kurse}.
\item[$\square$] J'accorde le participe épithète comme un adjectif (genre, nombre, cas).
\item[$\square$] Je distingue participe I (actif, simultané) et participe II (passif ou accompli).
\item[$\square$] Je conjugue le Perfekt avec le bon auxiliaire (\all{sein} pour \all{steigen, fallen}) et le participe II en fin de phrase.
\item[$\square$] J'exprime l'antériorité avec le Plusquamperfekt : \all{nachdem die Preise gefallen waren}.
\item[$\square$] Je place le verbe conjugué en 2\up{e} position dans toutes mes phrases de production.
\item[$\square$] Je mets une majuscule à tous les noms allemands et j'utilise les guillemets „ … “.
\item[$\square$] Je sais rédiger un résumé de 4 à 5 phrases avec des connecteurs logiques.
\end{itemize}
\end{aretenir}

\findecours

\end{document}
